\documentclass[a5paper,papersize]{corona-a5}

% URLの改行設定（参考文献の長いURLがはみ出さないようにする）
% url → hyperref の順で読み込む必要がある
\usepackage{url}
\makeatletter
\g@addto@macro{\UrlBreaks}{%
  \do\a\do\b\do\c\do\d\do\e\do\f\do\g\do\h\do\i\do\j%
  \do\k\do\l\do\m\do\n\do\o\do\p\do\q\do\r\do\s\do\t%
  \do\u\do\v\do\w\do\x\do\y\do\z\do\A\do\B\do\C\do\D%
  \do\E\do\F\do\G\do\H\do\I\do\J\do\K\do\L\do\M\do\N%
  \do\O\do\P\do\Q\do\R\do\S\do\T\do\U\do\V\do\W\do\X%
  \do\Y\do\Z\do\0\do\1\do\2\do\3\do\4\do\5\do\6\do\7%
  \do\8\do\9}
\makeatother
\usepackage[dvipdfmx,plainpages=false,pdfpagelabels,hypertexnames=false]{hyperref}
\usepackage{pxjahyper}
% corona-a5.cls のマクロがPDFしおりでタイトル重複を起こすため無害化
\pdfstringdefDisableCommands{%
  \def\mokuSEVENjidori#1{#1}%
  \def\SEVENjidori#1{#1}%
  \def\kintou#1#2{#2}%
}
\usepackage[dvipsnames]{xcolor}
\usepackage[dvipdfmx]{graphicx}
\usepackage{amsmath}
\usepackage{amssymb}
\usepackage{ascmac}
\usepackage{type1cm}
\usepackage{multicol}
\usepackage{array}
\usepackage{multirow}
\usepackage{booktabs}
\usepackage{longtable}
\usepackage[utf8]{inputenc}
\usepackage{float}
\usepackage{listings}

% 索引（mendex で処理）
\usepackage{makeidx}
\makeindex

% Pythonコードブロックの設定
\lstset{
  language=Python,
  basicstyle=\ttfamily\small,
  frame=single,
  breaklines=true,
  numbers=left,
  numberstyle=\tiny,
  keywordstyle=\color{blue},
  commentstyle=\color{green!50!black},
  stringstyle=\color{red!70!black},
  showstringspaces=false,
  tabsize=4
}

% AIプロンプト例のボックス環境
\newenvironment{promptbox}{\begin{itembox}[l]{プロンプト例}}{\end{itembox}}

% 章の学習目標ボックス用環境（ascmacのscreenを利用）
\newenvironment{chapterbox}{%
  \begin{screen}
}{%
  \end{screen}
}

% bibunits: 各章ごとに参考文献リストを配置
\usepackage{bibunits}
\defaultbibliographystyle{unsrt}
\defaultbibliography{references}

% bibunits + hyperref: 同一文献キーが章ごとに重複しても
% 参照先アンカーが衝突しないように bibunit ごとの接尾辞を付ける。
\makeatletter
\providecommand*{\@extra@b@citeb}{}
\gdef\@extra@binfo{}
\def\@citex[#1]#2{\leavevmode
  \let\@citea\@empty
  \@cite{\@for\@citeb:=#2\do
    {\@citea\def\@citea{,\penalty\@m\ }%
     \edef\@citeb{\expandafter\@firstofone\@citeb\@empty}%
     \if@filesw\immediate\write\@auxout{\string\citation{\@citeb}}\fi
     \@ifundefined{b@\@citeb\@extra@b@citeb}{\hbox{\reset@font\bfseries ?}%
       \G@refundefinedtrue
       \@latex@warning
         {Citation `\@citeb' on page \thepage \space undefined}}%
       {\@cite@ofmt{\csname b@\@citeb\@extra@b@citeb\endcsname}}}}{#1}}
\def\nocite#1{\@bsphack
  \ifx\@onlypreamble\document
    \@for\@citeb:=#1\do{%
      \edef\@citeb{\expandafter\@firstofone\@citeb}%
      \if@filesw\immediate\write\@auxout{\string\citation{\@citeb}}\fi
      \@ifundefined{b@\@citeb\@extra@b@citeb}{\G@refundefinedtrue
          \@latex@warning{Citation `\@citeb' undefined}}{}}%
  \else
    \AddToHook{begindocument/end}[kernel]{\nocite{#1}}%
  \fi
  \@esphack}
\let\std@nocite\nocite
\AtBeginDocument{%
  \renewcommand\bibcite[2]{\global\@namedef{b@#1\@extra@binfo}{#2}}%
}
\def\@startbibunit{%
  \global\let\@startbibunitorrelax\relax
  \global\let\@finishbibunit\@finishstartedbibunit
  \global\advance\@bibunitauxcnt 1
  \xdef\@extra@b@citeb{@\@bibunitname}%
  \xdef\@extra@binfo{\@extra@b@citeb}%
  \if@filesw
    {\endlinechar-1
    \@input{\@bibunitname.aux}}%
    \immediate\openout\@bibunitaux\@bibunitname.aux
    \immediate\write\@bibunitaux{\string\bibstyle{\@localbibstyle}}%
  \fi
}
\let\BU@orig@endbibunit\endbibunit
\def\endbibunit{%
  \BU@orig@endbibunit
  \gdef\@extra@b@citeb{}%
  \gdef\@extra@binfo{}%
}
\makeatother

% 全章の参考文献番号ラベル幅を統一（2桁幅）
\makeatletter
\let\corona@thebibliography\thebibliography
\renewcommand{\thebibliography}[1]{\corona@thebibliography{99}}
\makeatother

\newcommand{\AmSLaTeX}{\leavevmode\hbox{$\mathcal A\kern-.2em\lower.376ex%
        \hbox{$\mathcal M$}\kern-.2em\mathcal S$-\LaTeX}}
\newcommand{\sibun}{\hspace{0.25zw}}

% hyperref しおり階層: corona-a5.cls 独自セクション命令の登録
\makeatletter
\@ifundefined{toclevel@APPsection}{\def\toclevel@APPsection{1}}{}
\@ifundefined{toclevel@TUKIMONOchapter}{\def\toclevel@TUKIMONOchapter{0}}{}
\makeatother

\begin{document}

\frontmatter
\tableofcontents
\cleardoublepage
\begin{preface}{2026年}{松野裕}
2022年11月、OpenAIがChatGPTを公開したとき、世界は大きく変わった。
わずか2か月で1億人のユーザーを獲得したこの対話型AIは、研究者の日常にも急速に浸透し始めた。
文献調査、論文執筆、プログラミング、データ分析——あらゆる研究活動において、
生成AIは無視できない存在となっている。

しかし、その急速な普及に対して、大学院教育における体系的な指導は追いついていない。
多くの大学院生は、SNSやネット記事の断片的な情報を頼りに、手探りでAIツールを使っている。
ある学生は便利さに感動して過度に依存し、別の学生は「使ってはいけないもの」と
誤解して一切触れようとしない。どちらも、研究者として望ましい姿勢とは言えない。
本書は、この状況を踏まえて執筆したものである。大学院生が生成AIを批判的に、
かつ効果的に活用するための体系的なガイドとして、本書を執筆した。

筆者は2025年度より、日本大学大学院における「情報論」の授業で、
大学院生を対象にAI活用の指導を行ってきた。
授業を通じて痛感したのは、(1)研究活動におけるAI活用に焦点を当てた体系的な教材の不在、
(2)プロンプトエンジニアリングの理論と実践のギャップ、
(3)何をAIに任せてよいのかという倫理的な判断力の必要性、の3点である。
2025年度、2026年度の2年間にわたる授業実践を通じて蓄積した知見と、
受講生からのフィードバックが、本書の土台となっている。

本書を貫く基本哲学は、「批判的AI活用」（Critical AI Usage）である。
これは、AIを盲目的に信頼するのでも、頭ごなしに拒否するのでもない。
AIの出力を常に批判的に検証し、自らの専門知識と判断力をもって取捨選択する姿勢のことである。
具体的には、(1) AIは助手であり、研究者の代替ではない、
(2) AIの出力は常に検証する、(3) AIの使用を透明にする、という3つの原則に基づいている。

本書は、半期15回の授業に対応した構成となっている。第1章から第10章までの各章が
1〜2回の授業に対応し、講義と演習を組み合わせて進めることを想定している。
自習で読み進める場合は、まず第1章と第2章を通読し、基本的な考え方と技法を
身につけることを勧める。その後は、自分の研究活動で最も必要な章から読み進めてよい。

本書は、特定のAIツールに依存しない設計としている。
本文中ではChatGPT、Claude、Geminiなど多様なツールを例として取り上げるが、
個々のツールの操作方法よりも、どのツールにも共通する活用の原則と思考法を重視している。
AI分野の進化は速く、書籍の出版後にも新たなツールや技法が登場することが予想されるため、
GitHubリポジトリにて補足資料や更新情報を提供する予定である。

本書の執筆にあたり、多くの方々にお世話になった。
日本大学大学院理工学研究科の「情報論」受講生の皆さんには、
授業を通じて多くのフィードバックをいただいた。
学生の率直な質問や試行錯誤の経験が、本書の内容を大きく豊かにしてくれた。
心より感謝する。
コロナ社編集部には、企画段階から丁寧なご助言をいただいた。
Webサポートとの連携という出版形態を快くお認めいただいたことに感謝する。
最後に、執筆期間中、常に支えてくれた家族に感謝したい。
\end{preface}

\newpage
\mbox{}\thispagestyle{empty}
\mainmatter

\begin{bibunit}
\input{contents/chapter1}
\putbib
\end{bibunit}

\begin{bibunit}
\chapter{プロンプトエンジニアリング}
\label{chap2}

\begin{chapterbox}
\textbf{【この章で学ぶこと】}
\begin{itemize}
\item 効果的なプロンプトの設計原則（ROCFモデル）を理解する
\item Zero-shot、Few-shot、Chain-of-Thoughtなどの技法を使い分けられる
\item ハルシネーションを検出し対処する方法を身につける
\end{itemize}
\end{chapterbox}
\vspace{5mm}

%=========================================
\section{なぜプロンプトが重要か}
\label{sec:2.1}
%=========================================

\subsection{プロンプト次第で結果は大きく変わる}

生成AIから良い出力を得るためには、良い入力（プロンプト\index{ぷろんぷと@プロンプト}）が欠かせない。同じ目的であっても、プロンプトの書き方によって、得られる回答の質は大きく異なる。

以下の例を見てほしい。ある修士学生が、自分の研究テーマに関する先行研究を調べたい場面を想定する。

\begin{promptbox}
\textbf{【悪いプロンプトの例】}

深層学習の論文を教えて
\end{promptbox}

この質問に対して、AIは「深層学習」という広大な分野の中から、何を回答すべきか判断できない。結果として、一般的すぎる回答や、質問者の意図とずれた回答が返ってくることが多い。

\begin{promptbox}
\textbf{【良いプロンプトの例】}

私は修士1年で、「Few-shot学習を用いた製造業における\\
外観検査の自動化」をテーマに研究を始めます。

以下の条件で、関連する先行研究を5本紹介してください。

条件:
\begin{itemize}
\item 2021年以降に発表された査読付き論文
\item Few-shot learningまたはMeta-learningの手法を用いている
\item 製造業の品質管理に関連する応用研究
\item 各論文について、タイトル、著者、掲載誌、主要な貢献を\\
100字程度で説明してください
\end{itemize}

注意: 実在する論文のみを挙げてください。確信が持てない場合は「確認が必要」と明記してください。
\end{promptbox}

このプロンプトでは、背景（修士1年で研究を始める段階）、具体的な研究テーマ、求める情報の条件、出力形式が明確に指定されている。AIはこれらの情報を手がかりに、より適切な回答を生成できる。

\begin{screen}
\textbf{キーポイント}: プロンプトは「AIへの仕様書」である。ソフトウェア開発において曖昧な仕様書から良い製品が生まれないのと同様に、曖昧なプロンプトから良い回答は生まれない。
\end{screen}

\subsection{プロンプトエンジニアリングとは}

\textbf{プロンプトエンジニアリング}\index{ぷろんぷとえんじにありんぐ@プロンプトエンジニアリング}（Prompt Engineering）とは、AIから望ましい出力を得るために、入力（プロンプト）を設計・最適化する技術のことである\cite{liu2023pretrain}。

プロンプトエンジニアリングは、単なる「質問の仕方」にとどまらない。それは以下の要素を含む、体系的な設計プロセスである。

\begin{enumerate}
\item \textbf{目的の明確化}: AIに何をさせたいのかを明確にする
\item \textbf{文脈の設計}: AIが適切な回答を生成するために必要な背景情報を選定する
\item \textbf{制約の指定}: 出力の形式、長さ、スタイルなどの条件を設定する
\item \textbf{反復的改善}: 結果を評価し、プロンプトを修正するサイクルを回す
\end{enumerate}

研究者にとってプロンプトエンジニアリングは、今後ますます重要なスキルとなる。それは、AIの能力を引き出すためにも、AIの限界を見極めるためにも重要である。

%=========================================
\section{ROCFモデル：良いプロンプトの4要素}
\label{sec:2.2}
%=========================================

本書では、効果的なプロンプトを設計するためのフレームワークとして、\textbf{ROCFモデル}\index{ROCFもでる@ROCFモデル}を提案する。ROCFは以下の4要素の頭文字である。

\begin{itembox}[l]{ROCFモデルの4要素}
\begin{itemize}
\item \textbf{R}ole（役割）--- AIに専門家の役割を与える
\item \textbf{O}bjective（目的）--- 何をしてほしいか明確に
\item \textbf{C}ontext（文脈）--- 背景情報を提供する
\item \textbf{F}ormat（形式）--- 出力形式を指定する
\end{itemize}
\end{itembox}

\subsection{Role（役割）：AIに専門家の役割を与える}

AIに特定の専門家の役割を付与することで、その分野に適した語彙、視点、判断基準で回答を生成させることができる。

\begin{promptbox}
\textbf{【Role要素の例】}

例1: ``あなたは計算機科学分野の経験豊富な査読者です。''

例2: ``あなたは統計学の教授で、大学院生の研究指導を行っています。''

例3: ``あなたはテクニカルライターで、学術論文の英語校正を専門としています。''
\end{promptbox}

Roleの効果は、単に回答のトーンを変えるだけではない。たとえば、「査読者として」と指定すると、AIは論文の弱点や改善点を指摘する方向で回答を生成する傾向がある。「指導教員として」と指定すると、学生の理解を助けるような丁寧な説明を生成しやすくなる。

\begin{screen}
\textbf{ヒント}: Roleは具体的であるほど効果的である。「専門家として」よりも「機械学習分野で20年の研究経験を持つ教授として、修士学生の論文を指導する立場で」と指定した方が、より適切な回答を得やすい。
\end{screen}

\subsection{Objective（目的）：何をしてほしいか明確に}

Objectiveは、AIに実行してほしいタスクを明確に指定する要素である。曖昧な依頼は曖昧な結果を生む。

\begin{promptbox}
\textbf{【Objective要素の比較】}

曖昧な目的: ``この論文について教えてください''

明確な目的: ``この論文の研究手法と実験結果を、\\
以下の3つの観点から分析してください:\\
(1) 新規性、(2) 再現性、(3) 既存手法との比較''
\end{promptbox}

Objectiveを書く際のポイント:
\begin{enumerate}
\item \textbf{動詞を明確にする}: 「教えてください」ではなく「要約してください」「比較してください」「評価してください」など
\item \textbf{範囲を限定する}: 「すべて」ではなく「上位3つ」「主要な点」「特にXに関して」など
\item \textbf{成功基準を示す}: 何をもって良い回答とするかを伝える
\end{enumerate}

\subsection{Context（文脈）：背景情報を提供する}

Contextは、AIが適切な回答を生成するために必要な背景情報である。研究活動では特に重要な要素となる。

\begin{promptbox}
\textbf{【Context要素の例】}

``以下の背景を踏まえて回答してください:
\begin{itemize}
\item 私は情報工学専攻の修士1年で、深層学習を用いた自然言語処理が専門です
\item 対象とする言語は日本語です
\item 使用可能な計算資源はGPU 1台（NVIDIA A100）です
\item 投稿先はACL 2027を予定しています
\item 研究期間は1年半です''
\end{itemize}
\end{promptbox}

適切なContextを提供することで、AIはより現実的で実行可能な提案を行えるようになる。たとえば、「GPU 1台」という制約を伝えれば、計算コストが膨大なモデルの提案を避け、効率的なアプローチを推薦してくれる。

\subsection{Format（形式）：出力形式を指定する}

Formatは、AIの出力がどのような形式であるべきかを指定する要素である。

\begin{promptbox}
\textbf{【Format要素の例】}

例1: ``箇条書きで5項目にまとめてください''

例2: ``表形式で比較してください（列: 手法名、精度、計算コスト、利点、欠点）''

例3: ``\LaTeX 形式で数式を含めて記述してください''

例4: ``以下のJSON形式で出力してください''
\end{promptbox}

\subsection{ROCFモデルの実践——4要素を組み合わせる}

4つの要素を組み合わせた実践的なプロンプトの例を示す。

\begin{promptbox}
\textbf{【ROCFモデルの完全な適用例】}

\textbf{[Role]} あなたは自然言語処理分野の研究に精通した\\
大学教授です。修士学生の研究指導を日常的に行っています。

\textbf{[Objective]} 以下の研究テーマについて、修士論文の\\
構成案（章立て）を提案してください。各章の概要と、\\
そこで扱うべき内容を200字程度で説明してください。

\textbf{[Context]}
\begin{itemize}
\item 研究テーマ: 「事前学習済み言語モデルを用いた日本語法律文書の自動要約」
\item 修士2年の学生で、修士論文の提出は6か月後です
\item すでにベースラインの実験は完了し、提案手法の実装中です
\item 関連研究として、BERTを用いた法律文書分類の研究はすでにサーベイ済みです
\end{itemize}

\textbf{[Format]} 以下の形式で出力してください:

第X章 [章タイトル] / 概要: [200字程度の説明] / 含めるべき内容: [項目リスト]
\end{promptbox}

\begin{screen}
\textbf{キーポイント}: ROCFモデルの4要素すべてを毎回使う必要はない。タスクの性質に応じて、必要な要素を選択すればよい。ただし、ObjectiveとContextは常に含めることを推奨する。
\end{screen}

%=========================================
\section{基本的なプロンプト技法}
\label{sec:2.3}
%=========================================

\subsection{Zero-shot Prompting（例なしで指示）}

\textbf{Zero-shot prompting}\index{Zero-shot prompting}は、具体例を一切提示せず、指示文のみでタスクを実行させる技法である。最も基本的なプロンプト技法であり、明確なタスクには十分に機能する。

\begin{promptbox}
\textbf{【Zero-shot promptingの例】}

以下の論文タイトルを、学術的な日本語に翻訳してください。

``Retrieval-Augmented Generation for Knowledge-Intensive NLP Tasks''
\end{promptbox}

\textbf{適している場面:}
\begin{itemize}
\item タスクが明確で、AIが十分な学習データを持っている場合
\item 翻訳、要約、校正などの定型的なタスク
\item 簡単な質問応答
\end{itemize}

\textbf{注意点:}
\begin{itemize}
\item タスクの定義が曖昧な場合、AIが意図と異なる解釈をする可能性がある
\item 出力形式を厳密に制御したい場合は、追加の指定が必要
\end{itemize}

\subsection{Few-shot Prompting（例を添えて指示）}

\textbf{Few-shot prompting}\index{Few-shot prompting}は、望ましい入出力の例を数個提示した上で、タスクを実行させる技法である\cite{brown2020language}。AIに「こういう形式で、こういう品質の回答を期待している」というお手本を示すことで、出力の品質と一貫性が大幅に向上する。

\begin{promptbox}
\textbf{【Few-shot promptingの例：論文の分類】}

以下の論文タイトルを、研究分野に分類してください。

例1:

タイトル: ``BERT: Pre-training of Deep Bidirectional Transformers for Language Understanding''

分野: 自然言語処理 $>$ 事前学習モデル

例2:

タイトル: ``ResNet: Deep Residual Learning for Image Recognition''

分野: コンピュータビジョン $>$ 画像認識

例3:

タイトル: ``Playing Atari with Deep Reinforcement Learning''

分野: 強化学習 $>$ ゲームAI

では、以下の論文を分類してください:

タイトル: ``Generative Adversarial Nets''

分野:
\end{promptbox}

\textbf{適している場面:}
\begin{itemize}
\item 出力の形式や品質を厳密に制御したい場合
\item AIがあまり得意でないタスク（特殊な分類基準など）
\item 独自の評価基準や記述スタイルを適用したい場合
\end{itemize}

\textbf{Few-shotの例数の目安:}
\begin{itemize}
\item 2〜5個の例があれば、多くのタスクで十分
\item 例が多すぎると、コンテキスト長を浪費する
\item 例は多様性を持たせる（似た例ばかりにしない）
\end{itemize}

\subsection{Chain-of-Thought（ステップバイステップの思考を促す）}

\textbf{Chain-of-Thought（CoT）prompting}\index{Chain-of-Thought}は、AIに思考過程を段階的に出力させる技法である\cite{Wei2022ChainOfThought}。複雑な推論が必要なタスクで、回答の正確性が著しく改善される。

\begin{promptbox}
\textbf{【Chain-of-Thought promptingの例：研究手法の選択】}

以下の研究課題に対して、最適な実験手法を選択してください。\\
ステップバイステップで考え、各ステップの根拠を示してください。

研究課題: 「SNS投稿のテキストから、投稿者の感情状態を\\
4カテゴリ（喜び、怒り、悲しみ、恐怖）に分類するモデルを\\
開発したい。ただし、ラベル付きデータは各カテゴリ100件\\
しかない。」

ステップ1: まず、この課題の本質的な要件を整理してください

ステップ2: 考えられる手法を3つ以上列挙してください

ステップ3: 各手法のメリット・デメリットを分析してください

ステップ4: データ量の制約を考慮して、最適な手法を推薦してください

ステップ5: 推薦した手法の実装計画を簡潔に示してください
\end{promptbox}

Chain-of-Thoughtの効果が特に高い場面:
\begin{itemize}
\item 数学的な推論が必要な場合
\item 複数の要因を考慮した判断が必要な場合
\item 複雑な比較・分析が必要な場合
\end{itemize}

\begin{screen}
\textbf{キーポイント}: 「ステップバイステップで考えてください」という一文を追加するだけでも、回答の質が向上することが多い。これは、AIが中間的な推論過程を生成することで、最終的な結論の正確性が高まるためである。
\end{screen}

\subsection{各技法の比較と使い分け}

\begin{table}[htbp]
\centering
\caption{プロンプト技法の比較}
\label{tab:prompt_techniques}
\small
\begin{tabular}{llcl}
\toprule
技法 & 適用場面 & 例の必要数 & 効果 \\
\midrule
Zero-shot & 明確な定型タスク & 0 & 速い、簡潔 \\
Few-shot & 形式の制御が必要 & 2〜5 & 出力品質の安定化 \\
Chain-of-Thought & 複雑な推論 & 0〜3 & 推論精度の向上 \\
\bottomrule
\end{tabular}
\end{table}

実際の研究活動では、これらの技法を組み合わせて使うことが多い。たとえば、Few-shotの例にChain-of-Thoughtの思考過程を含めることで、「こういう思考プロセスで、こういう形式の回答を期待している」ということをAIに伝えることができる（Few-shot CoT）。

%=========================================
\section{高度なプロンプト技法}
\label{sec:2.4}
%=========================================

基本的な技法を習得したら、次はより高度な技法に挑戦しよう。これらの技法は、研究活動の中でも特に複雑なタスクで効果的である。

\subsection{ペルソナ設定}

ペルソナ設定\index{ぺるそなせってい@ペルソナ設定}は、ROCFモデルのRole要素をさらに発展させた技法である。AIに特定の視点から回答させることで、多角的な分析や批判的なフィードバックを得ることができる。

\begin{promptbox}
\textbf{【ペルソナ設定の実践例：論文の多角的レビュー】}

\textbf{--- ペルソナ1: 厳格な査読者 ---}

あなたは、トップカンファレンス（NeurIPS, ICML）の査読を\\
10年以上務めてきた厳格な研究者です。以下の論文概要に対して、\\
主要な弱点と改善すべき点を5つ指摘してください。\\
特に、技術的な新規性と実験の妥当性について厳しく評価してください。

{[}論文概要をここに貼る]
\end{promptbox}

\begin{promptbox}
\textbf{--- ペルソナ2: 好意的な指導教員 ---}

あなたは、修士学生の研究指導に情熱を持つ教授です。\\
以下の論文概要に対して、良い点を3つ、改善の余地がある点を\\
3つ指摘し、各改善点に対して具体的なアドバイスを提供して\\
ください。学生のモチベーションを維持しながら、\\
建設的なフィードバックを心がけてください。

{[}論文概要をここに貼る]
\end{promptbox}

\begin{promptbox}
\textbf{--- ペルソナ3: 隣接分野の研究者 ---}

あなたは、情報検索（IR）分野の研究者です。\\
自然言語処理（NLP）の専門家ではありませんが、テキスト処理に\\
関連する知識を持っています。以下の論文概要を読んで、\\
IR分野の研究者として疑問に思う点や、\\
IR分野との接点を指摘してください。

{[}論文概要をここに貼る]
\end{promptbox}

同じ論文概要に対して、異なるペルソナからのフィードバックを得ることで、自分では気づかなかった観点や問題点を発見できることがある。

\subsection{段階的精緻化（Iterative Refinement）}

段階的精緻化\index{だんかいてきせいちか@段階的精緻化}は、AIとの対話を複数ターンに分け、出力を少しずつ改善していく技法である。一度で完璧な出力を求めるのではなく、反復的に品質を高めていく。

\begin{promptbox}
\textbf{【段階的精緻化の例：研究要旨の作成】}

\textbf{--- ターン1: ドラフト生成 ---}

以下の情報をもとに、国際会議に投稿する論文のAbstract（250語以内、英語）のドラフトを作成してください。

研究テーマ: [テーマ] / 手法: [手法の概要] / 結果: [主要な結果] / 貢献: [主要な貢献]

\textbf{--- ターン2: 構造の改善 ---}

ありがとうございます。以下の点を修正してください:
(1) 最初の文で、研究の動機をより明確にしてください
(2) 手法の説明をもう少し具体的にしてください
(3) 結果の数値を追加してください: 精度92.3\%, ベースライン比+4.7ポイント

\textbf{--- ターン3: 言語の洗練 ---}

構造は良くなりました。次に、以下の観点で言語面を改善してください:
(1) 受動態を能動態に変えられる箇所を修正
(2) 冗長な表現を簡潔にする
(3) 学術的な語彙として不適切な表現があれば修正
\end{promptbox}

段階的精緻化のメリット:
\begin{itemize}
\item 各ステップで焦点を絞れるため、改善の方向が明確になる
\item AIの出力を確認しながら修正指示を出せるため、最終的な品質が高くなる
\item 人間の意図をAIに徐々に正確に伝えることができる
\end{itemize}

\subsection{メタプロンプティング}

\textbf{メタプロンプティング}\index{めたぷろんぷてぃんぐ@メタプロンプティング}は、AIに「プロンプト自体を考えさせる」技法である\cite{zhou2023large}。自分のタスクに最適なプロンプトがわからない場合に有効である。

\begin{promptbox}
\textbf{【メタプロンプティングの例】}

私は以下のタスクを実行したいのですが、最適なプロンプトの書き方がわかりません。

タスク: 修士論文の「関連研究」セクションを書くために、10本の論文を体系的に整理し、比較分析したい。

このタスクに対して、最も効果的なプロンプトを設計してください。以下の点を含めてください:
\begin{enumerate}
\item プロンプトの全文
\item なぜその書き方が効果的なのかの説明
\item プロンプトをさらに改善するためのヒント
\end{enumerate}
\end{promptbox}

メタプロンプティングは、プロンプトエンジニアリングの学習そのものにも有効である。AIが設計するプロンプトを分析することで、効果的なプロンプトの特徴を学ぶことができる。

\subsection{制約条件の活用}

AIの出力を精密に制御するためには、適切な制約条件\index{せいやくじょうけん@制約条件}を設定することが有効である。

\begin{promptbox}
\textbf{【制約条件の活用例】}

以下の制約条件に従って、研究計画書のドラフトを作成してください。

制約条件:
\begin{enumerate}
\item 文字数: 2000字以内
\item 構成: 背景、目的、手法、期待される成果の4セクション
\item 文体: 「である」調、学術的な表現
\item 含めるべきキーワード: 深層学習、転移学習、ドメイン適応
\item 避けるべき表現: 「画期的」「世界初」などの誇張表現
\item 参考にすべきスタイル: 科研費申請書の研究計画調書
\item 対象読者: 情報工学分野の審査委員（AI専門とは限らない）
\end{enumerate}
\end{promptbox}

効果的な制約条件のカテゴリを表\ref{tab:constraint_categories}に示す。

\begin{table}[htbp]
\centering
\caption{制約条件のカテゴリと例}
\label{tab:constraint_categories}
\begin{tabular}{ll}
\toprule
カテゴリ & 例 \\
\midrule
長さの制約 & 「300字以内」「5項目で」 \\
形式の制約 & 「表形式で」「\LaTeX で」「箇条書きで」 \\
内容の制約 & 「技術的な詳細は省く」「数式を含める」 \\
否定の制約 & 「専門用語を使わない」「主観的な評価は含めない」 \\
品質の制約 & 「学術論文レベルの正確さで」「初学者にもわかるように」 \\
スタイルの制約 & 「である調」「フォーマルな英語」 \\
\bottomrule
\end{tabular}
\end{table}

\begin{screen}
\textbf{ヒント}: 制約条件は多すぎると逆効果になることがある。AIが全ての制約を同時に満たそうとして、かえって不自然な出力になる場合がある。重要度の高い制約を3〜5個に絞ることを推奨する。
\end{screen}

%=========================================
\section{実践的なプロンプト設計技法}
\label{sec:practical_prompt}
%=========================================

前節までに紹介した技法に加え、AIサービス提供企業が公開しているベストプラクティスから、研究活動に役立つ実践的な技法を紹介する。本節では、特にAnthropicのプロンプトエンジニアリングガイドに基づく手法を取り上げる。

\subsection{XMLタグによるプロンプトの構造化}

複雑なプロンプトを作成する際、\textbf{XMLタグ}\index{XMLたぐ@XMLタグ}を用いてプロンプトの各部分を明示的に区切ることで、AIが指示・文脈・入力データを曖昧さなく解釈できるようになる。たとえば、指示部分を\texttt{<instructions>}タグで、分析対象の文書を\texttt{<document>}タグで囲むといった方法である。

この技法は、指示文、背景情報、具体的なデータ、出力形式の指定など、複数の異なる種類の情報を1つのプロンプトに含める場合に特に有効である。タグによって各要素の境界が明確になるため、AIが指示の一部をデータと誤認したり、データの一部を指示として解釈するといった問題を防ぐことができる。

研究活動では、論文をAIに分析させる場面でこの技法が威力を発揮する。論文テキストを\texttt{<document>}タグで囲み、分析の依頼を別のタグで明示することで、AIは「何が分析対象で、何が指示なのか」を正確に把握できる。

\begin{promptbox}
\textbf{【XMLタグを用いた構造化プロンプトの例】}

\texttt{<instructions>}

以下の論文のアブストラクトを読み、背景・手法・結果・限界の4点で要約してください。

\texttt{</instructions>}

\texttt{<document>}

{[}論文のアブストラクトをここに貼り付け]

\texttt{</document>}

\texttt{<output\_format>}

箇条書きで、各項目2--3文程度にまとめてください。

\texttt{</output\_format>}
\end{promptbox}

\subsection{長文コンテキストの効果的な配置}
\index{こんてきすとはいち@コンテキスト配置}

長い文書（論文全文やデータセット）をAIに与えて分析や質問応答を行う場合、\textbf{文書をプロンプトの先頭に配置し、質問や指示を末尾に配置する}ことで、回答品質が向上することがAnthropicの検証により報告されている。改善幅は最大30\%に達する場合もあるとされる。

この配置が効果的な理由は、LLMがプロンプト末尾の情報に対してより強い注意を向ける傾向があるためである。長い文書の後に質問を配置することで、AIは「いま何を求められているのか」を明確に認識した状態で文書内容を参照できる。

研究活動では、論文の全文をAIに読み込ませて質問する場面が多い。そのような場合、「論文テキスト→質問」の順にプロンプトを構成すると、回答の的確さが向上する。

\begin{screen}
\textbf{ヒント}: 長い論文やデータをAIに分析させるとき、文書を先に、質問を後に配置すると回答品質が向上する。「資料→指示」の順を心がけよう。
\end{screen}

\subsection{引用ベースの回答で正確性を高める}
\index{いんようべーすかいとう@引用ベース回答}

AIに文書の分析を依頼する際、まず\textbf{ソース文書から関連部分を引用させ、その後に分析を述べさせる}という手順を踏ませることで、ハルシネーションを大幅に抑制できる。この手法は\textbf{グラウンディング}\index{ぐらうんでぃんぐ@グラウンディング}（grounding）と呼ばれ、AIの回答を原文に「根拠づける」効果がある。

引用を先に行わせることで、AIは回答の根拠を原文に求めるようになり、「もっともらしいが原文にない情報」を生成するリスクが低減する。これは文献レビューや論文査読のように、原文の正確な参照が求められる研究場面で特に有用である。

\begin{promptbox}
\textbf{【引用ベースの回答を促すプロンプトの例】}

以下の論文から、研究手法に関する記述を引用してください。

引用部分を\texttt{<quotes>}タグで囲み、その後に\texttt{<analysis>}タグ内で分析を行ってください。

{[}論文テキスト]
\end{promptbox}

\subsection{指示の明確化：「やるな」より「代わりにこうして」}

AIに出力を制御する指示を与える際、\textbf{否定形の指示（「〜しないでください」）よりも、肯定形の指示（「代わりに〜してください」）の方が効果的}である。否定形の指示は「何をすべきでないか」を伝えるだけで、AIに明確な行動の方向を示さない。一方、肯定形の指示はAIが取るべき行動を具体的に示すため、出力がより安定する。

\vspace{3mm}
\textbf{悪い例と良い例の比較:}
\vspace{2mm}

\begin{itemize}
\item \textbf{悪い例}: 「箇条書きにしないでください」

$\rightarrow$ AIは箇条書き以外の何をすべきか判断に迷う

\item \textbf{良い例}: 「段落形式で、各段落に1つのポイントを含めて書いてください」

$\rightarrow$ AIは具体的な出力形式を把握できる
\end{itemize}

\vspace{2mm}

もう一つの例を挙げる。

\begin{itemize}
\item \textbf{悪い例}: 「マークダウンを使わないでください」
\item \textbf{良い例}: 「滑らかな段落構成の文章で書いてください」
\end{itemize}

\subsection{文脈と理由を添える}

AIへの指示では、\textbf{何をしてほしいかだけでなく、なぜそうしてほしいかを説明する}ことで、AIがその意図を汎化的に理解し、より適切に対応できるようになる。理由が明確であれば、指示で明示的にカバーしていない状況でも、AIは意図に沿った判断を下しやすくなる。

\vspace{3mm}
\textbf{悪い例と良い例の比較:}
\vspace{2mm}

\begin{itemize}
\item \textbf{悪い例}: 「省略記号（\ldots）を使わないでください」

$\rightarrow$ AIは指示には従えるが、なぜかわからないため類似の状況で判断を誤る可能性がある

\item \textbf{良い例}: 「この文章はテキスト読み上げソフトで使用するため、省略記号（\ldots）は使わないでください。読み上げソフトが省略記号を正しく発音できないためです。」

$\rightarrow$ AIは理由を理解し、省略記号以外にも読み上げに問題がありそうな表現を自主的に避けられる
\end{itemize}

\vspace{2mm}

研究活動では、たとえば「数式は\LaTeX 形式で書いてください。この文章をOverleafにそのまま貼り付けて使うためです。」のように理由を添えることで、数式以外の部分でも\LaTeX に適した記法をAIが選択するようになる。

\begin{screen}
本節の内容はAnthropicのプロンプトエンジニアリングガイド（Anthropic公式サイト「Prompting best practices」）を参考にしている。各AIツールの最新のベストプラクティスは、公式ドキュメントを確認されたい。
\end{screen}

%=========================================
\section{ハルシネーションの理解と対策}
\label{sec:2.5}
%=========================================

\subsection{ハルシネーションとは何か}

\textbf{ハルシネーション}\index{はるしねーしょん@ハルシネーション}（Hallucination）とは、AIがもっともらしいが事実に反する情報を生成する現象である。第\ref{chap1}章でも触れたが、プロンプトエンジニアリングの観点から、ここでより詳しく解説する。

ハルシネーションが生じる技術的な背景を簡潔に説明する。LLMは「次のトークン（単語）を予測する」モデルである。この予測は、学習データから得られた統計的パターンに基づいている。つまり、LLMは「事実を記憶している」のではなく、「もっともらしい文を生成している」のである。

この「もっともらしさ」と「正確さ」のギャップがハルシネーションの本質である。

\subsection{研究活動で特に危険なハルシネーション}

研究活動では、以下のタイプのハルシネーションが特に危険である。

\paragraph{タイプ1: 架空の論文引用}
最も頻繁に報告されるハルシネーションの一つである。AIは、実在しない論文のタイトル、著者名、雑誌名、出版年、さらにはDOI番号までもっともらしく生成することがある。

\begin{promptbox}
\textbf{【架空の論文引用の例】}

質問: ``ドメイン適応における最新の手法について、主要な論文を3本教えてください''

AIの回答（ハルシネーションを含む）:

1. Zhang, W., \& Liu, H. (2024). ``Progressive Domain\\
Adaptation with Contrastive Learning for Visual Recognition.''\\
IEEE TPAMI, 46(3), 1234--1248.

$\rightarrow$ この論文は実在しない可能性がある。著者名、タイトル、\\
巻号ページがもっともらしく生成されているが、\\
DOIで確認すると存在しない。
\end{promptbox}

\paragraph{タイプ2: 誤った数値データ}\mbox{}

\begin{promptbox}
\textbf{【誤った数値データの例】}

質問: ``ImageNetにおける最新の画像認識精度のSOTAは？''

AIの回答: ``2025年時点で、ImageNetにおける\\
Top-1精度のSOTAは94.7\%であり、\\
これはXYZモデルによって達成されました。''

$\rightarrow$ 数値やモデル名が不正確である可能性がある。\\
特にSOTAは頻繁に更新されるため、AIの学習データの\\
カットオフ以降の情報は誤っている可能性が高い。
\end{promptbox}

\paragraph{タイプ3: もっともらしい虚偽の説明}\mbox{}

\begin{promptbox}
\textbf{【もっともらしい虚偽の説明の例】}

質問: ``Batch Normalizationが有効な理由を教えてください''

AIの回答: ``Batch Normalizationは、内部共変量シフト\\
（Internal Covariate Shift）を軽減することで\\
学習を安定化させます。''

$\rightarrow$ 長年信じられてきた説明だが、\\
Santurkar et al.\ (2018)\cite{santurkar2018batch}の\\
研究により、\\
この説明が主要因ではないことが示されている。\\
AIは「広く引用されている説明」を生成しがちであり、\\
最新の理解を反映していない可能性がある。
\end{promptbox}

\subsection{5つの対策法}

ハルシネーションに対処するための5つの具体的な対策法を紹介する。

\begin{itembox}[l]{ハルシネーション対策の5つの方法}
\textbf{対策1: 出典を要求する}

プロンプトの中で、明示的に出典の提示を要求する。また、出典の信頼性を評価する基準も指定する。

\vspace{2mm}
\textbf{対策2: 複数ツールでクロスチェック}

重要な情報は、複数のAIツールや検索エンジンで確認する。

\vspace{2mm}
\textbf{対策3: 原典にあたって確認}

AIが提示した情報の元となる原典（論文、公式ドキュメント、データソース）を直接確認する。これは最も確実な対策であり、研究者の基本動作でもある。

\vspace{2mm}
\textbf{対策4: 「わからない」と言わせる}

AIに「わからない場合はわからないと回答してよい」と明示的に許可することで、ハルシネーションを抑制できることがある。

\vspace{2mm}
\textbf{対策5: 段階的に検証する}

AIの回答を一度に全て信頼するのではなく、小さな単位に分割して段階的に検証する。
\end{itembox}

\begin{promptbox}
\textbf{【出典を要求するプロンプト例】}

以下の質問に回答する際、各主張に対して出典を明記してください。出典は以下の形式で記載してください:
\begin{itemize}
\item 著者名、論文タイトル、掲載誌/会議名、年
\end{itemize}

また、以下の点に注意してください:
\begin{itemize}
\item 確実に実在する論文のみを引用してください
\item 記憶が曖昧な場合は「要確認」と明記してください
\item 出典を特定できない主張は、「一般的な見解」と注記してください
\end{itemize}
\end{promptbox}

\begin{promptbox}
\textbf{【「わからない」を許可するプロンプト例】}

以下の質問に回答してください。ただし、以下のルールに従ってください:
\begin{itemize}
\item 確信がある情報のみを回答してください
\item 不確実な情報には [要確認] タグを付けてください
\item わからないことは「この点についてはわかりません」と\\
正直に回答してください
\item 推測で回答する場合は、「推測ですが」と明記してください
\end{itemize}

決して架空の情報を生成しないでください。
\end{promptbox}

\begin{screen}
\textbf{キーポイント}: ハルシネーション対策の本質は、「AIの出力を無条件に信頼しない」という姿勢である。これは「AIを使わない」ということではなく、「AIの出力を仮説として扱い、検証してから採用する」ということである。この姿勢は、科学的な研究態度そのものとも言える。
\end{screen}

%=========================================
\section{研究活動におけるプロンプト設計の実践}
\label{sec:2.6}
%=========================================

ここまで学んだ技法を、研究活動の具体的な場面に適用する方法を示す。

\subsection{文献調査用プロンプトの設計}

文献調査は、研究の初期段階で最もAIの支援が有効な場面である。

\begin{promptbox}
\textbf{【文献調査用プロンプト：サーベイの構造化】}

\textbf{[Role]} あなたは情報工学分野の研究に精通した\\
図書館司書です。学術文献の体系的な整理と分類を専門と\\
しています。

\textbf{[Objective]} 以下の研究テーマに関連する研究の\\
分類体系を提案してください。

\textbf{[Context]} 研究テーマ: 「大規模言語モデルの\\
ファインチューニング手法」 / 目的: 修士論文の関連研究\\
セクションの構成を決めるため / 対象期間: 2020年〜2025年

\textbf{[Format]} 大分類 $>$ 中分類 $>$ 小分類。\\
各分類に対して: 定義（50字程度）、代表的な手法名を\\
2〜3個、この分類が重要な理由

注意: 実在する手法名のみを使用してください。
\end{promptbox}

\textbf{ベストプラクティス:}
\begin{itemize}
\item 最初は広く探索し、徐々に焦点を絞る
\item AIが提示した論文は必ずGoogle ScholarやSemantic Scholarで確認する
\item サーベイの構造はAIに提案させ、個々の論文の評価は自分で行う
\end{itemize}

\textbf{アンチパターン:}
\begin{itemize}
\item AIが生成した文献リストをそのまま使う（ハルシネーションのリスク）
\item AIの要約だけで論文を読んだ気になる（原典の精読は不可欠）
\item 1回のプロンプトで完璧なサーベイを期待する（段階的に進める）
\end{itemize}

\subsection{論文執筆用プロンプトの設計}

\begin{promptbox}
\textbf{【論文執筆用プロンプト：Abstractの推敲】}

\textbf{[Role]} あなたは英語ネイティブの学術論文校正者です。\\
計算機科学分野の論文を年間100本以上校正しています。

\textbf{[Objective]} 以下の英語Abstractを推敲してください。\\
3つのバージョンを提示し、それぞれの特徴を説明してください。

\textbf{[Context]} 投稿先: ACL 2027（自然言語処理のトップ\\
会議） / 文字数制限: 250語以内 / 論文の主要な貢献:\\
新しいファインチューニング手法の提案とその有効性の実証

\textbf{[Format]} 各バージョンについて:\\
(1) 修正後のAbstract全文、\\
(2) 主な変更点の一覧（変更前→変更後）、\\
(3) そのバージョンの長所と短所
\end{promptbox}

\textbf{ベストプラクティス:}
\begin{itemize}
\item 最初に構成（アウトライン）を固め、その後で各セクションを執筆する
\item AIに生成させた文章は必ず自分で読み直し、加筆修正する
\item 専門用語の使い方が正しいか確認する
\end{itemize}

\textbf{アンチパターン:}
\begin{itemize}
\item 論文全体をAIに一括で書かせる（品質が低下し、一貫性も失われる）
\item AIの文体をそのまま使う（自分の文体で書き直すべき）
\item 投稿先のスタイルガイドを無視する
\end{itemize}

\subsection{コード生成用プロンプトの設計}

\begin{promptbox}
\textbf{【コード生成用プロンプト：データ分析スクリプト】}

\textbf{[Role]} あなたはPythonのデータ分析に精通したシニアエンジニアです。

\textbf{[Objective]} 以下の仕様に基づいて、データ分析用のPythonスクリプトを作成してください。

\textbf{[Context]} データ: CSVファイル（列: id, age, gender,\\
score1, score2, score3, group） / 行数: 約500行 /\\
目的: グループ間のスコア比較 /\\
環境: Python 3.11, pandas, scipy, matplotlib

\textbf{[仕様]}
\begin{enumerate}
\item CSVファイルを読み込む
\item 基本統計量（平均、標準偏差、中央値）をグループ別に算出
\item グループ間の差の検定（t検定またはMann-Whitney U検定をデータの正規性に基づいて自動選択）
\item 結果を可視化（箱ひげ図、棒グラフ）
\item 結果を\LaTeX 形式の表として出力
\end{enumerate}
\end{promptbox}

\textbf{ベストプラクティス:}
\begin{itemize}
\item 仕様を明確に記述し、あいまいさを排除する
\item 生成されたコードは必ず自分で動作確認する
\item テストケースもAIに生成させ、網羅的にテストする
\item コードの各行が何をしているか理解してから使用する
\end{itemize}

\textbf{アンチパターン:}
\begin{itemize}
\item 理解していないコードをそのまま使う（バグの温床）
\item エラーが出た際に、エラーメッセージを読まずにAIに丸投げする
\item セキュリティに関わるコード（認証、暗号化など）をAIに任せる
\end{itemize}

\subsection{プロンプトテンプレートの管理}

研究活動で頻繁に使うプロンプトは、テンプレートとして保存・管理することを推奨する。

\begin{lstlisting}[language={}, caption=推奨フォルダ構成]
prompts/
  literature_review/
    survey_structure.md     # サーベイの構造化
    paper_summary.md        # 論文の要約
    comparison_table.md     # 論文の比較表作成
  writing/
    abstract_review.md      # Abstract推敲
    related_work.md         # 関連研究セクション
    english_editing.md      # 英語校正
  coding/
    data_analysis.md        # データ分析
    debug.md                # デバッグ支援
    test_generation.md      # テスト生成
  presentation/
    slide_outline.md        # スライド構成
    qa_preparation.md       # 質疑応答準備
\end{lstlisting}

テンプレートの各ファイルには、以下の情報を含めるとよい:
\begin{enumerate}
\item テンプレートの目的と適用場面
\item プロンプトの全文（変数部分は \texttt{[ここに入力]} で示す）
\item 使用時の注意点
\item 過去の使用実績と改善履歴
\end{enumerate}

\begin{screen}
\textbf{ヒント}: プロンプトテンプレートをGitで管理すると、改善の履歴を追跡できる。チーム（研究室）で共有すれば、メンバー全員のAI活用スキルの底上げにもつながる。
\end{screen}

%=========================================
\section*{演習問題}
\addcontentsline{toc}{section}{演習問題}
%=========================================

\begin{enumerate}
\item \textbf{問題2.1（ROCFモデルの適用）}

以下のタスクに対して、ROCFモデルの4要素（Role, Objective, Context, Format）を明示的に含むプロンプトを設計せよ。

タスク: あなたの研究テーマ（または興味のあるテーマ）に関連する先行研究を5本リストアップし、各論文の手法と結果を比較する表を作成したい。

\item \textbf{問題2.2（プロンプト技法の使い分け）}

以下の3つのタスクに対して、それぞれ最適なプロンプト技法（Zero-shot, Few-shot, Chain-of-Thought）を選択し、その技法を用いたプロンプトを作成せよ。なぜその技法を選んだかも説明せよ。

\begin{enumerate}
\item[(a)] 英語の論文タイトルを日本語に翻訳する
\item[(b)] 研究論文のAbstractを「新規性」「有用性」「信頼性」の3観点で評価する
\item[(c)] 2つの研究手法の違いを分析し、どちらが特定の問題設定に適しているか判断する
\end{enumerate}

\item \textbf{問題2.3（ハルシネーション検出の実践）}

任意の対話型AIを使い、以下を実行せよ。

\begin{enumerate}
\item 自分の研究分野に関する質問をし、AIに論文を5本推薦させる
\item 推薦された5本の論文すべてについて、Google ScholarまたはSemantic Scholarで実在を確認する
\item 実在しない論文があった場合、その割合を記録する
\item ハルシネーションを減らすためにプロンプトを改善し、再度実行する
\item 改善前後の結果を比較し、レポートにまとめる
\end{enumerate}

\item \textbf{問題2.4（段階的精緻化の実践）}

任意の対話型AIを使い、以下の手順で段階的精緻化を実践せよ。

\begin{enumerate}
\item 自分の研究テーマのAbstract（200〜300字、日本語）のドラフトをAIに生成させる
\item 生成されたドラフトの問題点を3つ以上指摘し、AIに修正を依頼する
\item 修正後のバージョンをさらに改善するため、別の観点から修正を依頼する
\item 最終バージョンと最初のドラフトを比較し、どのように改善されたかを分析する
\end{enumerate}

\item \textbf{問題2.5（メタプロンプティング）}

AIに「効果的なプロンプトを設計してもらう」というメタプロンプティングを実践せよ。

\begin{enumerate}
\item 以下のタスクをAIに伝える: 「修士論文の中間発表用スライド（10枚程度）の構成案を作成したい」
\item AIに、このタスクに最適なプロンプトを設計させる
\item AIが設計したプロンプトを実際に使用し、結果を評価する
\item AIが設計したプロンプトの良い点と改善すべき点を分析する
\end{enumerate}

\item \textbf{問題2.6（総合演習：プロンプトの比較実験）}

同じタスクに対して、3種類のプロンプト（シンプルなプロンプト、ROCFモデルを適用したプロンプト、Few-shot + Chain-of-Thoughtを組み合わせたプロンプト）を作成し、それぞれの出力を比較せよ。

タスク: 「あなたの研究分野における、今後5年間の主要な研究トレンドを予測する」

評価の観点: 回答の具体性、回答の正確性（検証可能な情報の割合）、回答の有用性（研究活動に実際に役立つか）、ハルシネーションの有無
\end{enumerate}

\putbib
\end{bibunit}

\begin{bibunit}
\chapter{AIを使った文献調査}
\label{chap3}

\begin{chapterbox}
\textbf{【この章で学ぶこと】}
\begin{itemize}
\item AIツールを活用した効率的な文献調査ワークフローを習得する --- 4ステップの体系的な調査手法を実践できるようになる
\item 情報の正確性を批判的に検証する方法を身につける --- AIが生成する情報の限界を理解し、ファクトチェックを行える
\item 研究の新規性を言語化する技法を理解する --- 先行研究との差分を明確に説明できるようになる
\end{itemize}
\end{chapterbox}
\vspace{5mm}

%=========================================
\section{研究調査の課題}
\label{sec:3.1}
%=========================================

\subsection{情報爆発の時代}

現代の研究者が直面する最大の課題は、学術情報の爆発的な増加である。プレプリントサーバーarXiv\index{arXiv}の投稿数は、2010年には年間約7万件であったが、2023年には年間約19万件に達し、わずか13年で約2.7倍に増加した。PubMedに登録される生物医学系論文も、年間約100万件に達している。

この情報爆発は、特に大学院生にとって深刻な問題である。自分の研究テーマに関連する論文をすべて読むことは物理的に不可能であり、重要な先行研究を見落とすリスクが常につきまとう。

\subsection{キーワードの壁}

文献調査\index{ぶんけんちょうさ@文献調査}で見落としが発生する大きな原因の一つは、「キーワードの壁」である。同じ概念でも分野によって異なる用語が使われることがある。例えば、以下のような例がある。

\begin{itemize}
\item 「転移学習（Transfer Learning）」と「ドメイン適応（Domain Adaptation）」
\item 「特徴量エンジニアリング（Feature Engineering）」と「表現学習（Representation Learning）」
\item 「異常検知（Anomaly Detection）」と「外れ値検出（Outlier Detection）」と「新奇性検出（Novelty Detection）」
\end{itemize}

自分が知らない用語で記述されている重要な研究を見逃す可能性は、特に分野横断的な研究において顕著である。従来の検索エンジンはキーワード一致に依存するため、この問題を解決できなかった。

\subsection{時間的制約}

大学院生の日常は、授業、研究、TA業務、学会準備など多忙を極める。その中で英語論文を精読する時間を確保することは容易ではない。1本の論文を深く理解するのに数時間かかることも珍しくなく、関連論文を網羅的に調査するには膨大な時間が必要である。

特に英語が母語でない日本語話者にとって、英語論文の読解には追加の負荷がかかる。専門用語の理解、複雑な文構造の把握、文脈に依存した表現の解釈など、言語的な障壁が調査効率を低下させる。

\subsection{思考の停滞}

文献調査を進める中で、多くの大学院生が経験するのが「思考の停滞」である。先行研究を読めば読むほど、「すでに研究されている」「自分の研究に新規性がない」と感じてしまうことがある。

この問題の根本は、自分の研究の新規性\index{しんきせい@新規性}をうまく言語化できないことにある。研究の新規性は、先行研究との差分を明確にすることで初めて可視化される。しかし、その差分を的確に表現する技術は、経験を積まなければ身につかないものであった。AIはこの言語化のプロセスを支援する。

%=========================================
\section{文献調査AIツールの使い分け}
\label{sec:3.2}
%=========================================

\subsection{対話型AI --- ブレスト・壁打ちの相棒}

\textbf{主なツール:} Gemini Pro、ChatGPT（GPT-4o）、Claude

対話型AIは、研究のアイデア出しや方向性の検討に最適である。人間の指導教員や研究仲間と壁打ちするように、AIと対話することで思考を整理できる。

\textbf{長所:}
\begin{itemize}
\item 自然言語で曖昧な質問ができる
\item 対話を通じて段階的に思考を深められる
\item 異なる視点からの提案を得られる
\item 複雑な概念の説明を求められる
\end{itemize}

\textbf{短所:}
\begin{itemize}
\item 出典が明示されないことが多い
\item 存在しない論文を生成する場合がある（ハルシネーション）
\item 学習データの時期によって最新情報が反映されない
\item 回答の正確性を自分で検証する必要がある
\end{itemize}

\subsection{検索特化型AI --- 出典付き回答でサーベイの初動}

\textbf{主なツール:} Perplexity\index{Perplexity}

検索特化型AIは、ウェブ検索とAIの回答生成を組み合わせたツールである。回答に出典リンクが付与されるため、情報の検証が容易である。

\textbf{長所:}
\begin{itemize}
\item 回答に出典URLが付く
\item リアルタイムのウェブ情報を反映できる
\item サーベイの初期段階で全体像を掴むのに適している
\item 英語と日本語の両方で検索可能
\end{itemize}

\textbf{短所:}
\begin{itemize}
\item 深い分析や批判的検討は苦手
\item 出典の質にばらつきがある（学術論文以外の情報も混在）
\item 対話の文脈保持能力は対話型AIに劣る
\end{itemize}

\subsection{論文分析AI --- 深掘りのための専門ツール}

\textbf{主なツール:} NotebookLM\index{NotebookLM}、Elicit\index{Elicit}、Connected Papers\index{Connected Papers}、Semantic Scholar\index{Semantic Scholar}

論文分析に特化したAIツールは、学術文献の深い理解を支援する。

\begin{table}[htbp]
\centering
\caption{論文分析AIツールの比較}
\label{tab:paper_analysis_tools}
\begin{tabular}{lll}
\toprule
ツール名 & 特徴 & 最適な用途 \\
\midrule
NotebookLM & アップロードした文書をもとに質問応答 & 特定論文の深い理解 \\
Elicit & 学術論文の検索・要約に特化 & 体系的な文献レビュー \\
Connected Papers & 引用関係の可視化 & 研究分野のマッピング \\
Semantic Scholar & AI支援の論文検索エンジン & 関連論文の網羅的発見 \\
\bottomrule
\end{tabular}
\end{table}

\subsection{目的別フローチャート}

文献調査の目的に応じて、最適なツールの組み合わせは異なる。以下に目的別の推奨フローを示す。

\begin{itembox}[l]{目的別の推奨フロー}
\textbf{パターンA: 新しい研究テーマの探索}
\begin{enumerate}
\item 対話型AI（Claude/ChatGPT）でブレインストーミング
\item Perplexityで関連分野の概要を把握
\item Connected Papersで研究分野のマッピング
\end{enumerate}

\vspace{3mm}
\textbf{パターンB: 既存テーマの先行研究調査}
\begin{enumerate}
\item Perplexityで主要な先行研究を特定
\item Semantic ScholarまたはElicitで網羅的な検索
\item NotebookLMで重要論文を深掘り
\item 対話型AIで新規性の言語化
\end{enumerate}

\vspace{3mm}
\textbf{パターンC: 特定論文の理解}
\begin{enumerate}
\item NotebookLMに論文をアップロードして質問
\item 対話型AIに要約・解説を依頼
\item Connected Papersで関連研究を確認
\end{enumerate}
\end{itembox}

%=========================================
\section{4ステップの文献調査ワークフロー}
\label{sec:3.3}
%=========================================

\subsection{Step 1: キーワード拡張 --- AIで関連キーワードを発見する}

文献調査の第一歩は、検索に使うキーワードの拡張\index{きーわーどかくちょう@キーワード拡張}である。自分が知っているキーワードだけでは、関連する研究を見落とす可能性がある。AIを使って関連キーワードを網羅的に発見しよう。

\begin{promptbox}
\textbf{【基本プロンプト例：キーワード拡張】}

私は「[研究テーマ]」について文献調査を行いたいと考えています。この研究テーマに関連する検索キーワードを、以下の観点から網羅的にリストアップしてください。

\begin{enumerate}
\item 同義語・類義語（英語・日本語両方）
\item 上位概念・下位概念
\item 関連する手法名・アルゴリズム名
\item 関連する応用分野
\item この分野で頻出するキーワード
\end{enumerate}

それぞれのキーワードについて、簡単な説明も付けてください。
\end{promptbox}

\begin{promptbox}
\textbf{【具体例 --- 「グラフニューラルネットワークによる薬物相互作用予測」の場合】}

私は「グラフニューラルネットワークによる薬物相互作用予測」について文献調査を行いたいと考えています。この研究テーマに関連する検索キーワードを、以下の観点から網羅的にリストアップしてください。

\begin{enumerate}
\item 同義語・類義語（英語・日本語両方）
\item 上位概念・下位概念
\item 関連する手法名・アルゴリズム名
\item 関連する応用分野
\item この分野で頻出するキーワード
\end{enumerate}

それぞれのキーワードについて、簡単な説明も付けてください。
\end{promptbox}

\textbf{出力の評価と活用法:}

AIが生成したキーワードリストは、そのまま使うのではなく、以下の手順で評価・活用する。

\begin{enumerate}
\item \textbf{知っているキーワードの確認} --- リストに自分が既知のキーワードが含まれているか確認する。含まれていない場合、プロンプトの改善が必要
\item \textbf{未知のキーワードの学習} --- 知らなかったキーワードについて、追加で調べる
\item \textbf{キーワードの優先順位付け} --- 自分の研究テーマとの関連度が高い順に並べ替える
\item \textbf{検索クエリの構築} --- 優先度の高いキーワードを組み合わせて、Google ScholarやSemantic Scholarの検索クエリを作成する
\end{enumerate}

\begin{screen}
\textbf{ヒント:} キーワード拡張は一度で完了するものではない。調査を進める中で新たなキーワードが見つかったら、このステップに戻って検索範囲を広げよう。反復的な調査が網羅性を高める。
\end{screen}

\subsection{Step 2: 全体像を掴むサーベイ --- Perplexity等で俯瞰する}

キーワードが揃ったら、次は研究分野の全体像を掴むサーベイ\index{さーべい@サーベイ}を行う。ここではPerplexityのような検索特化型AIが効果的である。

\paragraph{英語プロンプトの重要性}
学術情報の大部分は英語で公開されている。そのため、サーベイのプロンプトは英語で入力することを強く推奨する。日本語で入力すると、日本語のウェブページが優先的に検索され、最新の学術論文が見落とされる可能性がある。

\begin{promptbox}
\textbf{【英語でのサーベイプロンプト例】}

What are the current state-of-the-art approaches for drug-drug interaction prediction using graph neural networks? Please provide:

\begin{enumerate}
\item A brief overview of the field (2--3 paragraphs)
\item Key methods and their characteristics (with citations)
\item Major datasets used in this field
\item Open challenges and future directions
\item Key research groups and their contributions
\end{enumerate}
\end{promptbox}

\textbf{出典リンクの確認方法:}

Perplexityが提示する出典は、必ず以下の手順で確認する。

\begin{enumerate}
\item 出典リンクをクリックして、実際のページが存在するか確認する
\item 論文の場合、DOIまたはarXiv IDが正しいか確認する
\item Google Scholarで論文タイトルを検索し、被引用数や掲載ジャーナルを確認する
\item 情報が一次情報源に基づいているか確認する（ブログ記事やニュース記事ではなく、論文原文を参照する）
\end{enumerate}

\begin{screen}
\textbf{ヒント:} Perplexityの「Academic」モードを使うと、学術論文を優先的に検索してくれる。サーベイの初期段階では、このモードを活用しよう。
\end{screen}

\subsection{Step 3: 注目論文の深掘り --- アブストラクトから構造的要約へ}

全体像を掴んだら、特に重要と思われる論文を深く読み込む。AIを活用して、効率的に論文の要点を抽出しよう。

\begin{promptbox}
\textbf{【4点要約テンプレート】}

以下の論文のアブストラクト（および可能であれば本文）を読み、次の4点について日本語で構造的に要約してください。

【論文タイトル】: [タイトルを入力]

【アブストラクト】: [アブストラクトを貼り付け]

\begin{enumerate}
\item \textbf{背景（Background）}: この研究が取り組む問題は何か？なぜ重要か？
\item \textbf{手法（Method）}: どのようなアプローチ・手法を提案しているか？技術的な新規性は何か？
\item \textbf{結果（Results）}: 主要な実験結果は何か？既存手法と比較してどの程度の改善が得られたか？
\item \textbf{限界（Limitations）}: この研究の限界や未解決の課題は何か？今後の研究の方向性として何が示唆されているか？
\end{enumerate}

また、この論文の新規性を一文で表現してください。
\end{promptbox}

\paragraph{複数論文の比較要約}
重要な論文が複数見つかった場合、AIを使って比較表を作成できる。

\begin{promptbox}
\textbf{【複数論文の比較要約プロンプト】}

以下の3つの論文について、比較表を作成してください。

論文1: [タイトル] --- [アブストラクト]

論文2: [タイトル] --- [アブストラクト]

論文3: [タイトル] --- [アブストラクト]

比較の軸:
\begin{itemize}
\item 解決しようとしている問題
\item 提案手法の特徴
\item 使用データセット
\item 主要な結果（定量的な指標を含む）
\item 限界と課題
\item 自分の研究テーマ「[テーマ]」との関連性
\end{itemize}
\end{promptbox}

\begin{screen}
\textbf{ヒント:} NotebookLMを使えば、論文PDFを直接アップロードして質問できる。アブストラクトだけでなく、実験設定の詳細や図表の解釈についても質問できるため、深い理解が得られる。
\end{screen}

\subsection{Step 4: 研究の新規性を言語化する}

文献調査の最終目標は、自分の研究の新規性を明確に言語化することである。このステップが最も重要であり、同時に最も難しいステップである。

\begin{promptbox}
\textbf{【新規性言語化のためのプロンプト】}

私の研究テーマは「[テーマ]」です。以下に、関連する先行研究とその限界をまとめました。

先行研究1: [タイトル] / 手法: [手法の概要] / 限界: [限界]

先行研究2: [タイトル] / 手法: [手法の概要] / 限界: [限界]

先行研究3: [タイトル] / 手法: [手法の概要] / 限界: [限界]

私が提案する手法は「[自分の手法の概要]」です。

以下の点について、整理を手伝ってください。
\begin{enumerate}
\item 先行研究に共通する未解決の課題は何か？
\item 私の提案手法は、それらの課題をどのように解決するか？
\item 私の研究の新規性を、学術論文のIntroductionに書くような形式で3--4文にまとめてください。
\item この新規性の主張に対して、査読者が指摘しそうな弱点は何か？
\end{enumerate}
\end{promptbox}

\paragraph{差分マトリクスの作成}

\begin{promptbox}
\textbf{【差分マトリクス作成プロンプト】}

以下の先行研究と私の研究について、差分マトリクス\index{さぶんまとりくす@差分マトリクス}を作成してください。

比較軸として以下を含めてください:
\begin{itemize}
\item 対象とする問題設定
\item 使用するデータの種類
\item モデルアーキテクチャ
\item 学習方法
\item 評価指標
\item 計算コスト
\item 実用性
\end{itemize}

各セルには、具体的な違いを簡潔に記述してください。自分の研究の優位点には$\bigstar$をつけてください。
\end{promptbox}

\begin{screen}
\textbf{ヒント:} 新規性の言語化は、AIの出力をそのまま使うのではなく、あくまで「たたき台」として活用すること。最終的な新規性の主張は、指導教員や研究仲間と議論して磨き上げることが重要である。AIはあくまで思考の補助であり、研究の方向性を決めるのは自分自身である。
\end{screen}

%=========================================
\section{情報の批判的検証}
\label{sec:3.4}
%=========================================

\subsection{AIが出力した論文情報のファクトチェック手法}

AIは非常に有用なツールであるが、その出力を無批判に信頼することは危険である。特に文献調査においては、以下のような問題が発生する可能性がある。

\begin{itembox}[l]{ファクトチェック\index{ふぁくとちぇっく@ファクトチェック}のチェックリスト}
\begin{itemize}
\item[$\square$] 論文のタイトルと著者名が正確か（Google Scholarで検索）
\item[$\square$] 掲載ジャーナル/学会名が正しいか
\item[$\square$] 発表年が正確か
\item[$\square$] 引用されている数値や結果が論文原文と一致するか
\item[$\square$] 論文の主張がAIの要約と整合しているか
\end{itemize}
\end{itembox}

\textbf{具体的なファクトチェック手順:}

\begin{enumerate}
\item \textbf{Google Scholarでの検索:} AIが挙げた論文のタイトルをそのままGoogle Scholarに入力し、実在する論文かどうかを確認する
\item \textbf{DOIの確認:} DOIが提示されている場合、doi.orgで解決できるか確認する
\item \textbf{著者のホームページ確認:} 主要著者の研究室ページや個人ページで、該当論文が掲載されているか確認する
\item \textbf{数値の照合:} 重要な実験結果の数値は、必ず論文原文で確認する
\end{enumerate}

\subsection{存在しない論文への対処（ハルシネーション問題）}

AIが存在しない論文を「作り上げる」ことは、ハルシネーション\index{はるしねーしょん@ハルシネーション}の典型例である。もっともらしいタイトル、著者名、ジャーナル名が生成されるが、実際には存在しない論文であるケースが少なくない。

\textbf{ハルシネーションを疑うべきサイン:}
\begin{itemize}
\item タイトルが非常に一般的で、どの分野にも当てはまりそうな表現
\item 有名な研究者の名前が著者に含まれているが、その研究者の専門分野と一致しない
\item 掲載ジャーナルが非常に有名だが、分野が異なる
\item 発表年が最新すぎる、または古すぎる
\item DOIやURLが提示されない
\end{itemize}

\begin{promptbox}
\textbf{【ハルシネーション対処のプロンプト例】}

先ほど挙げていただいた論文について確認したいのですが、以下の論文はGoogle Scholarで見つかりませんでした。

「[論文タイトル]」

この論文は実在しますか？もし存在しない場合、同様の内容を扱う実在する論文を教えてください。その際、Google Scholarで確認できるように、正確なタイトルと著者名を記載してください。
\end{promptbox}

\begin{screen}
\textbf{ヒント:} Perplexityは出典リンク付きで回答を生成するため、ハルシネーションのリスクが比較的低い。論文の実在性が重要な場合は、Perplexityで確認するか、直接Google ScholarやSemantic Scholarで検索しよう。
\end{screen}

\subsection{バイアスの認識}

AIの回答には、学習データに起因するバイアス\index{ばいあす@バイアス}が含まれる可能性がある。文献調査において注意すべきバイアスには以下がある。

\paragraph{英語圏バイアス}
AIは英語のデータで主に学習されているため、英語圏の研究が過度に強調される傾向がある。日本語や他の言語で発表された重要な研究が見落とされる可能性がある。

\paragraph{引用数バイアス}
被引用数が多い「有名な」論文が優先的に紹介される傾向があり、最新の研究や小規模な研究グループの重要な成果が見落とされる可能性がある。

\paragraph{分野バイアス}
AIの学習データに含まれる論文の分野分布に偏りがあるため、特定の分野（例: 機械学習、自然言語処理）の研究が過度に強調される可能性がある。

\paragraph{出版バイアス}
肯定的な結果を報告する論文が優先される傾向があり、否定的な結果や再現性の問題を報告する論文が見落とされる可能性がある。

これらのバイアスを認識した上で、AIの出力を批判的に評価することが重要である。

%=========================================
\section{文献管理とまとめ}
\label{sec:3.5}
%=========================================

\subsection{調査結果の整理法}

文献調査で得られた情報は、体系的に整理することで初めて研究に活用できる。AIを使って、調査結果を表形式で整理しよう。

\begin{promptbox}
\textbf{【文献比較表の作成プロンプト】}

以下の論文情報をもとに、比較表を作成してください。

[論文情報をリストアップ]

表の列には以下を含めてください:
\begin{itemize}
\item 論文（著者, 年）
\item 研究目的
\item 提案手法
\item データセット
\item 主要結果
\item 限界・課題
\item 自分の研究との関連
\end{itemize}
\end{promptbox}

整理のポイントとして、以下の項目を意識すること。

\begin{itemize}
\item \textbf{時系列での整理:} 研究の発展の流れを把握するため、発表年順に並べる
\item \textbf{手法別の整理:} 異なるアプローチを分類し、各手法の特徴を比較する
\item \textbf{問題設定別の整理:} 同じ問題に取り組む研究をグループ化する
\end{itemize}

\subsection{AIで文献レビューの草稿を作成する}

文献調査の結果を論文の「関連研究」セクションにまとめる作業も、AIの支援を受けることができる。

\begin{promptbox}
\textbf{【文献レビュー草稿作成プロンプト】}

以下の先行研究の情報をもとに、論文の``Related Work''セクションの日本語草稿を作成してください。

[先行研究の要約リスト]

以下の要件を満たしてください:
\begin{enumerate}
\item 研究を意味のあるカテゴリに分類して整理する
\item 各カテゴリ内で、研究の発展を時系列で説明する
\item 各研究の貢献と限界を簡潔に述べる
\item 最後に、これらの先行研究と本研究の関係を明確にする
\item 学術論文にふさわしい客観的な文体を使用する
\end{enumerate}
\end{promptbox}

\begin{screen}
\textbf{ヒント:} AIが生成した文献レビューの草稿は、あくまで出発点である。必ず自分で読み直し、以下を確認すること。(1) 論文の内容が正確に要約されているか、(2) 研究間の関係が正しく記述されているか、(3) 自分の研究の位置づけが明確か。
\end{screen}

\subsection{文献管理ツールとの連携}

AIで発見した論文は、文献管理ツール\index{ぶんけんかんりつーる@文献管理ツール}に登録して一元管理しよう。主要な文献管理ツールとその特徴は表\ref{tab:reference_tools}の通りである。

\begin{table}[htbp]
\centering
\caption{主要な文献管理ツールの比較}
\label{tab:reference_tools}
\begin{tabular}{lll}
\toprule
ツール名 & 特徴 & 料金 \\
\midrule
Zotero\index{Zotero} & オープンソース、ブラウザ拡張が便利 & 無料（300MBまで） \\
Mendeley\index{Mendeley} & PDF管理が優秀、Elsevierとの連携 & 無料（2GBまで） \\
Paperpile & Google Docsとの連携が強力 & 有料（学生割引あり） \\
EndNote\index{EndNote} & 多くの大学で機関ライセンスあり & 有料 \\
\bottomrule
\end{tabular}
\end{table}

\begin{itembox}[l]{AIと文献管理ツールの連携ワークフロー}
\begin{enumerate}
\item AIで文献調査を行い、重要な論文を特定する
\item Google ScholarやSemantic Scholarで論文を確認し、文献管理ツールに登録する
\item 文献管理ツールでタグ付け・メモの追加を行う
\item 論文執筆時に、文献管理ツールから引用を挿入する
\end{enumerate}
\end{itembox}

\begin{screen}
\textbf{ヒント:} ZoteroのブラウザExtensionを使えば、Google ScholarやarXivのページからワンクリックで論文を登録できる。AIで見つけた論文を効率的に管理するためにも、文献管理ツールの導入を強く推奨する。
\end{screen}

%=========================================
\section*{演習問題}
\addcontentsline{toc}{section}{演習問題}
%=========================================

\begin{enumerate}
\item \textbf{演習1: キーワード拡張の実践}

自分の研究テーマ（または興味のある研究トピック）について、AIを使ってキーワード拡張を行いなさい。以下の手順に従うこと。

\begin{enumerate}
\item 対話型AIに研究テーマを説明し、関連キーワードを少なくとも20個生成させる
\item 生成されたキーワードを「知っていたもの」と「知らなかったもの」に分類する
\item 知らなかったキーワードのうち、最も重要だと思うもの3つについて、Google Scholarで検索して関連論文を見つける
\item この演習を通じて発見した新しい視点や気づきをまとめる
\end{enumerate}

\item \textbf{演習2: ファクトチェックの実践}

以下の手順で、AIのハルシネーション検出を体験しなさい。

\begin{enumerate}
\item 対話型AIに「[自分の研究分野]における重要な論文を10本挙げてください」と質問する
\item 挙げられた10本の論文すべてについて、Google Scholarで実在を確認する
\item 存在しない論文があった場合、その数と割合を記録する
\item ハルシネーションが発生した論文について、どのような特徴があったかを分析する
\item 結果をもとに、AIの出力を検証する際の注意点をまとめる
\end{enumerate}

\item \textbf{演習3: 4点要約の実践}

自分の研究テーマに関連する英語論文を1本選び、以下を行いなさい。

\begin{enumerate}
\item まず自分でアブストラクトを読み、4点（背景・手法・結果・限界）を日本語でまとめる
\item 次に、本章のテンプレートを使ってAIに4点要約を依頼する
\item 自分の要約とAIの要約を比較し、(a) AIが正確に要約できた点、(b) AIが見落とした点、(c) AIが誤って解釈した点、をそれぞれ挙げる
\item この比較を通じて、AI要約を活用する際のベストプラクティスをまとめる
\end{enumerate}

\item \textbf{演習4: 新規性の言語化}

以下の手順で、自分の研究の新規性を言語化しなさい。

\begin{enumerate}
\item 自分の研究テーマに最も近い先行研究を3本特定する
\item 各先行研究について、手法と限界を整理する
\item 本章のプロンプトテンプレートを使って、AIに新規性の言語化を依頼する
\item AIの出力を批判的に検討し、修正・改善する
\item 最終的な新規性の記述を、指導教員やゼミのメンバーに共有してフィードバックを得る
\end{enumerate}

\item \textbf{演習5: 文献レビューの草稿作成}

自分の研究テーマについて、以下の手順で文献レビューの草稿を作成しなさい。

\begin{enumerate}
\item これまでの演習で収集した文献情報をもとに、AIに文献レビューの草稿を作成させる
\item 草稿を読み、以下の観点で評価する: (a) 分類は適切か、(b) 時系列の流れは正確か、(c) 各研究の記述は正確か
\item 問題があった箇所を修正し、自分の研究の位置づけを加筆する
\item 完成した草稿を文献管理ツール（Zoteroなど）の引用情報と照合し、参考文献リストを作成する
\end{enumerate}
\end{enumerate}

\vspace{5mm}
\begin{screen}
\textbf{本章のまとめ:} AIは文献調査を効率化する強力なツールであるが、万能ではない。AIの出力を鵜呑みにせず、常にファクトチェックを行い、批判的思考を維持することが不可欠である。文献調査の目的は、先行研究を知ることだけでなく、自分の研究の新規性と価値を明確にすることにある。AIを「思考の補助輪」として活用し、自分自身の研究力を高めてほしい。
\end{screen}

\putbib
\end{bibunit}

\begin{bibunit}
% -*- coding: utf-8 -*-
\chapter{英語論文の執筆と推敲}
\label{chap4}

\begin{chapterbox}
\textbf{【この章で学ぶこと】}
\begin{itemize}
\item \textbf{英語論文の基本構造（IMRaD）を理解する} --- 各セクションの役割と書くべき内容を把握し、構造的な論文を書ける
\item \textbf{AIを活用した英訳・執筆・推敲ワークフローを習得する} --- 日本語原稿からAIの支援を受けて高品質な英語論文を仕上げるプロセスを実践できる
\item \textbf{学術的な英語表現を適切に使い分けられる} --- 分野に応じた定型表現を活用し、査読者に好印象を与える論文を執筆できる
\end{itemize}
\end{chapterbox}
\vspace{5mm}

%======================================================================
\section{なぜ英語論文が必要か}
\label{sec:4.1}
\index{えいごろんぶん@英語論文}

\subsection{国際会議・ジャーナルの重要性}
\index{こくさいかいぎ@国際会議}
\index{じゃーなる@ジャーナル}

研究の価値は、その成果が広く共有され、他の研究者に活用されて初めて発揮される。日本語で発表された研究は、国内の研究者には届くが、世界中の研究者にはほとんど届かない。トップカンファレンス（NeurIPS、ICML、ACL、CVPR、AAAI等）やトップジャーナル（Nature、Science、IEEE Transactions等）はすべて英語で発表されており、国際的な研究コミュニティの共通言語は英語である。

大学院生にとって英語論文の執筆は、以下の理由で極めて重要である。

\begin{itemize}
\item \textbf{研究成果の国際的な可視性:} 英語論文は世界中の研究者が読むことができ、被引用数の向上につながる
\item \textbf{キャリア形成:} 国際的な実績は、アカデミアでも産業界でも高く評価される
\item \textbf{研究コミュニティへの参加:} 英語論文を発表することで、国際的な研究ネットワークに参加できる
\item \textbf{フィードバックの質:} 国際会議の査読者は世界トップレベルの研究者であり、質の高いフィードバックが得られる
\end{itemize}

\subsection{研究成果の国際発信と評価}

英語論文の執筆は、単に日本語を英語に翻訳する作業ではない。国際的な読者に向けて、研究の背景、意義、手法、結果を論理的かつ説得力のある形で伝える作業である。

しかし、英語を母語としない研究者にとって、学術英語の壁は非常に高い。従来は英文校正サービスやネイティブの共同研究者に頼ることが一般的であったが、AIの登場により、英語論文執筆のハードルは大きく下がった。AIは万能ではないが、適切に活用すれば、執筆の効率と品質を大幅に向上させられる。

%======================================================================
\section{英語論文の基本構造}
\label{sec:4.2}

\subsection{IMRaD構造}
\index{IMRaD@IMRaD構造}

国際的な学術論文の標準的な構造はIMRaDと呼ばれ、以下の4つのセクションから構成される。

\begin{table}[htbp]
\centering
\caption{IMRaD構造の概要}
\label{tab:imrad}
\footnotesize
\begin{tabular}{llll}
\toprule
セクション & 英語名 & 主な内容 & 時制 \\
\midrule
はじめに & Introduction & 研究の背景、目的、貢献 & 現在形・過去形 \\
手法 & Methods & 提案手法の詳細 & 過去形 \\
結果 & Results & 実験結果の提示 & 過去形 \\
考察 & Discussion & 結果の解釈、限界、将来展望 & 現在形・過去形 \\
\bottomrule
\end{tabular}
\end{table}

これに加えて、以下のセクションが一般的に含まれる。

\begin{itemize}
\item \textbf{Abstract（要旨）:} 論文全体の要約（150--300語程度）
\item \textbf{Related Work（関連研究）:} 先行研究の整理と本研究の位置づけ
\item \textbf{Conclusion（結論）:} 研究の要約と将来展望
\item \textbf{References（参考文献）:} 引用した文献のリスト
\end{itemize}

\subsection{各セクションの役割と書くべき内容}

\begin{itembox}[l]{Introduction（はじめに）の構造}
Introductionは「逆三角形」の構造で書く。まず広い背景から始めて、徐々に焦点を絞り、最後に自分の研究の貢献を明確に述べる。
\begin{enumerate}
\item 研究分野の一般的な背景と重要性（1--2段落）
\item 既存研究の課題や未解決の問題（1--2段落）
\item 本研究のアプローチと貢献（1段落）
\item 論文の構成の概要（1段落、省略可）
\end{enumerate}
\end{itembox}

\textbf{Methods（手法）:}

提案手法を、読者が再現できるレベルで詳細に記述する。

\begin{enumerate}
\item 問題設定の形式的な定義
\item 提案手法の全体像
\item 各構成要素の詳細
\item 実装の詳細（ハイパーパラメータ、使用ライブラリ等）
\end{enumerate}

\textbf{Results（結果）:}

実験結果を客観的に提示する。解釈はDiscussionで行う。

\begin{enumerate}
\item 実験設定（データセット、評価指標、ベースライン）
\item 主要な結果（表やグラフ）
\item 結果の記述（数値の比較、傾向の指摘）
\end{enumerate}

\textbf{Discussion（考察）:}

結果の意味を解釈し、研究の意義と限界を議論する。

\begin{enumerate}
\item 主要な発見の解釈
\item 先行研究との比較
\item 研究の限界
\item 将来の研究方向
\end{enumerate}

\subsection{パラグラフの構成法}
\index{ぱらぐらふ@パラグラフ}

英語の学術論文では、パラグラフ（段落）が意味の基本単位である。各パラグラフは以下の構造に従う。

\begin{enumerate}
\item \textbf{トピックセンテンス:} パラグラフの主張を1文で述べる（段落の冒頭に置く）
\item \textbf{サポート:} トピックセンテンスを支える根拠、例、データを述べる（2--5文）
\item \textbf{結論/遷移:} パラグラフの内容をまとめ、次のパラグラフへつなげる（1文）
\end{enumerate}

\begin{promptbox}
以下の日本語の段落を、トピックセンテンス・サポート・結論の構造に
整理してください。必要に応じて、論理の飛躍を指摘してください。

{[}日本語の段落を貼り付け{]}
\end{promptbox}

\begin{itembox}[l]{Tip: 1パラグラフ1メッセージ}
英語論文では「1パラグラフ1メッセージ」の原則を守ることが重要である。複数のメッセージが混在するパラグラフは、読者の理解を妨げる。日本語の文章は1つの段落に複数の話題を含むことがあるが、英語論文ではそれぞれを独立したパラグラフに分けるのが原則である。
\end{itembox}

%======================================================================
\section{AIを使った英訳ワークフロー}
\label{sec:4.3}
\index{えいやくわーくふろー@英訳ワークフロー}

\subsection*{Step 1: 日本語原稿の準備と構造化}

AIに英訳を依頼する前に、日本語原稿の品質を高めることが重要である。曖昧な日本語からは曖昧な英語しか生まれない。

\textbf{日本語原稿の準備チェックリスト:}

\begin{itemize}
\item 各段落にトピックセンテンスがあるか
\item 論理の飛躍がないか
\item 主語と述語が明確か（日本語は主語省略が多いので注意）
\item 専門用語が一貫して使われているか
\item 図表への参照が正しいか
\end{itemize}

\textbf{構造化のためのプロンプト:}

\begin{promptbox}
以下の日本語原稿を、英語論文に翻訳する前の準備として、
構造的に整理してください。

具体的には:
1. 各段落のトピックセンテンスを明示してください
2. 論理の飛躍や不明確な箇所を指摘してください
3. 主語が省略されている文を特定し、主語を補ってください
4. 段落の分割・統合が必要な箇所を提案してください

{[}日本語原稿を貼り付け]
\end{promptbox}

\subsection*{Step 2: セクション単位でAIに英訳を依頼}

日本語原稿が整理できたら、セクション単位でAIに英訳を依頼する。論文全体を一度に訳すのではなく、セクション単位で行うことで、品質の管理がしやすくなる。

\textbf{効果的な英訳プロンプト:}

\begin{promptbox}
以下の日本語原稿を、学術英語に翻訳してください。

【条件】
- 分野: [例: コンピュータサイエンス / 機械学習]
- 対象読者: 国際会議の査読者（専門家）
- 文体: 学術論文（フォーマル）
- セクション: [例: Introduction / Methods / Results]
- 時制: [例: 過去形を基本とする（Methodsの場合）]

【注意事項】
- 以下の専門用語は指定の英語表現を使用してください:
  [用語1]: [英語表現1]
  [用語2]: [英語表現2]
- 受動態を過度に使用しないでください
- 1文が長くなりすぎないようにしてください（目安: 25語以内）

【日本語原稿】
{[}原稿を貼り付け]
\end{promptbox}

\begin{itembox}[l]{分野指定の重要性}
学術英語の表現は分野によって大きく異なる。例えば、同じ「提案する」でも分野によって適切な表現が異なる。
\begin{itemize}
\item コンピュータサイエンス: ``We propose...'' / ``We present...''
\item 医学: ``We hypothesize...'' / ``We examined...''
\item 物理学: ``We derive...'' / ``We demonstrate...''
\end{itemize}
\end{itembox}

\begin{itembox}[l]{Tip: 参考論文の活用}
翻訳の際は、自分の分野のトップ論文を2--3本手元に置いておき、表現や文体の参考にするとよい。AIが生成した英語が、その分野の慣習に合っているかを確認するための基準になる。
\end{itembox}

\subsection*{Step 3: 学術的表現への洗練}

初回の翻訳結果をさらに学術的に洗練する。

\begin{promptbox}
以下の英語論文の草稿を、より学術的で洗練された表現に改善してください。

【改善の方向性】
1. 曖昧な表現をより具体的・精確にする
2. 冗長な表現を簡潔にする
3. 接続語を適切に使い、論理の流れを明確にする
4. 学術論文で一般的な表現に置き換える
   （例: "very good" → "substantially improved"）
5. ヘッジング表現を適切に使用する
   （例: "This proves..." → "These results suggest..."）

改善箇所には【変更理由】を付けてください。

{[}英語草稿を貼り付け]
\end{promptbox}

\subsection*{Step 4: 文法・スタイルチェック}

最終段階として、文法とスタイルの統一性を確認する。

\begin{promptbox}
以下の英語論文について、文法・スタイルの最終チェックを行ってください。

【チェック項目】
1. 文法エラー（主語-動詞の一致、冠詞、前置詞など）
2. 時制の一貫性（セクション内で統一されているか）
3. 用語の一貫性（同じ概念に同じ用語が使われているか）
4. 数値の表記（有効数字、単位の表記）
5. 参照の形式（Figure 1、Table 2などの表記統一）
6. スペルの統一（British/American Englishの混在がないか）

エラーが見つかった場合、修正案と理由を示してください。

{[}英語論文を貼り付け]
\end{promptbox}

%======================================================================
\section{推敲テクニック}
\label{sec:4.4}
\index{すいこう@推敲}

\subsection{段階的推敲法}

英語論文の推敲は、一度に全てをチェックしようとすると、多くの問題を見落とす。以下の3段階に分けて、段階的に推敲を行うとよい。

\textbf{第1段階: 文法チェック}

\begin{promptbox}
以下の英語論文について、文法エラーのみに焦点を当ててチェックしてください。
表現の改善提案は不要です。純粋な文法エラー（主語-動詞の一致、冠詞、
前置詞、時制の誤りなど）のみを指摘し、修正案を示してください。

{[}英語論文を貼り付け]
\end{promptbox}

\textbf{第2段階: 表現チェック}

\begin{promptbox}
以下の英語論文について、表現の改善に焦点を当ててチェックしてください。
文法は正しいが、以下の観点で改善できる箇所を指摘してください。

1. より学術的な表現への置き換え
2. 冗長な表現の簡潔化
3. 曖昧な表現の明確化
4. 不自然な表現（日本語直訳調）の修正

{[}英語論文を貼り付け]
\end{promptbox}

\textbf{第3段階: 論理チェック}

\begin{promptbox}
以下の英語論文について、論理構造に焦点を当ててチェックしてください。

1. 段落間の論理的つながりは適切か
2. 主張に対する根拠は十分か
3. 飛躍や矛盾はないか
4. 読者が混乱する可能性のある箇所はどこか

{[}英語論文を貼り付け]
\end{promptbox}

\subsection{比較推敲法}

複数のAIツールで翻訳を行い、最良の表現を選ぶ方法である。

\begin{promptbox}
以下の日本語を3つの異なるスタイルで英訳してください。

スタイルA: 簡潔で直接的な表現
スタイルB: 詳細で丁寧な表現
スタイルC: [分野名]の国際会議で一般的な表現

それぞれの訳について、長所と短所を説明してください。

{[}日本語原稿を貼り付け]
\end{promptbox}

さらに、Claude、ChatGPT、Geminiなど異なるAIで同じ文章を翻訳し、表現を比較することも効果的である。各AIには得意な表現スタイルがあり、それぞれの良い部分を組み合わせることで、より高品質な英語論文を作成できる。

\subsection{逆翻訳チェック}
\index{ぎゃくほんやくちぇっく@逆翻訳チェック}

英訳が原文の意味を正確に反映しているかを確認する強力な方法が、逆翻訳チェックである。

\begin{promptbox}
以下の英語を日本語に翻訳してください。
できるだけ直訳に近い形で翻訳し、英語の構造がわかるようにしてください。

{[}英訳された論文の一部を貼り付け]
\end{promptbox}

逆翻訳された日本語を元の日本語原稿と比較し、意味のズレがないかを確認する。ズレが見つかった場合、英訳を修正する。

\textbf{逆翻訳で発見できる問題:}

\begin{itemize}
\item 原文にあった情報が英訳で欠落している
\item 英訳が原文とは異なる意味に解釈される
\item 曖昧な英語表現が原文の意図と異なる解釈を生んでいる
\item 専門用語の翻訳が不正確である
\end{itemize}

\subsection{ネイティブチェック風レビュー}

AIに査読者やネイティブスピーカーの役割を担わせることで、論文の品質を事前に確認できる。

\begin{promptbox}
あなたは英語を母語とする[分野名]の研究者です。
以下の英語論文を読み、非ネイティブ話者が書いたと思われる
不自然な表現を指摘してください。

指摘の際は以下の形式でお願いします:
- 原文: [不自然な表現]
- 修正案: [自然な表現]
- 理由: [なぜ不自然か、どう改善されるか]

特に以下の点に注意してください:
1. 冠詞の使い方
2. 前置詞の選択
3. コロケーション（単語の自然な組み合わせ）
4. 文のリズムと読みやすさ

{[}英語論文を貼り付け]
\end{promptbox}

\begin{itembox}[l]{Tip: 推敲のタイミング}
推敲は時間をおいて行うのが効果的である。書いた直後ではなく、翌日に見直すと、問題点が見つかりやすくなる。AIを使った推敲も同様に、初回の英訳から時間をおいてから行うとよい。新鮮な目で見ることで、より多くの改善点を発見できる。
\end{itembox}

%======================================================================
\section{Abstract作成の技法}
\label{sec:4.5}
\index{Abstract@Abstract}

\subsection{Abstractの構造}

Abstractは論文全体の縮図であり、読者が論文を読むかどうかを決める最も重要な部分である。以下の5要素を含めるのが標準的である。

\begin{enumerate}
\item \textbf{背景（Background）:} 研究分野の文脈と解決すべき問題（1--2文）
\item \textbf{目的（Objective）:} 本研究で何を達成するか（1文）
\item \textbf{手法（Method）:} 提案するアプローチの概要（2--3文）
\item \textbf{結果（Results）:} 主要な実験結果の数値（1--2文）
\item \textbf{意義（Significance）:} 研究の意義と影響（1文）
\end{enumerate}

\subsection{150--200語にまとめるプロンプト設計}

\begin{promptbox}
以下の情報をもとに、学術論文のAbstractを英語で作成してください。

【研究情報】
- 研究分野: [分野]
- 解決する問題: [問題の記述]
- 提案手法: [手法の概要]
- 主要な結果: [数値を含む結果]
- 研究の意義: [意義の記述]

【条件】
- 語数: 150-200語
- 構造: 背景→目的→手法→結果→意義の順
- 時制: 背景は現在形、手法と結果は過去形、意義は現在形
- 具体的な数値を必ず含める
- 略語を初出で定義する
\end{promptbox}

\subsection{良いAbstractと悪いAbstractの比較}

\textbf{悪い例（問題点を含む）:}

\begin{quote}
In this paper, we propose a new method for text classification. Our method uses deep learning. We conducted experiments and obtained good results. Our method outperformed baseline methods. This work contributes to the field of natural language processing.
\end{quote}

\textbf{問題点:}
\begin{itemize}
\item 具体性がない（``new method''、``good results''）
\item 数値が一切含まれていない
\item 手法の特徴が不明
\item 読者にとって有益な情報がほとんどない
\end{itemize}

\textbf{良い例:}

\begin{quote}
Text classification in low-resource languages remains challenging due to the scarcity of labeled training data. We propose CrossLingual-BERT, a cross-lingual transfer learning framework that leverages pre-trained multilingual representations to classify texts in languages with limited annotated corpora. Our approach fine-tunes a multilingual BERT model on high-resource source languages and adapts it to target languages using only 100 labeled examples per class. Experiments on five typologically diverse languages demonstrate that CrossLingual-BERT achieves an average F1 score of 87.3\%, outperforming monolingual baselines by 12.5 percentage points and the previous state-of-the-art cross-lingual method by 4.2 percentage points. These results suggest that effective cross-lingual transfer is possible even with minimal target-language supervision, opening new possibilities for NLP applications in under-resourced languages.
\end{quote}

\textbf{優れている点:}
\begin{itemize}
\item 背景が具体的で問題が明確
\item 手法の核心的な特徴が述べられている
\item 具体的な数値が含まれている
\item 研究の意義が最後に述べられている
\end{itemize}

\begin{itembox}[l]{Tip: Abstractの執筆タイミング}
Abstractは論文の最後に書くのが効率的である。全セクションを書き終えてからAbstractを作成すれば、研究の全体像を正確に反映できる。AIにAbstractの草稿を作成させる際も、論文の完成後に全セクションの内容を入力として与えるとよい。
\end{itembox}

%======================================================================
\section{よくあるミスと対策}
\label{sec:4.6}

本節では、日本語話者が英語論文を書く際に陥りやすいミスとその対策を、授業で扱った具体例とともに体系的に解説する。まず各トピックの概念的な基礎を理解し、その後に早見表を参照するという流れで学習を進めてほしい。

\begin{screen}
\textbf{松野先生の教え：} 結局、日本語（母国語）の能力が重要である。必要最小限の日本語にしてから英語にする。これが技術英語の基本である。
\end{screen}

\subsection{すっきり最小主義：冗長な日本語を削ぎ落とす}
\index{すっきりさいしょうしゅぎ@すっきり最小主義}

英語に翻訳する前に、まず日本語を簡潔にすることが最も重要なステップである。冗長な日本語からは冗長な英語しか生まれない。

\subsubsection{重複表現を削る}

日本語は同じ意味の言葉を重ねる傾向がある。英訳前にこれを削ぎ落とす。

\begin{itembox}[l]{例：重複表現の除去}
\textbf{原文：} 「リチャージャブルバッテリーの安定した充電のためには、電気が流れる端子部が清浄であることが重要です」\\
\textbf{問題：} 「リチャージャブル（再充電可能な）」と「充電」が重複。「電気が流れる」と「端子」も重複。\\
\textbf{整理後：} 「安定した充電のためには端子を清潔に保つこと」\\
\textbf{英訳：} Clean the terminals for stable charging.
\end{itembox}

\subsubsection{空っぽ表現を捨てる}

日本語には意味のない「飾り」表現が多い。これを除去してから英訳する。

\begin{itembox}[l]{例：空っぽ表現の除去}
\textbf{原文：} 「オフィスでの集中力の妨げの原因となる不快な騒音を排除することで、生産性の向上が期待できるようになる」\\
\textbf{整理：} 「妨げの原因となる」→「妨げる」、「期待できるようになる」→不要\\
\textbf{整理後：} 「騒音を排除すれば生産性が向上する」\\
\textbf{英訳：} Eliminating noise improves productivity.
\end{itembox}

確信度を調整したい場合は、助動詞を使う。
\begin{itemize}
\item 弱い主張（5\%〜10\%程度）: Eliminating noise \textit{can} improve productivity.
\item 中程度の主張: Eliminating noise \textit{may} improve productivity.
\item 強い主張: Eliminating noise \textit{will} improve productivity.
\end{itemize}

\subsubsection{無駄な丁寧表現を除去する}

日本語の丁寧表現は英語に直訳すると不自然になる。

\begin{itembox}[l]{例：丁寧表現の除去}
\textbf{原文：} 「本日、無事に設計の最終確認を終了することができました」\\
\textbf{問題：} 「本日」「無事に」「終了することができました」は英語では不要\\
\textbf{整理後：} 「設計の最終確認を完了した」\\
\textbf{英訳：} We completed the final design review.
\end{itembox}

\begin{promptbox}
以下の日本語から、重複表現・空っぽ表現・無駄な丁寧表現を\\
取り除き、簡潔な日本語に整理してください。\\
その後、整理した日本語を英訳してください。\\
\\
{[}日本語原文を貼り付け{]}
\end{promptbox}

%---
\subsection{まず具体化主義：曖昧な表現を具体的にする}
\index{ぐたいかしゅぎ@具体化主義}

日本語の曖昧な表現は、英語にすると意味が通らなくなる。英訳前に具体的な表現に置き換える。

\begin{itembox}[l]{例：曖昧表現の具体化}
\textbf{原文：} 「オイル量のチェックはできるだけ頻繁に行うことが望ましい」\\
\textbf{問題：} 「できるだけ頻繁に」は英語にしにくい（as frequently as possible は不自然）\\
\textbf{具体化：} 「オイル量は1日1回点検すること」\\
\textbf{英訳：} Check the oil level once a day.
\end{itembox}

数値の曖昧さにも注意が必要である。

\begin{itembox}[l]{例：数値の具体化}
\textbf{原文：} 「生産は数千万だ」\\
\textbf{問題：} 「数千万」は英語でどう訳すか不明確\\
\textbf{具体化：} 「年間生産台数は3,600万台である」\\
\textbf{英訳：} Annual production is 36 million units.
\end{itembox}

略語も曖昧表現の一種である。初出時には必ずスペルアウトする。

\begin{itembox}[l]{例：略語のスペルアウト}
\textbf{推奨：} an antilock braking system (ABS)\\
\textbf{許容：} an ABS (antilock braking system)\\
\textbf{不可：} いきなり ABS と書く（読者が意味を知らない可能性がある）
\end{itembox}

%---
\subsection{とことん動詞主義：名詞句を動詞に戻す}
\index{どうししゅぎ@動詞主義}

日本語では「〜の実施」「〜の確認」「〜を行う」のように動詞を名詞化する傾向がある。これをそのまま英訳すると、名詞句が連鎖した読みにくい英文になる。動詞に戻してから英訳する。

\begin{itembox}[l]{例：名詞句 → 動詞}
\textbf{原文：} 「付録1はユニットの分解手順の説明である」\\
\textbf{名詞的英訳：} Appendix 1 is a description of the disassembly procedure of the unit.\\
\textbf{問題：} 前置詞 of が3回連続し、読みにくい\\
\textbf{動詞的英訳：} Appendix 1 describes how to disassemble the unit.
\end{itembox}

\begin{itembox}[l]{例：「行う」を削る}
\textbf{原文：} 「両企業はライセンス契約の修正を行った」\\
\textbf{名詞的英訳：} The two companies made an amendment to the license agreement.\\
\textbf{動詞的英訳：} The two companies amended the license agreement.
\end{itembox}

動詞を使うメリットは2つある。
\begin{enumerate}
\item \textbf{前置詞が減る：} 名詞句には前置詞（of, to, for等）が必要だが、動詞にすれば不要になることが多い。
\item \textbf{文が短くなる：} 「〜の実施を行った」（名詞+動詞）が「〜を実施した」（動詞のみ）になる。
\end{enumerate}

\begin{screen}
\textbf{4つの原則のまとめ：} 英訳の前に日本語を整える4つのステップ：（1）すっきり最小主義で冗長さを削り、（2）具体化主義で曖昧さをなくし、（3）動詞主義で名詞句を動詞に戻し、（4）能動態主義で受動態を能動態に変える。この4ステップを踏むだけで、AIによる英訳の品質は格段に上がる。
\end{screen}

%----------------------------------------------------------------------
\subsection{能動態と受動態の使い分け}
\label{subsec:voice}
\index{のうどうたい@能動態}
\index{じゅどうたい@受動態}

\subsubsection{「まっすぐ能動態主義」の原則}

日本語は受動態を多用する言語であるが、英語---特に技術英語や学術英語---では能動態が好まれる。「すっきり最小主義」（余計な情報を削ぎ落とす）と「まっすぐ能動態主義」（まず能動態で書くことを第一選択とする）を基本原則として覚えておくとよい。

日本語の技術文書では、「〜されている」「〜が行われた」といった受動態が自然であり、多用される。しかし、これをそのまま英語にすると冗長で読みにくい文章になる。英語では「誰が何をしたか」を明確にする能動態がストレートで好まれる。

\begin{itembox}[l]{例1: 日本語的な受動態から英語的な能動態へ}
\begin{itemize}
\item \textbf{日本語:} 「作業油の流量は自動で調整されている」
\item \textbf{受動態英訳:} ``The flow of the fluid is automatically regulated.''
\item \textbf{能動態英訳:} ``Pressure-sensitive valves automatically regulate the flow of the fluid.'' （主語が明確でストレート）
\end{itemize}
\end{itembox}

\begin{itembox}[l]{例2: 学術論文での能動態と受動態}
\begin{itemize}
\item \textbf{受動態:} ``The experiment was conducted by the authors.'' （冗長）
\item \textbf{能動態:} ``The authors conducted the experiment.'' （よりストレート）
\item \textbf{受動態:} ``It was observed that the accuracy was improved.'' （回りくどい）
\item \textbf{能動態:} ``We observed that the accuracy improved.'' （明確）
\end{itemize}
\end{itembox}

\subsubsection{受動態を使うべき4つの場面}

「まっすぐ能動態主義」とはいえ、受動態が適切な場面もある。以下の4つのケースでは、受動態を積極的に使うべきである。

\begin{enumerate}
\item \textbf{主語の統一（Subject Consistency）:} パラグラフ内で主語を統一するために受動態を使う。文ごとに主語がころころ変わると読みにくくなる。
\item \textbf{主語の短縮（Shorter Subject）:} 能動態にすると主語が非常に長くなる場合、受動態を使って主語を短くする。
\item \textbf{責任のあいまい化（Ambiguating Agency）:} 行為者を明示したくない場合、または行為者が重要でない場合に受動態を使う。
\item \textbf{話題の流れ（Topic Flow）:} 既知情報を文頭に、新情報を文末に配置するために受動態を使う。英語では「旧情報→新情報」の流れが自然である。
\end{enumerate}

\begin{itembox}[l]{受動態が適切な場面の例}
\begin{itemize}
\item 主語の統一: ``The samples were collected from three sites. They were then filtered and analyzed.'' （``The samples''を主語に統一）
\item 手法の説明: ``The samples were analyzed using mass spectrometry.'' （誰がやったかより何をしたかが重要）
\item 責任のあいまい化: ``Mistakes were made in the initial calibration.'' （誰の責任かを明示しない）
\end{itemize}
\end{itembox}

\begin{itembox}[l]{注意: 「うっかり受け身」は通じない}
日本語では「うっかり壊してしまった」のように、意図せずにそうなったことを受身で表現する（迷惑の受身）。しかし英語には同じ機能の受動態がないため、日本語の感覚でそのまま受動態を使うと意味が通じなくなる。英語では意図の有無にかかわらず、行為者を主語にして能動態で書くのが原則である。
\end{itembox}

\begin{promptbox}
以下の英文の受動態を能動態に書き換えてください。ただし、受動態の方が適切な場合は
その理由（主語の統一・主語の短縮・責任のあいまい化・話題の流れのいずれか）とともに
残してください。

{[}英文を貼り付け]
\end{promptbox}

以下に、受動態から能動態への代表的な改善パターンを早見表としてまとめる。

\begin{table}[htbp]
\centering
\caption{受動態から能動態への改善例}
\label{tab:passive}
\footnotesize
\begin{tabular}{ll}
\toprule
受動態（冗長） & 能動態（簡潔） \\
\midrule
The model was trained by us using... & We trained the model using... \\
It was observed that the accuracy & We observed that the accuracy \\
\quad was improved by... & \quad improved by... \\
The data was collected from... & We collected data from... \\
\bottomrule
\end{tabular}
\end{table}

\begin{itembox}[l]{Tip: ``We''の使用}
``We''を主語にすることに抵抗がある方もいるが、現代の学術英語では``We''の使用は一般的に受け入れられている。ただし、過度に``We''を繰り返さないよう、受動態と能動態を適切に混ぜることが理想的である。
\end{itembox}

%----------------------------------------------------------------------
\subsection{冠詞（a/the）の使い分け}
\label{subsec:articles}
\index{かんし@冠詞}

\subsubsection{なぜ冠詞が難しいのか}

冠詞（a, an, the, 無冠詞）は、日本語にない概念であるため、日本語話者にとって最も間違いやすい文法事項である。冠詞の選択は、「読者がその名詞をどの程度特定できるか」という情報の共有状態に基づいている。

\subsubsection{theの3つの用法}

定冠詞theには、大きく分けて3つの用法がある。

\begin{enumerate}
\item \textbf{唯一照応:} 世界に一つしかないもの、文脈上唯一に定まるもの。\\
例: ``the sun''（太陽は一つ）、``the Internet''（インターネットは一つ）、``the proposed method''（提案手法は一つ）
\item \textbf{前方照応:} 前の文で既に導入されたものを指す。a/anで初出した名詞を、2回目以降theで受ける。\\
例: ``Click the button to open \textbf{a} new window. Drag the images onto \textbf{the} new window.''（aで導入→theで再参照）
\item \textbf{後方照応:} 名詞の直後に修飾語句（関係節、前置詞句など）があり、それによって特定されるもの。\\
例: ``Locate \textbf{the} four screws that are behind the lid.''（``that are behind the lid''によって特定される）
\end{enumerate}

\begin{itembox}[l]{学術論文での冠詞の具体例}
\begin{itemize}
\item 初出: ``We propose \textbf{a} method for text classification.''
\item 再出: ``\textbf{The} method consists of three modules.''（前方照応）
\item 特定: ``\textbf{The} method proposed by Smith et al. (2023)''（後方照応）
\item 唯一: ``\textbf{The} proposed method outperforms all baselines.''（唯一照応）
\item よくある間違い: ``We propose method'' → ``We propose \textbf{a} method''（冠詞の欠落）
\end{itemize}
\end{itembox}

\subsubsection{可算名詞と不可算名詞}

同じ単語でも、可算名詞と不可算名詞で意味が変わることがある。冠詞の有無が意味を決定するため、注意が必要である。

\begin{itemize}
\item \textbf{iron}（不可算）= 鉄（物質） / \textbf{an iron}（可算）= アイロン（道具）
\item \textbf{chicken}（不可算）= 鶏肉（食材） / \textbf{a chicken}（可算）= ニワトリ（動物）
\item \textbf{paper}（不可算）= 紙（素材） / \textbf{a paper}（可算）= 論文
\item \textbf{glass}（不可算）= ガラス（素材） / \textbf{a glass}（可算）= コップ
\end{itemize}

\begin{promptbox}
以下の英文の冠詞（a/an/the）の使い方を確認し、誤りがあれば修正してください。
修正した場合は、唯一照応・前方照応・後方照応のどれに該当するか説明してください。

{[}英文を貼り付け]
\end{promptbox}

以下に、冠詞の基本ルールを早見表としてまとめる。

\textbf{基本ルール:}
\begin{itemize}
\item \textbf{the:} 特定のもの、前出のもの、唯一のもの（唯一照応・前方照応・後方照応）
\item \textbf{a/an:} 不特定のもの、初出のもの
\item \textbf{無冠詞:} 複数形の一般論、不可算名詞の一般論
\end{itemize}

\begin{promptbox}
以下の英語論文の冠詞の使用を確認してください。
特に以下のパターンに注意してチェックしてください:

1. 初出の名詞に "the" を使っていないか
2. 前出の名詞に "a/an" を使っていないか
3. 数式や手法名に不要な冠詞がついていないか
4. "the proposed method" のように特定される名詞に
   "the" が付いているか

{[}英語論文を貼り付け]
\end{promptbox}

%----------------------------------------------------------------------
\subsection{時制の選択}
\label{subsec:tense}
\index{じせい@時制}

\subsubsection{論文のセクションと時制の関係}

英語論文では、セクションによって適切な時制が決まっている。これは単なる慣習ではなく、各セクションで述べる内容の性質に基づいている。

\begin{itemize}
\item \textbf{Methods（手法）:} \textbf{過去形}を使う。実験は過去に行ったことだからである。\\
例: ``We trained the model using 10,000 samples.''
\item \textbf{Results（結果）:} \textbf{過去形}を使う。結果も過去に得たものだからである。\\
例: ``The proposed method achieved an accuracy of 95.3\%.''
\item \textbf{Introduction / Discussion:} \textbf{現在形}を使う。一般的な事実や先行研究の知見は現在も成り立つものだからである。\\
例: ``Deep learning has become a dominant approach in NLP.''
\end{itemize}

\subsubsection{単純現在形の本質: 再現性100\%}

単純現在形は、「いつでも・誰がやっても成り立つ事実」を述べるときに使う。つまり、再現性が100\%であることを主張する時制である。

\begin{itemize}
\item ``Water boils at 100°C.''（水は100°Cで沸騰する = いつでも成り立つ事実）
\item ``This algorithm converges in $O(n \log n)$ time.''（このアルゴリズムは$O(n \log n)$で収束する = 数学的事実）
\end{itemize}

逆に言えば、単純現在形を使うということは「これは普遍的な事実である」と主張していることになる。自分の実験結果に単純現在形を使うと、「この結果はいつでも再現される」という強い主張になるので注意が必要である。

\subsubsection{現在完了形の3つの用法}

現在完了形は学術論文で頻繁に使われるが、以下の3つの用法を区別して使う必要がある。

\begin{enumerate}
\item \textbf{完了（already）:} 動作が完了したことを表す。\\
例: ``We have completed the analysis.''（分析を完了した）
\item \textbf{経験（times）:} 過去の経験を表す。\\
例: ``Several studies have shown that this approach is effective.''（複数の研究が示してきた）
\item \textbf{継続（for / since）:} 過去から現在まで続いていることを表す。\\
例: ``This method has been used for over a decade.''（10年以上使われ続けている）
\end{enumerate}

\subsubsection{助動詞の使い分け}

助動詞は、話者の「確信度」を表す重要な手段である。学術論文では、主張の強さを適切にコントロールするために助動詞を使い分ける。

\begin{itemize}
\item \textbf{will} = 確信・コミットメント（確信度が高い）\\
例: ``This approach will enable faster training.''
\item \textbf{can} = 可能性・能力\\
例: ``This method can handle large-scale datasets.''
\item \textbf{may} = 不確実性・推測（確信度が低い）\\
例: ``This discrepancy may be due to overfitting.''
\end{itemize}

\begin{itembox}[l]{注意: ``cannot'' と ``can not'' の違い}
``cannot''（1語）は「〜できない」という不可能の意味である。一方、``can not''（2語）は``can not only A but also B''のように、``not''が後続の語句を否定する場合に使う。通常の「〜できない」を表すときは必ず``cannot''と1語で書く。
\end{itembox}

以下に、セクション別の適切な時制を早見表としてまとめる。

\begin{table}[htbp]
\centering
\caption{セクション別の適切な時制}
\label{tab:tense}
\footnotesize
\begin{tabular}{lll}
\toprule
セクション & 基本の時制 & 例 \\
\midrule
Introduction & 現在形（一般的事実）/ & ``Deep learning has achieved...'' \\
 & 過去形（先行研究の報告） & ``Smith et al. (2023) proposed...'' \\
Methods & 過去形 & ``We trained the model using...'' \\
Results & 過去形 & ``The proposed method achieved...'' \\
Discussion & 現在形（解釈）/ & ``These results suggest...'' \\
 & 過去形（結果への言及） & ``Our model showed...'' \\
\bottomrule
\end{tabular}
\end{table}

\begin{promptbox}
以下の英語論文のMethodsセクションについて、
時制が適切に使われているか確認してください。
Methodsでは基本的に過去形を使用します。
不適切な時制が使われている箇所を指摘し、修正してください。

{[}Methodsセクションを貼り付け]
\end{promptbox}

%----------------------------------------------------------------------
\subsection{自動詞と他動詞の注意}
\label{subsec:transitive}
\index{じどうし@自動詞}
\index{たどうし@他動詞}

英語では、同じ動詞でも自動詞として使うか他動詞として使うかで意味が大きく変わることがある。日本語では助詞（「を」「に」など）で関係を示すため、動詞の自他をあまり意識しないが、英語では目的語の有無が意味を決定する。

\begin{itembox}[l]{自動詞と他動詞で意味が変わる例}
\begin{itemize}
\item \textbf{submit}
  \begin{itemize}
  \item 他動詞: ``We submitted the paper.''（論文を提出した）
  \item 自動詞: ``The army submitted.''（軍が降伏した・屈服した）
  \end{itemize}
\item \textbf{search}
  \begin{itemize}
  \item 他動詞: ``I searched him.''（彼をボディチェックした = 身体検査した）
  \item 自動詞+前置詞: ``I searched for him.''（彼を探した）
  \end{itemize}
\item \textbf{run}
  \begin{itemize}
  \item 自動詞: ``The program runs.''（プログラムが動作する）
  \item 他動詞: ``We run the program.''（プログラムを実行する）
  \end{itemize}
\end{itemize}
\end{itembox}

特に``search''は要注意である。``search the database''は「データベースの中を探す」（正しい用法）だが、``search him''は「彼の身体を捜索する」という意味になり、「彼を探す」とは全く異なる。「〜を探す」と言いたい場合は``search for''を使う。

学術論文で動詞を使う際は、辞書で自動詞・他動詞の区別を必ず確認すること。特に、日本語で「〜を」と言える動詞が英語でも他動詞とは限らない点に注意が必要である。

%----------------------------------------------------------------------
\subsection{直訳による不自然な表現}
\index{ちょくやく@直訳}

日本語を直訳すると、文法的には正しくても不自然な英語になることがある。

\begin{table}[htbp]
\centering
\caption{直訳による不自然な表現と自然な英語}
\label{tab:unnatural}
\footnotesize
\begin{tabular}{ll}
\toprule
日本語直訳（不自然） & 自然な英語 \\
\midrule
``We could confirm that...'' & ``We confirmed that...'' \\
``It is considered that...'' & ``We consider...'' / ``This suggests...'' \\
``From this result...'' & ``Based on these results...'' \\
``We performed the experiment...'' & ``We conducted experiments...'' \\
``The accuracy became high...'' & ``The accuracy improved...'' \\
``As a merit of our method...'' & ``An advantage of our method is...'' \\
\bottomrule
\end{tabular}
\end{table}

\begin{promptbox}
以下の英語論文について、日本語の直訳に起因する不自然な表現を
特定してください。各箇所について、より自然な英語表現を提案し、
なぜその表現の方が自然かを説明してください。

{[}英語論文を貼り付け]
\end{promptbox}

%----------------------------------------------------------------------
\subsection{専門用語の確認方法}
\index{せんもんようご@専門用語}

専門用語の英語表現が分からない場合や、正しいかどうか確認したい場合は、以下の方法が有効である。

\begin{enumerate}
\item \textbf{Google Scholarでの用例検索:} ``用語''を含む論文を検索し、トップ論文でどのように使われているか確認する
\item \textbf{AIに確認を依頼:}
\end{enumerate}

\begin{promptbox}
「[日本語の専門用語]」の英語表現として、以下のどちらが
{[}分野名]の学術論文で一般的ですか？

A: [候補1]
B: [候補2]

それぞれの使用頻度や文脈の違いがあれば教えてください。
また、他に一般的な表現があれば提示してください。
\end{promptbox}

%======================================================================
\section{学術英語表現集}
\label{sec:4.7}
\index{がくじゅつえいごひょうげん@学術英語表現}

\subsection{Introductionの定型表現}

\begin{table}[htbp]
\centering
\caption{Introductionで使える定型表現}
\label{tab:intro-expressions}
\footnotesize
\begin{tabular}{lp{0.55\textwidth}}
\toprule
目的 & 表現例 \\
\midrule
分野の重要性を述べる & ``X has attracted significant attention in recent years.'' \\
 & ``X plays a crucial role in...'' \\
 & ``X is of fundamental importance to...'' \\
\midrule
先行研究を紹介する & ``Several studies have investigated...'' \\
 & ``Previous work has shown that...'' \\
 & ``Smith et al. (2023) demonstrated that...'' \\
\midrule
問題点を指摘する & ``However, existing methods suffer from...'' \\
 & ``Despite these advances, ... remains challenging.'' \\
 & ``A major limitation of prior work is...'' \\
\midrule
本研究の目的を述べる & ``In this paper, we propose...'' \\
 & ``This work addresses the problem of...'' \\
 & ``We present a novel approach to...'' \\
\midrule
貢献を列挙する & ``Our main contributions are as follows:'' \\
 & ``The key contributions of this work include:'' \\
\bottomrule
\end{tabular}
\end{table}

\subsection{Methodsの定型表現}

\begin{table}[htbp]
\centering
\caption{Methodsで使える定型表現}
\label{tab:methods-expressions}
\footnotesize
\begin{tabular}{lp{0.55\textwidth}}
\toprule
目的 & 表現例 \\
\midrule
手法の全体像を述べる & ``Our approach consists of three main components:'' \\
 & ``The overall architecture is illustrated in Figure X.'' \\
 & ``We formulate the problem as...'' \\
\midrule
手順を説明する & ``First, we... Then, we... Finally, we...'' \\
 & ``The input is processed through...'' \\
 & ``We apply X to obtain Y.'' \\
\midrule
数式を導入する & ``The objective function is defined as:'' \\
 & ``We optimize the following loss function:'' \\
 & ``where X denotes the...'' \\
\midrule
実装の詳細を述べる & ``We implemented our method using...'' \\
 & ``All experiments were conducted on...'' \\
 & ``The hyperparameters were tuned via...'' \\
\bottomrule
\end{tabular}
\end{table}

\subsection{Resultsの定型表現}

\begin{table}[htbp]
\centering
\caption{Resultsで使える定型表現}
\label{tab:results-expressions}
\footnotesize
\begin{tabular}{lp{0.55\textwidth}}
\toprule
目的 & 表現例 \\
\midrule
結果を提示する & ``Table X summarizes the results of...'' \\
 & ``As shown in Figure X, our method...'' \\
 & ``The results are presented in Table X.'' \\
\midrule
比較する & ``Our method outperforms the baseline by X\%.'' \\
 & ``Compared to the state-of-the-art, our approach achieves...'' \\
 & ``The improvement is statistically significant ($p < 0.05$).'' \\
\midrule
傾向を述べる & ``We observe a consistent improvement across...'' \\
 & ``The performance increases as X grows.'' \\
 & ``Notably, our method shows robust performance on...'' \\
\bottomrule
\end{tabular}
\end{table}

\subsection{Discussion / Conclusionの定型表現}

\begin{table}[htbp]
\centering
\caption{Discussion / Conclusionで使える定型表現}
\label{tab:discussion-expressions}
\footnotesize
\begin{tabular}{lp{0.55\textwidth}}
\toprule
目的 & 表現例 \\
\midrule
結果を解釈する & ``These results suggest that...'' \\
 & ``One possible explanation is that...'' \\
 & ``This finding is consistent with...'' \\
\midrule
限界を述べる & ``A limitation of our study is...'' \\
 & ``Our approach may not generalize to...'' \\
 & ``One potential concern is...'' \\
\midrule
将来の方向を述べる & ``Future work could explore...'' \\
 & ``A promising direction is to...'' \\
 & ``We plan to extend this work by...'' \\
\midrule
結論を述べる & ``In summary, we have presented...'' \\
 & ``This work demonstrates the effectiveness of...'' \\
 & ``Our findings highlight the potential of...'' \\
\bottomrule
\end{tabular}
\end{table}

\begin{itembox}[l]{Tip: 定型表現の活用法}
これらの表現を暗記する必要はない。執筆時に参照表として使い、自分の分野の論文でよく見かける表現を徐々に身につけていくとよい。定型表現は「型」として活用し、その中に自分の研究の独自の内容を盛り込むことが重要である。
\end{itembox}

%======================================================================
\section*{演習問題}

\begin{enumerate}

\item \textbf{演習1: IMRaD構造の分析}

自分の研究分野のトップ会議またはジャーナルから論文を1本選び、以下を行いなさい。
\begin{enumerate}
\item 各セクション（Introduction, Methods, Results, Discussion）の段落数と概算語数を調べる
\item Introductionの「逆三角形構造」（背景→問題→本研究の貢献）を具体的に特定する
\item 各段落のトピックセンテンスを抜き出し、それだけを読んで論文の流れが理解できるか確認する
\item この分析から得られた、英語論文の構造に関する気づきをまとめる
\end{enumerate}

\item \textbf{演習2: AIを使った英訳の実践}

自分の研究に関する日本語の文章（300--500字）を用意し、以下を行いなさい。
\begin{enumerate}
\item 本章のStep 1のチェックリストを使って日本語原稿を確認・修正する
\item Step 2のプロンプトテンプレートを使って、AIに英訳を依頼する
\item Step 3のプロンプトを使って、英訳を学術的に洗練する
\item Step 2の英訳とStep 3の洗練後の英訳を比較し、どのような改善がなされたかを分析する
\item 自分の分野の論文と表現を比較し、さらなる改善点がないか検討する
\end{enumerate}

\item \textbf{演習3: 逆翻訳チェックの実践}

演習2で作成した英語論文の一部について、以下を行いなさい。
\begin{enumerate}
\item 別のAIツール（演習2でClaudeを使った場合はChatGPT等）で日本語に逆翻訳する
\item 逆翻訳された日本語を元の日本語原稿と比較し、意味のズレを特定する
\item ズレが見つかった箇所について、英訳のどこに問題があるかを分析する
\item 英訳を修正し、再度逆翻訳して意味のズレが解消されたか確認する
\end{enumerate}

\item \textbf{演習4: Abstract作成の実践}

自分の研究（または研究計画）について、以下を行いなさい。
\begin{enumerate}
\item 本章のAbstract構造（背景→目的→手法→結果→意義）に従い、各要素を日本語で箇条書きにする
\item 本章のプロンプトテンプレートを使って、AIに英語のAbstractを作成させる
\item 作成されたAbstractが150--200語に収まっているか確認する
\item 本章の「良いAbstractの条件」（具体性、数値、構造）を満たしているか評価する
\item 必要に応じて修正し、最終版を完成させる
\end{enumerate}

\item \textbf{演習5: 総合推敲の実践}

演習2--4で作成した英語テキストをまとめ、以下の段階的推敲を行いなさい。
\begin{enumerate}
\item 第1段階として、文法チェックのプロンプトを使って文法エラーを修正する
\item 第2段階として、表現チェックのプロンプトを使って学術的表現に洗練する
\item 第3段階として、論理チェックのプロンプトを使って論理構造を確認する
\item 最後に、ネイティブチェック風レビューのプロンプトを使って自然さを確認する
\item 各段階で見つかった問題の種類と数を記録し、自分の英語ライティングの傾向を分析する
\end{enumerate}

\end{enumerate}

\vspace{10mm}
\begin{itembox}[l]{本章のまとめ}
英語論文の執筆は、日本語話者にとって大きな挑戦であるが、AIを活用することでそのハードルを大幅に下げることができる。重要なのは、AIを「翻訳機」としてではなく「執筆パートナー」として活用することである。AIの出力をそのまま使うのではなく、段階的な推敲プロセスを通じて品質を高め、最終的には自分の言葉として責任を持てる英語論文に仕上げることが大切である。繰り返し実践することで、AIに頼らなくても書ける力が自然と身についていく。
\end{itembox}

\putbib
\end{bibunit}

\begin{bibunit}
% -*- coding: utf-8 -*-
\chapter{研究発表スライドの作成}
\label{chap5}

\begin{chapterbox}
\textbf{【この章で学ぶこと】}
\begin{itemize}
\item \textbf{AIを活用して論理的で説得力のある研究発表スライドを作成できる} --- 発表の構成案からスライド内容まで、AIを効果的に活用した作成プロセスを実践できる
\item \textbf{デザインAIツールを使った効率的なスライド作成法を習得する} --- GammaやCanva AIなどのツールを活用し、視覚的に優れたスライドを短時間で作成できる
\item \textbf{図表・概念図の生成と批判的な検討ができる} --- AIによる図表生成の可能性と限界を理解し、適切に活用・検証できる
\end{itemize}
\end{chapterbox}
\vspace{5mm}

%======================================================================
\section{研究発表スライドの基本}
\label{sec:5.1}
\index{けんきゅうはっぴょうすらいど@研究発表スライド}

\subsection{学会発表の基本構成}
\index{がっかいはっぴょう@学会発表}

研究発表のスライドは、論文と同様に論理的な構成が求められます。ただし、スライドは「読む」ものではなく「見せる」ものであるため、論文とは異なるアプローチが必要です。

\textbf{標準的な発表構成:}

\begin{table}[htbp]
\centering
\caption{研究発表スライドの標準的な構成}
\label{tab:slide-structure}
\begin{tabular}{lll}
\toprule
スライド & 内容 & 目的 \\
\midrule
タイトル & 研究タイトル、著者名、所属 & 研究の第一印象を決める \\
背景 & 研究分野の現状と課題 & 聴衆に問題の重要性を理解させる \\
目的 & 本研究で解決する問題 & 研究のゴールを明確にする \\
関連研究 & 先行研究と本研究の位置づけ & 新規性の根拠を示す \\
提案手法 & アプローチの概要と詳細 & 技術的な貢献を説明する \\
実験設定 & データセット、評価指標、比較手法 & 実験の妥当性を示す \\
結果 & 主要な実験結果 & 提案手法の有効性を示す \\
考察 & 結果の解釈、分析 & 結果の意味を深く理解させる \\
まとめ & 貢献の要約と将来展望 & 研究のインパクトを印象づける \\
\bottomrule
\end{tabular}
\end{table}

\subsection{時間配分の目安（10分発表の場合）}

10分間の口頭発表は、学会で最も一般的な発表形式の一つです。以下に推奨される時間配分を示します。

\begin{table}[htbp]
\centering
\caption{10分発表の推奨時間配分}
\label{tab:time-allocation}
\begin{tabular}{lll}
\toprule
セクション & 時間 & スライド枚数（目安） \\
\midrule
タイトル・自己紹介 & 0.5分 & 1枚 \\
背景・問題設定 & 2分 & 2--3枚 \\
提案手法 & 3分 & 3--4枚 \\
実験結果 & 2.5分 & 2--3枚 \\
考察・まとめ & 2分 & 2枚 \\
\midrule
\textbf{合計} & \textbf{10分} & \textbf{10--13枚} \\
\bottomrule
\end{tabular}
\end{table}

1枚のスライドに費やす時間は、平均して45秒から1分程度です。情報を詰め込みすぎると、聴衆は内容を理解する前に次のスライドに移ってしまいます。

\subsection{良いスライドの原則}

\begin{itembox}[l]{1スライド1メッセージの原則}
各スライドが伝えるべきメッセージは1つだけです。複数のメッセージを1枚に詰め込むと、聴衆はどこに注目すべきかわからなくなります。スライドのタイトルに、そのスライドのメッセージを端的に表現しましょう。

\textbf{悪い例のタイトル:} 「実験結果」

\textbf{良い例のタイトル:} 「提案手法は全データセットでベースラインを上回った」
\end{itembox}

\textbf{文字量の制限:}

\begin{itemize}
\item タイトル: 1行（最大15語程度）
\item 本文: 1スライドあたり最大6行
\item フォントサイズ: タイトル28pt以上、本文20pt以上
\item 箇条書きは3--5項目まで
\end{itemize}

\textbf{視覚的な原則:}

\begin{itemize}
\item 図表を積極的に使う（テキストより図の方が記憶に残る）
\item 色の使いすぎを避ける（メインカラー + アクセントカラーの2--3色）
\item 余白を十分にとる（スライドの30\%以上は空白にする）
\item アニメーションは必要最小限にする
\end{itemize}

\begin{itembox}[l]{Tip: スライド作成の出発点}
スライドを作成する前に、まず「各スライドで伝えたいメッセージ」を箇条書きで書き出しましょう。メッセージの一覧が論理的に流れていれば、スライド作成はスムーズに進みます。この段階でAIを活用すると効率的です。
\end{itembox}

%======================================================================
\section{AIによる構成案の生成}
\label{sec:5.2}
\index{こうせいあん@構成案}

\subsection{研究テーマの構造化}
\index{えれべーたーぴっち@エレベーターピッチ}

AIに構成案を依頼する前に、自分の研究を1--2文で簡潔に表現できるようにしましょう。これは「エレベーターピッチ」とも呼ばれ、短時間で研究の本質を伝える技術です。

\begin{promptbox}
私の研究を以下の情報をもとに、1-2文の簡潔な説明にまとめてください。

- 研究分野: [分野]
- 解決する問題: [問題]
- 提案するアプローチ: [手法の概要]
- 主な結果: [成果]
- 研究の意義: [意義]

対象聴衆: [例: 同じ分野の研究者 / 異分野の研究者 / 一般的な理系の聴衆]
\end{promptbox}

\subsection{AIに構成案を提案させるプロンプト設計}

研究テーマが整理できたら、AIに発表スライドの構成案を提案させます。

\begin{promptbox}
以下の研究について、10分間の学会口頭発表のスライド構成案を
作成してください。

【研究概要】
タイトル: [研究タイトル]
分野: [研究分野]
概要: [1-2文の研究概要]

【詳細情報】
背景: [研究の背景と動機]
提案手法: [手法の概要]
主要結果: [主な実験結果]
結論: [研究の結論]

【発表条件】
発表時間: 10分（質疑5分）
対象聴衆: [聴衆の専門性レベル]
発表言語: [日本語/英語]

以下の形式で構成案を作成してください:
- スライド番号
- スライドタイトル（メッセージ型）
- 含めるべき内容（箇条書き3-5項目）
- 想定所要時間
- 視覚要素の提案（図表、グラフなど）
\end{promptbox}

\subsection{生成された構成案の批判的検討}

AIが生成した構成案は、必ず以下の4つの観点から批判的に検討します。

\textbf{1. 論理一貫性の確認:}

スライドのタイトルだけを順番に読んで、研究のストーリーが自然に流れているか確認します。

\begin{promptbox}
以下のスライド構成案について、論理的な一貫性を確認してください。
スライドタイトルの順序を読むだけで、研究のストーリーが
伝わるかどうかを評価してください。

論理の飛躍、不足している要素、順序の改善案があれば指摘してください。

[スライドタイトルのリスト]
\end{promptbox}

\textbf{2. 網羅性の確認:}

必要な情報がすべて含まれているか確認します。特に以下の要素は必須です。

\begin{itemize}
\item 研究の動機（なぜこの研究が重要か）
\item 先行研究との差別化（何が新しいのか）
\item 手法の核心的なアイデア（どのように解決するか）
\item 定量的な結果（どの程度うまくいったか）
\item 限界と将来展望（残された課題は何か）
\end{itemize}

\textbf{3. 独自性の確認:}

構成案が一般的すぎないか、自分の研究の独自性が反映されているか確認します。AIは汎用的な構成を提案する傾向があるため、自分の研究の強みが際立つように修正を加えます。

\textbf{4. 時間配分の確認:}

各スライドの想定所要時間を合計し、発表時間内に収まるか確認します。実際に声に出してリハーサルすると、想定よりも時間がかかることが多いため、余裕を持った時間配分にしましょう。

\begin{itembox}[l]{Tip: 早い段階でのフィードバック}
構成案の段階で指導教員やゼミのメンバーにフィードバックをもらうのが効果的です。スライドを作り込んでからの大幅な修正は時間のロスが大きいため、早い段階で方向性を固めましょう。
\end{itembox}

%======================================================================
\section{各スライドの内容生成}
\label{sec:5.3}

\subsection{研究背景の文章生成}

研究背景のスライドは、聴衆の注意を引き、研究の重要性を理解してもらうための出発点です。

\begin{promptbox}
以下の研究テーマの背景を、学会発表スライド用の簡潔な文章として
生成してください。

【研究テーマ】: [テーマ]
【解決したい問題】: [問題]
【対象聴衆】: [聴衆の属性]

以下の条件を守ってください:
1. 箇条書き形式で、1項目あたり1行（15-20文字程度）
2. 専門用語は[聴衆の属性]が理解できるレベルにする
3. 問題の重要性が直感的に伝わるデータや事実を含める
4. 最後に「この問題をどう解決するか？」への橋渡しとなる文で締める
5. 合計5-6行以内
\end{promptbox}

\subsection{研究手法の説明文生成}

提案手法の説明は、発表の核心部分です。複雑な手法を聴衆に理解してもらうためには、段階的な説明が重要です。

\begin{promptbox}
以下の研究手法を、学会発表スライド用に段階的に説明する文章を
生成してください。

【手法の概要】: [手法の詳細な説明]

以下の構成で、3枚のスライドに分けてください:

スライド1: 手法の全体像（Overview）
- 手法の核心的なアイデアを1文で
- 全体のパイプラインを箇条書き3-4項目で

スライド2: 技術的な詳細（Key Technique）
- 最も重要な技術的貢献に焦点
- 数式は最小限（1-2個まで）
- 直感的な説明を優先

スライド3: 実装のポイント
- 実用的な工夫
- 計算コストや効率性

各スライドは5行以内で、図表の提案も含めてください。
\end{promptbox}

\subsection{研究結果のまとめ生成}

\begin{promptbox}
以下の実験結果を、学会発表スライド用にまとめてください。

【実験結果】:
[結果のデータを貼り付け]

以下の形式でお願いします:
1. 主要な発見を重要度順に3つ挙げる
2. 各発見を1行で端的に表現する
3. 数値は効果的に強調する（例: "精度12.5\%向上"）
4. 図表（グラフの種類）の提案
5. 聴衆が「なるほど」と思えるような解釈を1文添える
\end{promptbox}

\subsection{専門用語と平易な表現のバランス調整}

学会発表では、聴衆の専門性に応じて表現のレベルを調整する必要があります。

\begin{promptbox}
以下のスライド内容について、対象聴衆に合わせて表現を調整してください。

【現在の内容】:
[スライドの内容]

【対象聴衆】: [例: 情報科学を専攻する修士学生]

調整の方針:
1. 聴衆が確実に知っている用語はそのまま使用
2. 聴衆が知らない可能性のある用語は、初出時に簡単な説明を付加
3. 略語は初出でフルスペルを示す
4. 過度に専門的な詳細は省略し、直感的な説明に置き換える
\end{promptbox}

\subsection{AIの出力を自分の言葉に修正する重要性}

AIが生成したスライド内容は、あくまで出発点です。以下の理由から、必ず自分の言葉に修正してください。

\begin{itembox}[l]{なぜ自分の言葉に修正が必要か}
\begin{enumerate}
\item \textbf{発表の自然さ:} AIが生成した文章をそのまま読み上げると、不自然な発表になります。自分が普段使う言葉遣いに修正することで、自信を持って発表できます。
\item \textbf{質疑への対応:} 自分の言葉でスライドを作成していれば、質疑応答で深い質問をされても自然に回答できます。AIの文章をそのまま使っていると、内容を十分に理解していないことが露呈するリスクがあります。
\item \textbf{独自の視点の反映:} AIは汎用的な表現を生成しますが、あなた自身の研究に対する独自の洞察や情熱は、自分の言葉でしか表現できません。
\end{enumerate}
\end{itembox}

\begin{itembox}[l]{Tip: 違和感を大切にする}
AIが生成した内容を読んで「これは自分の考えと違う」と感じた箇所は、積極的に修正しましょう。その違和感こそが、あなた自身の研究に対する理解の深さを反映しています。
\end{itembox}

%======================================================================
\section{デザインAIによるスライド作成}
\label{sec:5.4}
\index{でざいんAI@デザインAI}

\subsection{Gamma --- テキストからスライドを自動生成}
\index{Gamma@Gamma}

Gamma（gamma.app）は、テキスト入力から高品質なプレゼンテーションスライドを自動生成するAIツールです。研究発表スライドの初稿を短時間で作成するのに適しています。

\textbf{基本的な使い方:}

\begin{enumerate}
\item Gammaのウェブサイトにアクセスし、アカウントを作成する
\item 「Create new」から「Presentation」を選択する
\item 研究の概要テキストを入力するか、アウトライン形式で構成を指定する
\item AIがスライドのデザインと内容を自動生成する
\item 生成されたスライドを編集・カスタマイズする
\end{enumerate}

\textbf{研究発表に使う際のプロンプト例:}

\begin{promptbox}
Create a 12-slide academic research presentation:

Title: [研究タイトル]

Slide 1: Title slide
Slide 2: Research Background - [背景の要点]
Slide 3: Problem Statement - [問題の記述]
Slide 4: Related Work - [先行研究の位置づけ]
Slide 5-7: Proposed Method - [手法の説明]
Slide 8: Experimental Setup - [実験設定]
Slide 9-10: Results - [結果の要点]
Slide 11: Discussion - [考察の要点]
Slide 12: Conclusion and Future Work - [結論]

Style: Clean, minimal, academic
Color scheme: Blue and white
\end{promptbox}

\textbf{Gammaの長所:}
\begin{itemize}
\item テキストから短時間でスライドを生成できる
\item デザインが洗練されている
\item 図表のプレースホルダーが自動配置される
\item エクスポート形式が豊富（PDF, PPTX等）
\end{itemize}

\textbf{Gammaの短所:}
\begin{itemize}
\item 細かいレイアウトの調整には限界がある
\item 数式の表示には対応していない場合がある
\item 学術発表特有のフォーマット（ヘッダー/フッターなど）のカスタマイズが限定的
\item 無料プランでは一部機能が制限される
\end{itemize}

\subsection{Canva AI --- テンプレートベースのデザイン支援}
\index{Canva AI@Canva AI}

Canva（canva.com）は、テンプレートベースのデザインツールで、AI機能（Magic Design）によるスライド作成支援を提供しています。

\textbf{基本的な使い方:}

\begin{enumerate}
\item Canvaにアクセスし、「プレゼンテーション」を選択する
\item テンプレートを選択する（学術発表向けテンプレートが複数用意されている）
\item Magic Designにテーマを入力すると、テンプレートに合わせたスライドが提案される
\item 各スライドの内容を編集する
\item 図表やアイコンをCanvaの素材ライブラリから追加する
\end{enumerate}

\textbf{学術発表に適したテンプレートの選び方:}

\begin{itemize}
\item シンプルで余白が十分にあるもの
\item 色数が少ないもの（2--3色）
\item フォントが読みやすいもの（サンセリフ体推奨）
\item 装飾的な要素が少ないもの
\item グラフや表のプレースホルダーがあるもの
\end{itemize}

\begin{itembox}[l]{Tip: 最終調整はPowerPointで}
デザインAIで作成したスライドは、必ずPowerPointやKeynote形式でエクスポートし、最終的な微調整はこれらのツールで行いましょう。特に数式、特殊な図表、アニメーションについては、PowerPoint等で仕上げた方が柔軟性が高いです。
\end{itembox}

\subsection{デザインAIの使い分け}

\begin{table}[htbp]
\centering
\caption{デザインAIツールの比較}
\label{tab:design-ai-comparison}
\begin{tabular}{lll}
\toprule
項目 & Gamma & Canva AI \\
\midrule
強み & テキストからの自動生成 & テンプレートの豊富さ \\
作成速度 & 非常に速い & 速い \\
カスタマイズ性 & 中程度 & 高い \\
数式対応 & 限定的 & 限定的 \\
共同編集 & 対応 & 対応 \\
学術向けテンプレート & 少なめ & 多い \\
推奨用途 & 初稿の高速作成 & デザインの洗練 \\
\bottomrule
\end{tabular}
\end{table}

%======================================================================
\section{図表・概念図の生成}
\label{sec:5.5}
\index{がいねんず@概念図}

\subsection{画像生成AIによる概念図}
\index{がぞうせいせいAI@画像生成AI}

研究発表では、手法の概念図やシステムの全体像を示す図が非常に重要です。画像生成AI（DALL-E、Stable Diffusion等）を活用して、概念図のアイデアを得ることができます。

\textbf{活用できる場面:}

\begin{itemize}
\item 研究のイメージを伝えるための概要図
\item 応用分野のイラスト（医療、環境、ロボティクスなど）
\item スライドの背景画像やアイコン
\end{itemize}

\textbf{活用が難しい場面:}

\begin{itemize}
\item 正確な技術的アーキテクチャ図（ネットワーク構造など）
\item 数値データを含むグラフや表
\item 数式や記号を含む図
\end{itemize}

\textbf{プロンプト例:}

\begin{promptbox}
Create a clean, minimalist scientific illustration showing the concept
of transfer learning in natural language processing.
Show knowledge flowing from a large pre-trained model to a smaller
task-specific model. Use a blue and white color scheme.
Style: technical diagram, flat design, no text.
\end{promptbox}

\begin{itembox}[l]{Tip: 2段階アプローチ}
画像生成AIは概念的なイラストの「アイデア出し」には有用ですが、最終的な研究発表用の図は、専用のツール（PowerPoint、Illustrator、draw.io等）で自作することを推奨します。生成画像には不正確な要素が含まれている可能性があり、学術発表での使用には注意が必要です。
\end{itembox}

\subsection{テキストベース図表ツール --- Mermaid.jsの活用}
\index{Mermaid.js@Mermaid.js}

Mermaid.jsは、テキストベースで図表を生成するツールです。AIにMermaidコードを生成させることで、フローチャート、シーケンス図、状態遷移図などを効率的に作成できます。

\textbf{フローチャートの生成プロンプト:}

\begin{promptbox}
以下の研究手法のパイプラインを、Mermaid.jsのフローチャートとして
表現してください。

【手法のパイプライン】:
1. 入力データの前処理（テキストのトークナイズ、正規化）
2. 特徴抽出（BERTによるエンコーディング）
3. 分類層（全結合層 + Softmax）
4. 出力（カテゴリラベル）

各ステップ間のデータの流れが明確にわかるようにしてください。
重要なステップは色を変えて強調してください。
\end{promptbox}

\textbf{AIが生成するMermaidコードの例:}

\begin{lstlisting}[language={}]
graph TD
    A[入力テキスト] --> B[前処理]
    B --> B1[トークナイズ]
    B --> B2[正規化]
    B1 --> C[特徴抽出]
    B2 --> C
    C --> C1[BERT エンコーダ]
    C1 --> D[分類層]
    D --> D1[全結合層]
    D1 --> D2[Softmax]
    D2 --> E[カテゴリラベル出力]

    style C1 fill:#4A90D9,stroke:#333,color:#fff
    style D2 fill:#4A90D9,stroke:#333,color:#fff
\end{lstlisting}

\textbf{シーケンス図の生成プロンプト:}

\begin{promptbox}
以下のシステムの処理フローを、Mermaid.jsのシーケンス図として
表現してください。

【処理フロー】:
1. ユーザーがクエリを入力
2. フロントエンドがAPIサーバーにリクエストを送信
3. APIサーバーがデータベースに問い合わせ
4. データベースが結果を返す
5. APIサーバーがAIモデルに推論を依頼
6. AIモデルが結果を返す
7. APIサーバーがフロントエンドにレスポンスを返す
8. フロントエンドがユーザーに結果を表示
\end{promptbox}

\textbf{研究手法の比較図:}

\begin{promptbox}
以下の3つの手法の比較を、Mermaid.jsの図として表現してください。
各手法の入力、処理ステップ、出力を並列に表示し、
共通点と相違点が一目でわかるようにしてください。

手法A: [概要]
手法B: [概要]
手法C（提案手法）: [概要]
\end{promptbox}

\textbf{Mermaidコードの活用方法:}

生成されたMermaidコードは、以下の方法でスライドに組み込むことができます。

\begin{enumerate}
\item \textbf{Mermaid Live Editor}（mermaid.live）でSVGまたはPNG画像としてエクスポート
\item \textbf{VS Code拡張機能}（Markdown Preview Mermaid Support）でプレビュー
\item \textbf{GitHubのMarkdown}はMermaidを直接レンダリング可能
\item エクスポートした画像をPowerPointやKeynoteに貼り付け
\end{enumerate}

\subsection{生成画像の著作権と正確性の検証}
\index{ちょさくけん@著作権}

AIで生成した図表をスライドに使用する際は、以下の点に注意が必要です。

\textbf{著作権に関する注意:}

\begin{itemize}
\item 画像生成AIの利用規約を確認する（商用利用、学術利用の可否）
\item DALL-Eで生成した画像はOpenAIの利用規約に基づいて使用可能
\item Stable Diffusionはモデルのライセンスによって異なる
\item 学会の規定によっては、AI生成画像の使用に制限がある場合がある
\end{itemize}

\textbf{正確性の検証:}

\begin{itemize}
\item AIが生成した概念図の要素が正確かどうかを確認する
\item 生成画像に含まれるテキストは不正確なことが多いため、自分で修正する
\item 技術的な図（ネットワーク構造、数式を含む図）は特に注意して検証する
\item 不正確な図を使用すると、査読者や聴衆の信頼を損なう可能性がある
\end{itemize}

\begin{itembox}[l]{Tip: Overview図の作成}
学会発表では、手法の全体像を示す「Overview図」が特に重要です。AIで概念図のラフ案を生成し、それを参考にしてPowerPointやdraw.ioで正確な図を自作するという2段階のアプローチが効果的です。最終的な図には、自分で正確性を保証できるものだけを使用しましょう。
\end{itembox}

%======================================================================
\section{相互レビューと改善}
\label{sec:5.6}
\index{そうごれびゅー@相互レビュー}

\subsection{ピアレビューのチェックポイント}

スライドが完成したら、発表前にピアレビュー（相互レビュー）を行いましょう。以下のチェックポイントを活用してください。

\textbf{内容面のチェック:}

\begin{itemize}
\item 研究の動機が明確に伝わるか
\item 提案手法の核心的なアイデアが理解できるか
\item 実験結果が説得力を持って提示されているか
\item 結論が実験結果に基づいているか
\item スライド間の論理的なつながりがスムーズか
\end{itemize}

\textbf{デザイン面のチェック:}

\begin{itemize}
\item テキストが読みやすいか（フォントサイズ、コントラスト）
\item 図表が鮮明で、説明なしでも理解できるか
\item 色の使い方が一貫しているか
\item 余白が十分にあるか
\item アニメーションが適切か（過度でないか）
\end{itemize}

\textbf{発表面のチェック:}

\begin{itemize}
\item 制限時間内に収まるか
\item 各スライドのメッセージが明確か
\item 専門用語が適切に説明されているか
\item 質疑応答で予想される質問への準備ができているか
\end{itemize}

\subsection{AIに「聴衆として」フィードバックさせる}

人間のレビューに加えて、AIにも聴衆の視点からフィードバックを求めることができます。

\textbf{基本プロンプト:}

\begin{promptbox}
あなたは[分野名]の学会に参加している研究者です。
以下のスライド内容（テキスト版）を読み、発表に対する
フィードバックを提供してください。

【スライド内容】
スライド1: [タイトルと内容]
スライド2: [タイトルと内容]
...

以下の観点からフィードバックをお願いします:
1. ストーリーの流れは自然か？論理的な飛躍はないか？
2. 背景と動機は十分に説得力があるか？
3. 提案手法の説明で不明確な点はないか？
4. 実験結果の提示は効果的か？
5. 聴衆として、質疑応答で聞きたい質問は何か？（3つ挙げてください）
6. 全体的な改善提案
\end{promptbox}

\textbf{異なる聴衆レベルでのフィードバック:}

\begin{promptbox}
以下のスライド内容について、異なる立場の聴衆からのフィードバックを
シミュレーションしてください。

聴衆A: この分野の専門家（教授レベル）
聴衆B: 同じ学科の別分野の大学院生
聴衆C: 理系だが異なる学部の学部生

それぞれの聴衆が「理解できなかった」と感じる可能性が
高いポイントを指摘してください。

[スライド内容]
\end{promptbox}

\textbf{質疑応答の準備:}

\begin{promptbox}
以下の研究発表スライドの内容をもとに、学会の質疑応答で
想定される質問を10個生成してください。

質問は以下のカテゴリに分類してください:
1. 手法に関する技術的な質問（3問）
2. 実験設定に関する質問（2問）
3. 結果の解釈に関する質問（2問）
4. 研究の限界や将来展望に関する質問（2問）
5. 厳しい指摘・批判的な質問（1問）

各質問に対する推奨回答の骨子も示してください。

[スライド内容]
\end{promptbox}

\begin{itembox}[l]{Tip: AIと人間のフィードバックの組み合わせ}
AIからのフィードバックと人間からのフィードバックは、異なる強みを持っています。AIは論理的な一貫性や網羅性のチェックに優れていますが、発表の「温度感」や「説得力」は人間のフィードバックの方が的確です。両方を組み合わせることで、より完成度の高い発表になります。
\end{itembox}

%======================================================================
\section*{演習問題}

\begin{enumerate}

\item \textbf{演習1: 1スライド1メッセージの実践}

自分の研究（または興味のある研究トピック）について、以下を行いなさい。
\begin{enumerate}
\item 10分間の口頭発表を想定し、各スライドで伝えたいメッセージを箇条書きで12個書き出す
\item 各メッセージが「1スライド1メッセージ」の原則を満たしているか確認する（複数のメッセージが混在していたら分割する）
\item メッセージの順序を読むだけで、研究のストーリーが自然に流れるか確認する
\item 必要に応じてメッセージの追加・削除・順序変更を行い、最終的な構成案を完成させる
\end{enumerate}

\item \textbf{演習2: AIによる構成案の比較}

本章のプロンプトテンプレートを使って、2つの異なるAI（例: Claude と ChatGPT）にスライド構成案を生成させ、以下を行いなさい。
\begin{enumerate}
\item 2つの構成案の共通点と相違点を整理する
\item 各構成案の長所と短所を分析する
\item 2つの構成案の良い部分を組み合わせて、最終的な構成案を作成する
\item 最終構成案が、本章で述べた4つの検討観点（論理一貫性、網羅性、独自性、時間配分）を満たしているか確認する
\end{enumerate}

\item \textbf{演習3: Mermaid.jsによる図表作成}

自分の研究手法のパイプライン（または興味のある手法）について、以下を行いなさい。
\begin{enumerate}
\item AIにMermaid.jsのフローチャートコードを生成させる
\item Mermaid Live Editor（mermaid.live）でコードを貼り付け、図を確認する
\item 図が正確に手法を表現しているか検証し、必要に応じてコードを修正する
\item SVGまたはPNG形式でエクスポートし、スライドに挿入する
\item テキストベース図表ツールの長所と短所をまとめる
\end{enumerate}

\item \textbf{演習4: デザインAIの活用}

演習1または2で作成した構成案をもとに、以下を行いなさい。
\begin{enumerate}
\item Gamma（gamma.app）を使ってスライドの初稿を自動生成する
\item 生成されたスライドのデザインと内容を評価する（良い点、改善すべき点）
\item 内容を自分の言葉に修正する
\item デザインの改善（色、フォント、レイアウト、余白）を行う
\item PowerPointまたはKeynote形式でエクスポートし、最終的な微調整を行う
\end{enumerate}

\item \textbf{演習5: AIを活用した質疑応答準備}

完成したスライド（演習4の成果物、またはこれまでに作成したスライド）について、以下を行いなさい。
\begin{enumerate}
\item 本章のプロンプトを使って、AIに想定質問を10個生成させる
\item 各質問に対する回答を自分で準備する（AIの推奨回答を参考にしつつ、自分の言葉で）
\item 特に難しいと感じた質問3つについて、回答を声に出して練習する
\item ゼミや研究室メンバーの前でリハーサル発表を行い、実際にどのような質問が出るかを記録する
\item AIが予測した質問と実際の質問を比較し、AIによる質疑準備の有効性を評価する
\end{enumerate}

\end{enumerate}

\vspace{10mm}
\begin{itembox}[l]{本章のまとめ}
研究発表スライドの作成は、研究内容の理解と伝達力の両方が求められる総合的な作業です。AIは構成案の生成、内容の整理、図表の作成、フィードバックの提供など、多くの段階で強力な支援を提供してくれます。しかし、最も重要なのは「自分の研究を自分の言葉で伝える」ことです。AIの出力は必ず自分で検証・修正し、発表のリハーサルを十分に行いましょう。良い発表は、優れたスライドと十分な練習の両方から生まれます。
\end{itembox}

\putbib
\end{bibunit}

\begin{bibunit}
% -*- coding: utf-8 -*-
\chapter{AIによるコード生成とテスト}
\label{chap6}

\begin{chapterbox}
\textbf{【この章で学ぶこと】}
\begin{itemize}
\item \textbf{AIを活用したコード生成の基本フローとプロンプト設計を習得する} --- 自然言語からコードを生成するプロセスを理解し、効果的なプロンプトを書けるようになる
\item \textbf{テストコードの自動生成とAIコードレビューを実践できる} --- テスト駆動開発の考え方を学び、AIを活用した品質向上手法を身につける
\item \textbf{AI生成コードの品質・セキュリティリスクを理解し対処できる} --- 生成コードを盲信せず、適切に検証・修正する力を養う
\end{itemize}
\end{chapterbox}
\vspace{5mm}

%======================================================================
\section{AIとソフトウェア開発の現在}
\label{sec:6.1}
\index{こーどせいせいAI@コード生成AI}

\subsection{コード生成AIの普及状況}

2024年以降、AIによるコード生成は研究者・開発者の日常的なツールとなった。主要なツールとその特徴を以下に整理する。

\begin{table}[htbp]
\centering
\caption{主要なコード生成AIツール}
\label{tab:code-ai-tools}
\footnotesize
\begin{tabular}{lp{4cm}p{3.5cm}}
\toprule
ツール & 特徴 & 料金 \\
\midrule
GitHub Copilot & VS Code統合、コード補完に特化 & 学生無料（GitHub Education） \\
Cursor & AI統合型エディタ、対話的開発 & 無料枠あり、Pro月\$20 \\
ChatGPT / GPT-4 & 汎用的、説明付きコード生成 & 無料枠あり、Plus月\$20 \\
Gemini Pro & Google連携、長文コンテキスト & 無料枠あり \\
Claude & 長文コンテキスト、正確な推論 & 無料枠あり、Pro月\$20 \\
\bottomrule
\end{tabular}
\end{table}

\subsection{開発者のAI利用統計}

GitHubの2024年開発者調査によると、回答者の92\%がAIコーディングツールを職場内外で利用しており、70\%以上が「生産性が向上した」と回答している。特に以下の場面での活用が報告されている。

\begin{itemize}
\item \textbf{定型的なコードの作成}（ボイラープレートコード）: 87\%
\item \textbf{新しいプログラミング言語の学習}: 73\%
\item \textbf{バグの特定と修正}: 68\%
\item \textbf{テストコードの作成}: 62\%
\end{itemize}

\subsection{AIが変える開発ワークフロー}

従来の開発ワークフローでは、プログラマが仕様を読み、設計し、コードを一行ずつ書いていた。AIの登場により、このプロセスは大きく変わりつつある。

\textbf{従来のワークフロー:}
\begin{verbatim}
仕様理解 → 設計 → コーディング → テスト → デバッグ → レビュー
\end{verbatim}

\textbf{AI活用ワークフロー:}
\begin{verbatim}
仕様記述（プロンプト） → AI生成 → 人間がレビュー・修正
    → AIでテスト生成 → 検証
\end{verbatim}

重要なのは、AIは「コードを書く」部分を加速するが、「何を作るか」「なぜ作るか」を考えるのは依然として人間の仕事であるということだ。大学院生にとって、研究のロジックを明確にする力がますます重要になっている。

%======================================================================
\section{AIによるコード生成の基本}
\label{sec:6.2}
\index{こーどせいせい@コード生成}

\subsection{基本フロー}

AIによるコード生成は、以下の3ステップで行う。

\begin{verbatim}
Step 1: 自然言語で仕様を記述する（プロンプト作成）
Step 2: AIがコードを生成する
Step 3: 人間がレビュー・テスト・修正する
\end{verbatim}

このうち、最も重要なのは\textbf{Step 1（プロンプト作成）}である。指示が曖昧だと、生成されるコードも意図に沿わないものになりやすい。明確で具体的な仕様を書くことが、良いコードを得る鍵である。

\subsection{効果的なプロンプトの4要素}

AIに高品質なコードを生成させるためには、以下の4つの要素をプロンプトに含めることが重要である。

\begin{table}[htbp]
\centering
\caption{効果的なプロンプトの4要素}
\label{tab:prompt-elements}
\resizebox{\textwidth}{!}{%
\begin{tabular}{lll}
\toprule
要素 & 内容 & 例 \\
\midrule
1. 言語・フレームワーク指定 & 使用する言語やライブラリを明示 & 「Python 3.10、pandasを使用」 \\
2. 入出力仕様 & 何を受け取り、何を返すか & 「CSVファイルパスを受け取り、 \\
 & & DataFrameを返す」 \\
3. 制約条件 & パフォーマンス、エラー処理等 & 「メモリ効率を重視、 \\
 & & 欠損値はNaNで保持」 \\
4. 例示 & 具体的な入出力例 & 「入力: 'data.csv' → \\
 & & 出力: 100行×5列のDataFrame」 \\
\bottomrule
\end{tabular}}
\end{table}

\subsection{Bad prompt vs Good prompt の比較}

\textbf{Bad prompt（悪い例）:}

\begin{promptbox}
CSVを読み込む関数を書いて
\end{promptbox}

この指示では、言語もライブラリも不明であり、エラー処理の有無も分からない。AIは何らかのコードを返すが、研究の要件に合うかは運次第である。

\textbf{Good prompt（良い例）:}

\begin{promptbox}
以下の仕様でPython関数を作成してください。

\#\# 関数名
load\_experiment\_data

\#\# 目的
実験結果のCSVファイルを読み込み、前処理済みのDataFrameを返す

\#\# 使用ライブラリ
- Python 3.10
- pandas 2.0以上

\#\# 入力
- file\_path: str --- CSVファイルのパス

\#\# 出力
- pd.DataFrame --- 以下の前処理を適用済み
  1. 列名を英語の小文字に統一（スペースはアンダースコアに変換）
  2. 数値列の欠損値は列の中央値で補完
  3. 重複行を削除

\#\# エラー処理
- ファイルが存在しない場合: FileNotFoundError を発生
- CSVの解析に失敗した場合: ValueError を発生（メッセージにファイル名を含む）
\end{promptbox}

このように具体的な仕様を書くことで、AIはより正確で実用的なコードを生成できる。

\textbf{AIが生成するコードの例:}

\begin{lstlisting}[language=Python]
import pandas as pd
import os


def load_experiment_data(file_path: str) -> pd.DataFrame:
    """
    実験結果のCSVファイルを読み込み、前処理済みのDataFrameを返す。

    Parameters
    ----------
    file_path : str
        CSVファイルのパス

    Returns
    -------
    pd.DataFrame
        前処理済みのDataFrame

    Raises
    ------
    FileNotFoundError
        ファイルが存在しない場合
    ValueError
        CSVの解析に失敗した場合
    """
    # ファイル存在確認
    if not os.path.exists(file_path):
        raise FileNotFoundError(
            f"ファイルが見つかりません: {file_path}"
        )

    # CSV読み込み
    try:
        df = pd.read_csv(file_path)
    except Exception as e:
        raise ValueError(
            f"CSVの解析に失敗しました ({file_path}): {e}"
        )

    # 1. 列名を英語の小文字に統一
    df.columns = (
        df.columns
        .str.strip()
        .str.lower()
        .str.replace(" ", "_")
    )

    # 2. 数値列の欠損値を中央値で補完
    numeric_cols = df.select_dtypes(include="number").columns
    for col in numeric_cols:
        median_val = df[col].median()
        df[col] = df[col].fillna(median_val)

    # 3. 重複行を削除
    df = df.drop_duplicates().reset_index(drop=True)

    return df
\end{lstlisting}

%======================================================================
\section{主要ツールの活用法}
\label{sec:6.3}

\subsection{Gemini Pro / ChatGPT / Claude でのコード生成}

ブラウザベースのAIチャットツールは、最も手軽にコード生成を始められる方法である。

\textbf{活用のコツ:}

\begin{enumerate}
\item \textbf{新しいチャットで始める} --- 過去の文脈が混ざらないよう、コード生成タスクは新しいセッションで行う
\item \textbf{コンテキストを最初に提供する} --- 「私は大学院生で、心理学実験のデータを分析しています」のように背景を伝える
\item \textbf{段階的に依頼する} --- 一度に全てを求めず、まず骨格を作り、次にエラー処理を追加するなど段階的に進める
\end{enumerate}

\textbf{プロンプト例（Claude）:}

\begin{promptbox}
私は神経科学の大学院生です。脳波（EEG）データの解析スクリプトを
作成しています。

以下の関数を Python で作成してください。
- MNE-Python ライブラリを使用
- .fif 形式のEEGデータを読み込む
- バンドパスフィルタ（1-40Hz）を適用
- アーティファクト除去（閾値: ±100μV）を行う
- エポック分割（-0.2〜0.8秒、イベントID指定可能）を行う

各処理ステップにログ出力を含めてください。
\end{promptbox}

\subsection{GitHub Copilot の活用}
\index{GitHub Copilot@GitHub Copilot}

GitHub Copilotは、VS Code上でリアルタイムにコード補完を行うツールである。\textbf{大学院生はGitHub Educationを通じて無料で利用できる}。

\textbf{セットアップ手順:}

\begin{enumerate}
\item GitHub Education（\texttt{https://education.github.com}）で学生認証を行う
\item VS Codeの拡張機能から「GitHub Copilot」をインストール
\item GitHubアカウントでサインイン
\item \texttt{.py}ファイルを開いてコメントを書くと、自動的にコード候補が表示される
\end{enumerate}

\textbf{効果的な使い方:}

\begin{lstlisting}[language=Python]
# --- コメントで意図を書くと、Copilotが補完する ---

# CSVファイルを読み込み、欠損値の数を列ごとに表示する関数
def show_missing_values(file_path):
    # ← ここでTabキーを押すと、Copilotがコードを提案する
\end{lstlisting}

\textbf{Tips:}
\begin{itemize}
\item コメントは英語でも日本語でも良いが、具体的に書くほど精度が上がる
\item \texttt{Alt + ]} で次の候補、\texttt{Alt + [} で前の候補を表示できる
\item Copilot Chat（サイドパネル）を使えば、対話的にコードを改善できる
\end{itemize}

\subsection{対話的コード開発のテクニック}
\index{たいわてきこーどかいはつ@対話的コード開発}

AIとの対話を通じてコードを段階的に改善する方法を、「対話的コード開発」と呼ぶ。以下にその流れを示す。

\textbf{Step 1: 最初のプロンプト}
\begin{promptbox}
二群の平均値を比較する統計検定を行うPython関数を書いてください。
scipy.stats を使用してください。
\end{promptbox}

\textbf{Step 2: 結果を見て追加指示}
\begin{promptbox}
ありがとうございます。以下の改善をお願いします：
1. 正規性の検定（Shapiro-Wilk検定）を先に行い、結果に応じて
   t検定またはMann-Whitney U検定を選択する
2. 効果量（Cohen's d）も計算する
3. 結果を辞書形式で返す
\end{promptbox}

\textbf{Step 3: さらに改善}
\begin{promptbox}
結果の辞書に以下も追加してください：
- 検定名（使用した検定の名前）
- 自由度（該当する場合）
- 95\%信頼区間
\end{promptbox}

このように、段階的に要件を追加することで、自分の研究に最適なコードに近づけていく。

%======================================================================
\section{実践テクニック}
\label{sec:6.4}

\subsection{段階的開発: 関数単位で生成・検証}

AIにコード全体を一度に生成させるのではなく、\textbf{関数単位で生成し、一つずつ動作確認する}方法が最も効率的である。

\begin{verbatim}
手順:
1. まず、プログラム全体の構成を相談する
2. 各関数を一つずつ生成させる
3. 生成された関数を実行して動作確認する
4. 問題があれば修正を依頼する
5. 全関数が揃ったら統合テストを行う
\end{verbatim}

\textbf{プロンプト例:}

\begin{promptbox}
以下の研究データ分析パイプラインを設計したいです。
まず、全体の構成（関数の一覧と役割）を提案してください。
コードはまだ書かなくて良いです。

処理の流れ:
1. 複数のCSVファイルを結合
2. データのクリーニング
3. 記述統計量の算出
4. 群間比較（統計検定）
5. 結果のグラフ作成
6. レポートの自動生成
\end{promptbox}

\subsection{テスト駆動: テストケースを先に書かせる}
\index{てすとくどうかいはつ@テスト駆動開発}

テスト駆動開発（TDD）の考え方をAIと組み合わせると、コードの品質向上が期待できる。

\textbf{手順:}

\begin{verbatim}
1. 関数の仕様を決める
2. AIにテストケースを先に生成させる
3. AIに関数本体を生成させる
4. テストを実行して合格を確認する
\end{verbatim}

\textbf{プロンプト例:}

\begin{promptbox}
以下の仕様の関数について、pytestのテストケースを先に書いてください。
関数本体はまだ書かないでください。

関数名: calculate\_effect\_size
入力: group1 (list[float]), group2 (list[float])
出力: dict（cohen\_d, interpretation を含む）
仕様:
- Cohen's d を計算する
- 解釈: |d| < 0.2 → "negligible", < 0.5 → "small",
  < 0.8 → "medium", >= 0.8 → "large"
- 空リストが渡された場合は ValueError を発生
\end{promptbox}

\subsection{リファクタリング依頼}
\index{りふぁくたりんぐ@リファクタリング}

既存のコードをAIに改善させることも有用な活用法である。

\textbf{プロンプト例:}

\begin{promptbox}
以下のPythonコードをリファクタリングしてください。
改善ポイント:
1. 可読性の向上（変数名の改善、コメント追加）
2. 関数への分割（1関数1責任の原則）
3. エラー処理の追加
4. 型ヒントの追加

{[}既存のコードをここに貼り付ける]
\end{promptbox}

\subsection{デバッグ支援: エラーメッセージの活用}
\index{でばっぐ@デバッグ}

エラーに遭遇したとき、AIは頼りになるデバッグ支援ツールとなる。

\textbf{プロンプト例:}

\begin{promptbox}
以下のPythonコードを実行したところ、エラーが発生しました。
原因と修正方法を教えてください。

\#\# コード
{[}コードを貼り付け]

\#\# エラーメッセージ
Traceback (most recent call last):
  File "analysis.py", line 42, in <module>
    result = df.groupby("condition").apply(analyze\_group)
ValueError: unequal length arrays

\#\# 追加情報
- Python 3.10, pandas 2.0, scipy 1.11
- データは条件ごとに被験者数が異なります
\end{promptbox}

\textbf{Tips:} エラーメッセージだけでなく、コード全体と実行環境の情報を含めることで、AIの回答精度が大幅に向上する。

%======================================================================
\section{テストコードの自動生成}
\label{sec:6.5}
\index{てすとこーど@テストコード}

\subsection{pytestによるテストケース生成}
\index{pytest@pytest}

テストコードの作成は、研究コードの信頼性を担保する上で極めて重要である。AIを活用すれば、包括的なテストケースを短時間で生成できる。

\textbf{プロンプト例:}

\begin{promptbox}
以下のPython関数に対するpytestのテストケースを生成してください。
正常系、異常系、境界値テストを含めてください。

{[}関数のコードを貼り付け]
\end{promptbox}

\textbf{対象関数の例:}

\begin{lstlisting}[language=Python]
def normalize_scores(
    scores: list[float],
    method: str = "min-max"
) -> list[float]:
    """スコアのリストを正規化する。"""
    if not scores:
        raise ValueError("スコアリストが空です")
    if method == "min-max":
        min_val, max_val = min(scores), max(scores)
        if min_val == max_val:
            return [0.0] * len(scores)
        return [(x - min_val) / (max_val - min_val)
                for x in scores]
    elif method == "z-score":
        mean_val = sum(scores) / len(scores)
        std_val = (sum((x - mean_val) ** 2
                   for x in scores) / len(scores)) ** 0.5
        if std_val == 0:
            return [0.0] * len(scores)
        return [(x - mean_val) / std_val for x in scores]
    else:
        raise ValueError(f"未対応の正規化方法: {method}")
\end{lstlisting}

\textbf{AIが生成するテストコードの例:}

\begin{lstlisting}[language=Python]
import pytest
from normalize import normalize_scores


class TestNormalizeScoresMinMax:
    """min-max正規化のテスト"""

    def test_basic_normalization(self):
        """基本的な正規化が正しく動作すること"""
        result = normalize_scores([0, 50, 100])
        assert result == [0.0, 0.5, 1.0]

    def test_negative_values(self):
        """負の値を含む場合"""
        result = normalize_scores([-10, 0, 10])
        assert result == [0.0, 0.5, 1.0]

    def test_single_element(self):
        """要素が1つの場合（min == max）"""
        result = normalize_scores([42.0])
        assert result == [0.0]

    def test_all_same_values(self):
        """全て同じ値の場合"""
        result = normalize_scores([5.0, 5.0, 5.0])
        assert result == [0.0, 0.0, 0.0]


class TestNormalizeScoresZScore:
    """z-score正規化のテスト"""

    def test_basic_z_normalization(self):
        """基本的なz-score正規化"""
        result = normalize_scores(
            [2, 4, 6], method="z-score"
        )
        assert pytest.approx(
            result[1], abs=1e-6
        ) == 0.0


class TestNormalizeScoresErrors:
    """異常系テスト"""

    def test_empty_list(self):
        """空リストでValueError"""
        with pytest.raises(
            ValueError, match="スコアリストが空です"
        ):
            normalize_scores([])

    def test_invalid_method(self):
        """不正なメソッド名でValueError"""
        with pytest.raises(
            ValueError, match="未対応の正規化方法"
        ):
            normalize_scores([1, 2, 3], method="invalid")


class TestNormalizeScoresBoundary:
    """境界値テスト"""

    def test_very_large_values(self):
        """非常に大きな値"""
        result = normalize_scores([0, 1e10])
        assert result == [0.0, 1.0]

    def test_two_elements(self):
        """要素が2つの最小ケース"""
        result = normalize_scores([10, 20])
        assert result == [0.0, 1.0]
\end{lstlisting}

\subsection{テストカバレッジの考え方}
\index{てすとかばれっじ@テストカバレッジ}

テストカバレッジとは、テストがコードのどの程度をカバーしているかを示す指標である。

\begin{lstlisting}[language=bash]
# カバレッジの測定
pip install pytest-cov
pytest --cov=my_module --cov-report=html tests/
\end{lstlisting}

研究コードでは、最低でも以下をカバーすることを目指す。

\begin{itemize}
\item \textbf{正常系}: 期待通りの入力に対する動作
\item \textbf{異常系}: 不正な入力に対するエラー処理
\item \textbf{境界値}: 最小値、最大値、空データ、要素1つなどの特殊なケース
\end{itemize}

\subsection{AI生成テストの検証ポイント}

AIが生成したテストコードは、以下の点を必ず確認する。

\begin{enumerate}
\item \textbf{テストが実際に実行できるか} --- インポートパスやファイル名が正しいか
\item \textbf{期待値が正しいか} --- AIが計算した期待値が本当に合っているか手計算で確認する
\item \textbf{テストの独立性} --- 各テストが他のテストに依存していないか
\item \textbf{意味のあるテストか} --- \texttt{assert True} のような無意味なテストが含まれていないか
\end{enumerate}

%======================================================================
\section{AIコードレビューとリファクタリング}
\label{sec:6.6}
\index{こーどれびゅー@コードレビュー}

\subsection{AIにコードレビューを依頼するプロンプト}

AIはコードレビューの補助ツールとしても活用できる。特に一人で研究を進める大学院生にとって、客観的なフィードバックを得る貴重な手段である。

\textbf{プロンプト例:}

\begin{promptbox}
以下のPythonコードをレビューしてください。
以下の観点でチェックし、改善提案をお願いします。

1. **可読性**: 変数名、関数名は適切か。コメントは十分か。
2. **効率性**: 不要なループや非効率な処理はないか。
3. **セキュリティ**: 入力検証は適切か。脆弱性はないか。
4. **エラー処理**: 例外処理は適切か。エッジケースへの対応は十分か。
5. **保守性**: コードの構造は適切か。将来の変更に対応しやすいか。

問題の深刻度を「重大」「中程度」「軽微」の3段階で分類してください。

{[}コードをここに貼り付け]
\end{promptbox}

\subsection{コードレビューチェックリスト}

AIにレビューを依頼する際の観点をチェックリストとしてまとめる。

\begin{itembox}[l]{コードレビューチェックリスト}
\textbf{可読性}
\begin{itemize}
\item 変数名・関数名が意図を表しているか
\item 適切なコメント・docstringがあるか
\item コードの構造が論理的か
\end{itemize}

\textbf{効率性}
\begin{itemize}
\item 不要なループや計算がないか
\item 適切なデータ構造を使っているか
\item メモリ使用量は適切か
\end{itemize}

\textbf{セキュリティ}
\begin{itemize}
\item ユーザー入力を適切に検証しているか
\item ファイルパスの検証があるか
\item SQLインジェクション等の脆弱性がないか
\end{itemize}

\textbf{エラー処理}
\begin{itemize}
\item 想定されるエラーを適切にキャッチしているか
\item エラーメッセージが分かりやすいか
\item リソースの解放が確実に行われるか
\end{itemize}

\textbf{テスト可能性}
\begin{itemize}
\item 関数が適切な粒度に分割されているか
\item 外部依存を注入可能か
\item 副作用が最小限か
\end{itemize}
\end{itembox}

\subsection{リファクタリングの種類と適用場面}

\begin{table}[htbp]
\centering
\caption{リファクタリングの種類と適用場面}
\label{tab:refactoring}
\footnotesize
\begin{tabular}{lll}
\toprule
種類 & 適用場面 & プロンプト例 \\
\midrule
関数の抽出 & 長い関数を分割 & 「この関数を責務ごとに \\
 & & 小さな関数に分割して」 \\
変数名の改善 & 意味不明な変数名 & 「変数名をより意味が \\
 & & 通るように改善して」 \\
重複の排除 & 類似コードの繰り返し & 「重複しているコードを \\
 & & 共通関数にまとめて」 \\
条件分岐の簡略化 & 複雑なif-else & 「この条件分岐を \\
 & & 簡潔に書き直して」 \\
データ構造の変更 & 不適切な構造 & 「リストの代わりに辞書を \\
 & & 使うように変更して」 \\
\bottomrule
\end{tabular}
\end{table}

\subsection{ドキュメント自動生成}

AIを使って、コードのドキュメントを自動生成できる。

\textbf{docstring生成プロンプト:}

\begin{promptbox}
以下のPython関数にNumPyスタイルのdocstringを追加してください。
Parameters、Returns、Raises、Examples セクションを含めてください。

{[}関数のコードを貼り付け]
\end{promptbox}

\textbf{README生成プロンプト:}

\begin{promptbox}
以下のPythonプロジェクトのREADME.mdを作成してください。

プロジェクト名: eeg-analysis-toolkit
概要: EEGデータの前処理・分析パイプライン
主要ファイル:
- preprocess.py: データ前処理
- analyze.py: 統計分析
- visualize.py: 結果の可視化
- requirements.txt: 依存パッケージ

以下のセクションを含めてください:
1. プロジェクト概要
2. インストール方法
3. 使い方（コード例付き）
4. ファイル構成
5. ライセンス
\end{promptbox}

%======================================================================
\section{AI生成コードの品質と注意点}
\label{sec:6.7}
\index{せきゅりてぃ@セキュリティ}

\subsection{セキュリティリスク}

AI生成コードには、以下のようなセキュリティリスクが含まれる場合がある。

\textbf{1. SQLインジェクション}

\begin{lstlisting}[language=Python]
# 危険: AI が生成しがちなコード
def get_user(name):
    query = f"SELECT * FROM users WHERE name = '{name}'"
    cursor.execute(query)

# 安全: パラメータ化クエリを使用
def get_user(name):
    query = "SELECT * FROM users WHERE name = ?"
    cursor.execute(query, (name,))
\end{lstlisting}

\textbf{2. 不適切な入力検証}

\begin{lstlisting}[language=Python]
# 危険: ファイルパスを検証していない
def read_file(path):
    with open(path) as f:
        return f.read()

# 安全: パスの検証を含む
def read_file(path, allowed_dir="/data/experiments"):
    real_path = os.path.realpath(path)
    if not real_path.startswith(
        os.path.realpath(allowed_dir)
    ):
        raise ValueError("許可されていないパスです")
    with open(real_path) as f:
        return f.read()
\end{lstlisting}

\textbf{3. ハードコードされた秘密情報}

\begin{lstlisting}[language=Python]
# 危険: APIキーがコードに含まれている
api_key = "sk-1234567890abcdef"

# 安全: 環境変数から読み込む
import os
api_key = os.environ.get("API_KEY")
if not api_key:
    raise RuntimeError(
        "API_KEY 環境変数が設定されていません"
    )
\end{lstlisting}

\subsection{ライセンス問題}
\index{らいせんす@ライセンス}

AIが生成したコードには、学習データに含まれるオープンソースコードの断片が含まれる可能性がある。

\begin{itemize}
\item \textbf{研究コードの場合}: 論文の再現性のために公開する場合、ライセンスを明記する
\item \textbf{推奨}: MITライセンスやApache 2.0ライセンスなど、許容的なライセンスを選択する
\item \textbf{注意}: AI生成コードの著作権は法的にグレーゾーンであることを認識しておく
\end{itemize}

\subsection{「信頼するが検証する」の原則}

AI生成コードに対する基本姿勢は\textbf{「Trust, but verify（信頼するが検証する）」}である。

\begin{itembox}[l]{検証チェックリスト}
\begin{itemize}
\item コードが意図通りに動作するか（手動テスト）
\item テストケースが全て通過するか（自動テスト）
\item エッジケースでも正しく動作するか
\item エラー処理が適切か
\item パフォーマンスが許容範囲内か
\item セキュリティ上の問題がないか
\item コードの論理が自分で理解できるか $\leftarrow$ \textbf{最も重要}
\end{itemize}
\end{itembox}

最後の項目が最も重要である。\textbf{自分で理解できないコードは使わない}。AIに説明を求め、納得してから使用すること。

\subsection{研究コードの再現性確保}
\index{さいげんせい@再現性}

研究で使用するコードは、再現性が欠かせない。以下の対策を講じる。

\begin{enumerate}
\item \textbf{依存関係の固定}: \texttt{requirements.txt}や\texttt{pyproject.toml}でバージョンを固定する
\item \textbf{乱数シードの設定}: 結果の再現性のために、乱数シードを明示的に設定する
\item \textbf{環境情報の記録}: Pythonバージョン、OS、主要ライブラリのバージョンを記録する
\item \textbf{バージョン管理}: Gitでコードの履歴を管理する（第8章で詳述）
\end{enumerate}

\begin{lstlisting}[language=Python]
# 再現性確保のテンプレート
import random
import numpy as np

RANDOM_SEED = 42
random.seed(RANDOM_SEED)
np.random.seed(RANDOM_SEED)

# 実行環境の記録
import sys
import platform
print(f"Python: {sys.version}")
print(f"NumPy: {np.__version__}")
print(f"OS: {platform.platform()}")
\end{lstlisting}

%======================================================================
\section*{演習問題}

\begin{enumerate}

\item \textbf{演習6-1: 基本的なコード生成}

以下の仕様で、AIを使ってPython関数を生成しなさい。生成されたコードを実行し、正しく動作することを確認しなさい。

\textbf{仕様:}
\begin{itemize}
\item 関数名: \texttt{summarize\_data}
\item 入力: pandasのDataFrame
\item 出力: 各数値列について、平均、中央値、標準偏差、最小値、最大値を含む辞書
\item エラー処理: DataFrameが空の場合はValueErrorを発生
\end{itemize}

使用したプロンプトと、生成されたコードを提出すること。

\item \textbf{演習6-2: テスト駆動開発}

演習6-1で作成した\texttt{summarize\_data}関数に対して、AIを使ってpytestのテストケースを生成しなさい。以下のケースを最低限含めること。

\begin{itemize}
\item 正常なDataFrameに対するテスト
\item 空のDataFrameに対するテスト
\item 数値列のないDataFrameに対するテスト
\item 欠損値を含むDataFrameに対するテスト
\end{itemize}

\item \textbf{演習6-3: コードレビューとリファクタリング}

以下のコードをAIにレビューさせ、改善版を作成しなさい。どのような問題が指摘され、どう改善されたかをレポートにまとめること。

\begin{lstlisting}[language=Python]
def proc(d):
    r = []
    for i in range(len(d)):
        if d[i] != None and d[i] > 0:
            x = d[i] * 2.5 + 10
            r.append(x)
    s = 0
    for j in r:
        s = s + j
    a = s / len(r)
    return a
\end{lstlisting}

\item \textbf{演習6-4: デバッグ支援}

以下のコードにはバグが含まれている。AIを使ってバグを特定し、修正しなさい。修正の過程で使用したプロンプトも提出すること。

\begin{lstlisting}[language=Python]
import pandas as pd

def calculate_growth_rate(df, value_col, time_col):
    """時系列データの成長率を計算する"""
    df = df.sort_values(time_col)
    growth_rates = []
    for i in range(len(df)):
        current = df.iloc[i][value_col]
        previous = df.iloc[i-1][value_col]
        rate = (current - previous) / previous * 100
        growth_rates.append(rate)
    df["growth_rate"] = growth_rates
    return df

# テストデータ
data = pd.DataFrame({
    "year": [2020, 2021, 2022, 2023],
    "revenue": [100, 120, 150, 180]
})

result = calculate_growth_rate(data, "revenue", "year")
print(result)
\end{lstlisting}

\item \textbf{演習6-5: 研究データ分析スクリプトの作成}

自分の研究テーマ（または以下のサンプルテーマ）に関連するデータ分析スクリプトをAIと対話的に作成しなさい。

\textbf{サンプルテーマ}: 「オンライン授業と対面授業の学習効果比較」

以下の機能を含むスクリプトを作成すること。
\begin{enumerate}
\item サンプルデータの生成（被験者50名、2群、テストスコアあり）
\item 記述統計量の算出
\item 適切な統計検定の実施
\item 結果のグラフ作成
\end{enumerate}

AIとの対話履歴（プロンプトとレスポンスの要約）、最終的なコード、実行結果を提出すること。

\end{enumerate}

\vspace{10mm}
\begin{itembox}[l]{本章のまとめ}
AIは頼りになるコーディングパートナーであるが、最終的な品質責任は使用者にある。「AIに書かせて、人間が検証する」という姿勢を常に持ち、特に研究コードでは再現性とセキュリティに十分注意すること。プロンプトの書き方一つでコードの品質が大きく変わるため、明確で具体的な仕様記述のスキルを磨くことが重要である。
\end{itembox}

\putbib
\end{bibunit}

\begin{bibunit}
\chapter{データ分析と可視化}
\label{chap7}

\begin{chapterbox}
\textbf{【この章で学ぶこと】}
\begin{itemize}
\item AIを活用したデータ分析ワークフロー（前処理→分析→可視化）を実践できる
\item 論文品質のグラフを効率的に作成できる
\item 統計的検定の基礎と適切な手法選択を理解する
\end{itemize}
\end{chapterbox}
\vspace{5mm}

\section{研究におけるデータ分析}
\index{でーたぶんせき@データ分析}

\subsection{なぜデータ分析が重要か}

大学院での研究において、データ分析\index{でーたぶんせき@データ分析}は以下の場面で欠かせない。

\begin{enumerate}
\item \textbf{実験結果の解釈}: 収集したデータから意味のあるパターンを見出す
\item \textbf{仮説の検証}: 統計的手法により、仮説が支持されるかを客観的に判断する
\item \textbf{論文の説得力}: 適切な分析と可視化は、論文の査読者を説得する論文の説得力を大きく左右する
\item \textbf{新たな発見}: 探索的データ分析（EDA）により、予想外の知見が得られることがある
\end{enumerate}

データ分析のスキルは、理系・文系を問わず、現代の大学院生に求められる基本能力である。心理学の実験データ、社会調査の回答データ、生物学の測定データ、工学のシミュレーション結果——あらゆる研究分野でデータ分析が行われている。

\subsection{AIがデータ分析をどう変えるか}

従来、データ分析を行うには、プログラミング言語（Python, R等）の習得に多くの時間を要した。AIの登場により、このハードルは大幅に下がっている。

\begin{itembox}[l]{AIが特に役立つ場面}
\begin{itemize}
\item \textbf{分析コードの生成}: 「このデータでt検定を行いたい」と伝えるだけで、コードが得られる
\item \textbf{手法の選択支援}: データの特性を伝えると、適切な統計手法を提案してくれる
\item \textbf{グラフの作成}: 「論文に載せる棒グラフを作って」という指示で、出版品質のグラフコードが得られる
\item \textbf{結果の解釈}: 統計結果の意味を平易な言葉で説明してもらえる
\end{itemize}
\end{itembox}

ただし、\textbf{AIはあくまで道具であり、分析の設計と結果の解釈は研究者の責任である}。本章では、AIの力を正しく活用しつつ、データ分析の基礎力を養うことを目指す。

\section{データ分析ワークフロー}
\index{わーくふろー@ワークフロー}

\subsection{Step 1: データの理解と前処理}

データ分析の第一歩は、手元のデータを正しく理解することである。

\begin{promptbox}
以下のCSVファイルを読み込み、データの概要を把握する\\
Pythonスクリプトを作成してください。

ファイル: experiment\_results.csv\\
使用ライブラリ: pandas

以下の情報を出力してください:
\begin{enumerate}
\item データの形状（行数×列数）
\item 各列のデータ型
\item 最初の5行の表示
\item 基本統計量（数値列: 平均・標準偏差・最小・最大、\\
カテゴリ列: ユニーク数・最頻値）
\item 欠損値の数と割合（列ごと）
\item 各数値列のヒストグラム
\end{enumerate}
\end{promptbox}

AIが生成するコードの例:

\begin{lstlisting}[language=Python]
import pandas as pd
import matplotlib.pyplot as plt
import numpy as np


def explore_data(file_path: str) -> pd.DataFrame:
    """CSVファイルを読み込み、データの概要を表示する"""

    # データ読み込み
    df = pd.read_csv(file_path)

    # 1. データの形状
    print("=" * 50)
    print(f"データの形状: {df.shape[0]}行 × {df.shape[1]}列")
    print("=" * 50)

    # 2. 各列のデータ型
    print("\n【データ型】")
    print(df.dtypes.to_string())

    # 3. 最初の5行
    print("\n【最初の5行】")
    print(df.head().to_string())

    # 4. 基本統計量
    print("\n【基本統計量（数値列）】")
    print(df.describe().to_string())

    # 5. 欠損値
    print("\n【欠損値】")
    missing = df.isnull().sum()
    missing_pct = (missing / len(df) * 100).round(2)
    missing_df = pd.DataFrame({
        "欠損数": missing,
        "欠損率(%)": missing_pct
    })
    print(missing_df[missing_df["欠損数"] > 0].to_string())
    if missing.sum() == 0:
        print("  欠損値なし")

    # 6. ヒストグラム
    num_cols = df.select_dtypes(include=[np.number]).columns
    if len(num_cols) > 0:
        fig, axes = plt.subplots(
            1, len(num_cols),
            figsize=(5 * len(num_cols), 4))
        if len(num_cols) == 1:
            axes = [axes]
        for ax, col in zip(axes, num_cols):
            df[col].hist(ax=ax, bins=20, edgecolor="black")
            ax.set_title(col)
            ax.set_xlabel("値")
            ax.set_ylabel("頻度")
        plt.tight_layout()
        plt.savefig("data_overview_histograms.png", dpi=150)
        plt.show()

    return df
\end{lstlisting}

\begin{itembox}[l]{欠損値処理の主な方法}
\begin{table}[H]
\centering
\footnotesize
\begin{tabular}{lll}
\toprule
方法 & 適用場面 & Pythonコード \\
\midrule
行の削除 & 欠損が少ない場合 & \texttt{df.dropna()} \\
平均値補完 & 正規分布に近いデータ & \texttt{df.fillna(df.mean())} \\
中央値補完 & 外れ値がある場合 & \texttt{df.fillna(df.median())} \\
前方補完 & 時系列データ & \texttt{df.ffill()} \\
多重代入法 & 欠損が多い場合 & \texttt{IterativeImputer} \\
\bottomrule
\end{tabular}
\end{table}
\end{itembox}

\begin{promptbox}
以下のDataFrameの数値列について、外れ値を検出する\\
コードを作成してください。\\
外れ値の定義: IQR法（Q1 - 1.5*IQR 未満、\\
または Q3 + 1.5*IQR 超過）\\
外れ値がある行を表示し、箱ひげ図で可視化してください。
\end{promptbox}

\subsection{Step 2: 探索的データ分析（EDA）}
\index{たんさくてきでーたぶんせき@探索的データ分析}

探索的データ分析（Exploratory Data Analysis, EDA）\index{EDA@EDA}は、データの構造やパターンを視覚的に把握するプロセスである。

\begin{promptbox}
以下のDataFrameについて、探索的データ分析（EDA）を行う\\
Pythonコードを作成してください。

分析内容:
\begin{enumerate}
\item 数値列間の相関行列の計算とヒートマップの作成
\item 各数値列の分布（ヒストグラム + 正規分布フィット）
\item カテゴリ列ごとのグループ別統計量
\item 主要な変数間の散布図マトリクス
\end{enumerate}

使用ライブラリ: pandas, matplotlib, seaborn\\
グラフは日本語フォントに対応させてください。
\end{promptbox}

\subsection{Step 3: 統計的検定}
\index{とうけいてきけんてい@統計的検定}

研究の仮説を検証するためには、適切な統計的検定を行う必要がある。以下に、よく使われる検定手法とその選び方を示す。

\begin{itembox}[l]{統計的検定の選び方フローチャート}
\begin{verbatim}
データは正規分布に従うか？
├── はい → 2群比較か？
│   ├── はい → 対応のあるデータか？
│   │   ├── はい → 対応のあるt検定
│   │   └── いいえ → 独立2標本t検定
│   └── いいえ → 3群以上の比較
│       ├── 対応あり → 反復測定分散分析
│       └── 対応なし → 一元配置分散分析
└── いいえ → 2群比較か？
    ├── はい → 対応のあるデータか？
    │   ├── はい → Wilcoxon符号順位検定
    │   └── いいえ → Mann-Whitney U検定
    └── いいえ → Kruskal-Wallis検定
カテゴリデータ同士 → カイ二乗検定
\end{verbatim}
\end{itembox}

\begin{promptbox}
以下のデータについて、2群間の平均値の差を検定する\\
Pythonコードを作成してください。

データ構造:
\begin{itemize}
\item DataFrame に ``group'' 列（``control'' と ``treatment''）と ``score'' 列がある
\item 各群の被験者数は異なる可能性がある
\end{itemize}

処理手順:
\begin{enumerate}
\item 各群のデータの正規性を Shapiro-Wilk 検定で確認
\item 等分散性を Levene 検定で確認
\item 結果に応じて適切な検定を自動選択
\item 効果量（Cohen's d または r）を計算
\item 結果を分かりやすくレポート出力
\end{enumerate}
\end{promptbox}

\begin{itembox}[l]{p値の正しい解釈}
\index{pち@$p$値}
\begin{table}[H]
\centering
\begin{tabular}{p{4.5cm}p{5cm}}
\toprule
誤った解釈 & 正しい解釈 \\
\midrule
$p = 0.03$は「差がある確率が97\%」 & 帰無仮説の下で、観測データ以上に極端な結果が得られる確率が3\% \\
$p < 0.05$は「重要な差がある」 & 統計的に有意だが、実用的な重要性は効果量で判断 \\
$p = 0.06$は「差がない」 & 有意水準を超えただけで、差がないとは限らない \\
$p$値が小さいほど効果が大きい & $p$値はサンプルサイズの影響を受ける。効果量を見るべき \\
\bottomrule
\end{tabular}
\end{table}
\end{itembox}

\section{論文品質のグラフ作成}
\index{ぐらふさくせい@グラフ作成}

\subsection{matplotlib / seaborn の基本}

Pythonでグラフを作成するための2大ライブラリがmatplotlib\index{matplotlib@matplotlib}とseaborn\index{seaborn@seaborn}である。

\begin{itemize}
\item \textbf{matplotlib}: 低レベルの描画ライブラリ。細かいカスタマイズが可能
\item \textbf{seaborn}: matplotlibをベースにした高レベルの可視化ライブラリ。統計的な可視化に特化し、美しいデフォルト設定を持つ
\end{itemize}

\subsection{効果的なグラフの選び方}

\begin{table}[H]
\centering
\small
\begin{tabular}{lll}
\toprule
データの性質 & 推奨グラフ & 用途 \\
\midrule
カテゴリ別の量的比較 & 棒グラフ（エラーバー付き） & 群間比較の結果表示 \\
2変数の関係 & 散布図 & 相関関係の可視化 \\
分布の比較 & 箱ひげ図 / バイオリンプロット & 群間の分布の違い \\
変数間の関連 & ヒートマップ & 相関行列の可視化 \\
時系列データ & 折れ線グラフ & 時間変化の表示 \\
割合・構成比 & 積み上げ棒グラフ & カテゴリの内訳表示 \\
\bottomrule
\end{tabular}
\end{table}

\subsection{論文に必要なグラフの要素}

学術論文や学会発表のグラフには、以下の要素が不可欠である。

\begin{itembox}[l]{必須要素}
\begin{itemize}
\item 明確な軸ラベル（単位を含む）
\item 適切なタイトル（Figure 1: ... の形式）
\item 凡例（複数系列の場合）
\item エラーバー（平均値を示す場合、標準誤差または95\%信頼区間）
\item 適切なフォントサイズ（本文と同程度、最低8pt）
\item 高解像度（300dpi以上）
\end{itemize}
\end{itembox}

\subsection{白黒印刷対応のグラフ設計}

多くの学術論文は白黒で印刷されることがあるため、色だけに頼らないデザインが重要である。

\begin{enumerate}
\item \textbf{ハッチパターン}: 棒グラフにストライプや格子模様を加える
\item \textbf{マーカーの形状}: 散布図で○、△、□など異なるマーカーを使う
\item \textbf{線のスタイル}: 折れ線グラフで実線、破線、点線を使い分ける
\item \textbf{コントラスト}: 色を使う場合も、白黒にしたときに区別できるコントラストを確保する
\end{enumerate}

\subsection{AIでグラフ作成コードを生成するプロンプト例}

\begin{promptbox}
以下の仕様で論文品質の棒グラフを作成する\\
Pythonコードを書いてください。

データ:
\begin{itemize}
\item 3群（Control, Treatment A, Treatment B）
\item 各群の平均値と標準誤差がある
\end{itemize}

グラフの要件:
\begin{itemize}
\item matplotlib + seaborn を使用
\item 白背景、枠線あり
\item エラーバー（標準誤差）付き
\item 各棒に異なるハッチパターン（白黒印刷対応）
\item 有意差を示すブラケットと * マーク
\item フォントサイズ: タイトル14pt、軸ラベル12pt、目盛り10pt
\item 解像度 300dpi で PNG 保存
\item 図のサイズ: 幅8cm × 高さ6cm（単カラム論文用）
\end{itemize}
\end{promptbox}

\section{インタラクティブ可視化}
\index{いんたらくてぃぶかしか@インタラクティブ可視化}

\subsection{Plotlyの紹介}

Plotly\index{Plotly@Plotly}は、インタラクティブなグラフを作成するライブラリである。マウスオーバーで値を表示したり、ズーム・パンができるため、データの探索段階で特に有用である。

\begin{lstlisting}[language=Python]
import plotly.express as px

# インタラクティブ散布図
fig = px.scatter(
    df, x="variable_a", y="variable_b",
    color="group",
    hover_data=["subject_id", "age", "score"],
    title="変数Aと変数Bの関係",
    labels={
        "variable_a": "変数A",
        "variable_b": "変数B"
    }
)
fig.update_layout(
    font=dict(size=14),
    hoverlabel=dict(font_size=12)
)
fig.write_html("interactive_scatter.html")
fig.show()
\end{lstlisting}

\subsection{データ探索での活用}

インタラクティブ可視化は、以下の場面で特に有用である。

\begin{itemize}
\item \textbf{外れ値の特定}: マウスオーバーで外れ値の詳細情報を確認できる
\item \textbf{パターンの発見}: ズーム機能で特定の領域を拡大して観察できる
\item \textbf{プレゼンテーション}: ゼミ発表でリアルタイムにデータを探索できる
\item \textbf{共同研究者との共有}: HTMLファイルを送るだけで、相手もインタラクティブに操作できる
\end{itemize}

ただし、\textbf{論文に掲載するグラフは静的な画像（matplotlib等）で作成する}こと。インタラクティブ可視化はあくまで探索・確認用である。

\section{AIを使ったデータ分析の落とし穴}

\subsection{統計的な誤解}
\index{pはっきんぐ@$p$ハッキング}

AIがデータ分析を手軽にする一方で、統計的な誤解に基づく分析が行われるリスクも増大している。

\begin{itembox}[l]{p-hacking（p値ハッキング）}
多数の検定を繰り返し、有意な結果が出たものだけを報告する行為。AIを使うと、様々な分析を手軽に試せるため、このリスクが高まる。

有意水準5\%で20回検定すれば、本当は差がなくても平均1回は「有意」になる（偽陽性）。
\end{itembox}

\begin{itembox}[l]{多重比較問題}
\index{たじゅうひかく@多重比較}
複数の群を同時に比較する場合、検定の回数が増えるほど偽陽性の確率が高くなる。3群を比較する場合はBonferroni補正（補正後の有意水準: $0.05 / 3 = 0.0167$）やTukey HSDなどの多重比較法を用いる。
\end{itembox}

\subsection{過学習への注意}
\index{かがくしゅう@過学習}

AIが提案する分析手法が必要以上に複雑になることがある。特に機械学習を使う場合、過学習（overfitting）のリスクに注意が必要である。

対策:
\begin{itemize}
\item 交差検証（Cross-validation）を必ず行う
\item 単純なモデルから始め、必要な場合にのみ複雑化する
\item サンプルサイズに見合った分析手法を選ぶ
\end{itemize}

\subsection{AIが出した統計結果を鵜呑みにしない}

AIが生成した分析コードは、一見正しく見えても論理的な誤りを含むことがある。

\begin{itembox}[l]{自分で確認すべきポイント}
\begin{itemize}
\item 検定の前提条件は満たされているか（正規性、等分散性）
\item 実験計画（対応あり/なし）に合った検定を使っているか
\item 有意水準は事前に設定されているか
\item 効果量も報告しているか
\item 多重比較を行う場合、補正は適切か
\item サンプルサイズは十分か
\item 結果の解釈は論理的か
\end{itemize}
\end{itembox}

\section*{演習問題}

\begin{enumerate}
\item[\textbf{演習7-1}] \textbf{データの前処理と概要把握}\\
サンプルデータを生成し、AIを使ってデータの概要把握と前処理を行いなさい。データの概要把握（形状、型、分布、欠損値）、外れ値の検出と対処、欠損値の処理、処理前後の比較を実施し、レポートにまとめなさい。

\item[\textbf{演習7-2}] \textbf{探索的データ分析}\\
演習7-1のデータについて、グループごとの記述統計量、相関分析、各変数の分布の可視化を実施し、発見した特徴やパターンを3つ以上報告しなさい。

\item[\textbf{演習7-3}] \textbf{適切な統計検定の実施}\\
「treatment群はcontrol群と比べて、pre\_scoreからpost\_scoreへの変化量が大きい」という仮説を検定しなさい。正規性の検定、適切な統計検定の選択と実施、効果量の計算を行い、結果を論文の「結果」セクションの書き方で報告しなさい。

\item[\textbf{演習7-4}] \textbf{論文品質のグラフ作成}\\
演習7-3の結果を以下のグラフで可視化しなさい。全てのグラフは論文品質（300dpi、適切なフォントサイズ、軸ラベル）で作成すること。(1) 群ごとのpre/postスコアの棒グラフ（エラーバー付き）、(2) 群ごとの変化量の箱ひげ図（個別データ点付き）、(3) pre\_scoreとpost\_scoreの散布図（群ごとに色分け、回帰直線付き）。

\item[\textbf{演習7-5}] \textbf{総合分析レポート}\\
自分の研究テーマについて、データ分析の全工程を実施し、レポートにまとめなさい。サンプルデータの生成、前処理とEDA、適切な統計分析、論文品質のグラフ3枚以上、結果のまとめ（発見した知見を3つ以上）を含めること。AIとの対話履歴の要約も合わせて提出すること。
\end{enumerate}

\putbib
\end{bibunit}

\begin{bibunit}
\chapter{GitHubによる研究成果管理}
\label{chap8}

\begin{chapterbox}
\textbf{【この章で学ぶこと】}
\begin{itemize}
\item Git/GitHubの基本操作を習得する
\item AIを活用したWebページ作成とGitHub Pagesでの公開を実践できる
\item 研究ポートフォリオの構築と活用方法を理解する
\end{itemize}
\end{chapterbox}
\vspace{5mm}

\section{なぜ大学院生がGitHubを使うべきか}
\index{GitHub@GitHub}

\subsection{研究コードのバージョン管理}

大学院での研究では、コードを何度も修正する。バージョン管理\index{ばーじょんかんり@バージョン管理}なしでは、以下のような問題が生じる。

\begin{verbatim}
よくある悲劇:
analysis.py
analysis_v2.py
analysis_v2_final.py
analysis_v2_final_revised.py
analysis_v2_final_revised_REAL_FINAL.py
\end{verbatim}

このようなファイル管理では、「どの時点で何を変更したか」「なぜ変更したか」が全く分からなくなる。Gitを使えば、コミットした変更の履歴を追跡でき、いつでも過去の状態に戻れる。

\subsection{共同研究での活用}

研究は一人で行うものではない。GitHubを使えば、次のようなことができる。

\begin{itemize}
\item \textbf{コードの共有}: URLを送るだけで共同研究者がコードにアクセスできる
\item \textbf{同時作業}: 複数人が同じプロジェクトを同時に編集できる
\item \textbf{変更の追跡}: 誰が、いつ、何を変更したかが一目で分かる
\item \textbf{レビュー}: 共同研究者の変更をレビューしてからマージできる
\end{itemize}

\subsection{就職活動でのポートフォリオ効果}

技術系の就職活動において、GitHubのプロフィールは「第二の履歴書」になる。特にIT企業やデータサイエンス関連のポジションでは、GitHubでの活動が評価される。

\begin{itembox}[l]{GitHubプロフィールで評価されるポイント}
\begin{itemize}
\item 定期的なコミット（継続的に開発・研究している証拠）
\item README.md が整備されたリポジトリ（ドキュメンテーション能力）
\item コードの品質（命名規則、コメント、テストの有無）
\item オープンソースプロジェクトへの貢献
\end{itemize}
\end{itembox}

\subsection{オープンサイエンスへの貢献}
\index{おーぷんさいえんす@オープンサイエンス}

近年、研究のオープン化（Open Science）が世界的なトレンドとなっている。論文のデータや分析コードをGitHubで公開することで、以下のメリットがある。

\begin{itemize}
\item \textbf{研究の透明性向上}: 分析手法が第三者に検証可能になる
\item \textbf{再現性の担保}: 他の研究者が同じ分析を再現できる
\item \textbf{引用の増加}: 公開されたコードやデータは他の研究者に利用され、引用が増える傾向がある
\item \textbf{学術コミュニティへの貢献}: 研究ツールやデータセットの共有
\end{itemize}

\section{Gitの基礎}
\index{Git@Git}

\subsection{バージョン管理とは何か}

バージョン管理とは、ファイルの変更履歴を記録し、過去の任意の時点の状態を復元できるようにする仕組みである。

\subsection{リポジトリ、コミット、ブランチの基本概念}

Gitを理解するための3つの核心概念を、ゲームのアナロジーで説明する。

\begin{itembox}[l]{Gitの基本概念}
\index{りぽじとり@リポジトリ}\index{こみっと@コミット}\index{ぶらんち@ブランチ}

\textbf{リポジトリ（Repository）= セーブデータフォルダ}

プロジェクトの全ファイルとその変更履歴を格納する場所である。

\textbf{コミット（Commit）= セーブポイント}

ある時点でのファイルの状態を保存する操作である。各コミットには「何を変更したか」を説明するメッセージを付ける。

\textbf{ブランチ（Branch）= 並行世界}

メインの開発ラインから分岐して、独立に作業できる仕組みである。
\end{itembox}

\begin{table}[H]
\centering
\caption{Git用語とゲームのアナロジー}
\begin{tabular}{lll}
\toprule
Git 用語 & ゲームのアナロジー & 説明 \\
\midrule
リポジトリ & セーブデータフォルダ & プロジェクト全体の格納場所 \\
コミット & セーブポイント & ある時点の状態を記録 \\
ブランチ & 並行ルート & 分岐して独立に作業 \\
マージ & ルートの合流 & 分岐した作業を統合 \\
プッシュ & クラウドに同期 & ローカルの変更をGitHubに送る \\
プル & クラウドから同期 & GitHubの変更をローカルに取得 \\
クローン & セーブデータをコピー & GitHubからプロジェクトをダウンロード \\
\bottomrule
\end{tabular}
\end{table}

\section{GitHubの基本操作}

\subsection{GitHubアカウントの作成}

\begin{enumerate}
\item \texttt{https://github.com} にアクセスする
\item 「Sign up」をクリックし、メールアドレス・ユーザー名・パスワードを入力する
\item メール認証を完了する
\end{enumerate}

ユーザー名は本名またはそれに近い名前が望ましい（就職活動で使うため）。GitHub Education（\texttt{https://education.github.com}）に大学のメールアドレスで学生認証を行うと、GitHub Copilotが無料で使えるなどの特典がある。

\subsection{GitHub Desktop のセットアップ}
\index{GitHub Desktop@GitHub Desktop}

コマンドラインが苦手な人には、GitHub Desktop（GUIアプリ）が推奨される。\texttt{https://desktop.github.com} からダウンロードし、GitHubアカウントでサインインする。

\subsection{リポジトリの作成とクローン}

GitHub上でリポジトリを作成する手順:
\begin{enumerate}
\item GitHub にログイン
\item 右上の「+」→「New repository」をクリック
\item Repository name、Description を入力し、Public または Private を選択
\item 「Add a README file」にチェックし、「Create repository」をクリック
\end{enumerate}

\subsection{コミット→プッシュの基本フロー}

日常的な作業の流れは以下の通りである。

\begin{enumerate}
\item ファイルを編集する（VS Code等で）
\item 変更をステージングする（保存したい変更を選択）
\item コミットメッセージを書く
\item コミットする（ローカルにセーブ）
\item プッシュする（GitHubにアップロード）
\end{enumerate}

\begin{itembox}[l]{良いコミットメッセージの書き方}
\begin{table}[H]
\centering
\begin{tabular}{ll}
\toprule
悪い例 & 良い例 \\
\midrule
「更新」 & 「t検定の結果を棒グラフで可視化」 \\
「修正」 & 「欠損値処理のバグを修正: NaNが残る問題」 \\
「追加」 & 「相関分析のコードとヒートマップを追加」 \\
\bottomrule
\end{tabular}
\end{table}
コミットメッセージは「何を、なぜ変更したか」が分かるように書く。
\end{itembox}

\subsection{VS Code との連携}

VS Code\index{VS Code@VS Code}にはGitの操作機能が組み込まれている。左サイドバーの「ソース管理」アイコンからファイルのステージング、コミット、プッシュが可能である。拡張機能として\textbf{GitLens}（コードの各行の変更履歴表示）や\textbf{Git Graph}（コミット履歴のグラフ表示）の導入を推奨する。

\section{AIでWebページを作成する}
\index{Webページ@Webページ}

\subsection{AIでHTML/CSSコードを生成する}

AIを使えば、プログラミングの知識がなくても自己紹介Webページを作成できる。

\begin{promptbox}
大学院生の自己紹介Webページを作成するHTML/CSSコードを書いてください。

内容:
\begin{itemize}
\item 名前: 田中太郎
\item 所属: ○○大学大学院 情報科学研究科 修士2年
\item 研究テーマ: 「深層学習を用いた自然言語処理の研究」
\item スキル: Python, PyTorch, 統計分析
\end{itemize}

デザイン要件:
\begin{itemize}
\item モダンでシンプルなデザイン
\item レスポンシブ対応（スマホでも見やすい）
\item セクション: ヘッダー、自己紹介、研究概要、業績、スキル、連絡先
\item CSSは同一ファイルに含める（1ファイルで完結）
\end{itemize}
\end{promptbox}

\subsection{VS Codeでの編集とプレビュー}

生成されたHTMLをVS Codeで編集する手順は以下の通りである。

\begin{enumerate}
\item VS Codeでプロジェクトフォルダを開く
\item \texttt{index.html}という名前でファイルを作成
\item AIが生成したコードを貼り付ける
\item 拡張機能「Live Server」をインストール
\item \texttt{index.html}を右クリック→「Open with Live Server」
\item ブラウザでリアルタイムにプレビューが表示される
\end{enumerate}

\subsection{AIでカスタマイズする}

Webページの修正やカスタマイズもAIに依頼できる。デザインの変更やセクションの追加など、自然言語で指示するだけで対応してもらえる。

\section{GitHub Pagesで公開する}
\index{GitHub Pages@GitHub Pages}

\subsection{GitHub Pagesとは}

GitHub Pagesは、GitHubのリポジトリに置いたHTMLファイルを無料でWebサイトとして公開できるサービスである。

\begin{itemize}
\item 無料で利用可能
\item \texttt{https://ユーザー名.github.io/リポジトリ名/} のURLで公開される
\item リポジトリにプッシュするだけで自動的に更新される
\item カスタムドメインも設定可能
\end{itemize}

\subsection{GitHub Pages の設定手順}

\begin{enumerate}
\item GitHubで新しいリポジトリを作成する（Publicに設定）
\item リポジトリに \texttt{index.html} をアップロードする
\item リポジトリの「Settings」→「Pages」をクリック
\item 「Source」でBranch: \texttt{main}、Folder: \texttt{/ (root)} を設定し「Save」
\item 数分後に \texttt{https://ユーザー名.github.io/リポジトリ名/} で公開される
\end{enumerate}

\subsection{更新フロー}

ローカルで\texttt{index.html}を編集し、Live Serverでプレビュー確認した後、GitHub Desktopでコミット・プッシュすると、数分後にWebサイトに反映される。

\section{研究ポートフォリオの構築}
\index{ぽーとふぉりお@ポートフォリオ}

\subsection{ポートフォリオに含めるもの}

\begin{table}[H]
\centering
\begin{tabular}{llc}
\toprule
セクション & 内容 & 重要度 \\
\midrule
自己紹介 & 名前、所属、研究分野、顔写真 & 必須 \\
研究概要 & 研究テーマの説明（一般向け） & 必須 \\
業績一覧 & 論文、学会発表、受賞歴 & 必須 \\
スキル & プログラミング言語、ツール、資格 & 推奨 \\
プロジェクト & 研究プロジェクトの詳細と成果物 & 推奨 \\
連絡先 & メールアドレス、SNS & 必須 \\
CV/履歴書 & PDF形式でダウンロード可能 & 推奨 \\
\bottomrule
\end{tabular}
\end{table}

\subsection{README.md の書き方}

README.md\index{README@README.md}は、リポジトリの「顔」である。GitHubでリポジトリを開いたとき最初に表示されるため、丁寧に書くことが重要である。プロジェクト概要、特徴、インストール方法、使い方、ファイル構成、依存ライブラリ、ライセンス、著者情報を含めるとよい。

\subsection{フォルダ構成のベストプラクティス}

\begin{verbatim}
my-research-project/
├── README.md
├── LICENSE
├── requirements.txt
├── .gitignore
├── src/              # ソースコード
├── data/
│   ├── raw/          # 生データ（変更不可）
│   └── processed/    # 処理済みデータ
├── notebooks/        # Jupyter Notebook
├── figures/          # 生成されたグラフ
├── tests/            # テストコード
└── docs/             # ドキュメント
\end{verbatim}

\begin{itembox}[l]{.gitignore ファイル}
\texttt{.gitignore}は、Gitで追跡すべきでないファイルを指定するファイルである。大きなデータファイル、Pythonのキャッシュ（\texttt{\_\_pycache\_\_/}）、環境変数（\texttt{.env}）、OS固有ファイル（\texttt{.DS\_Store}）などを指定する。
\end{itembox}

\section{共同研究でのGitHub活用}

\subsection{Issue による課題管理}
\index{Issue@Issue}

Issueは、プロジェクトの課題やタスクを管理する機能である。タイトルと説明を記入し、ラベル（bug, enhancement, question等）を付与し、担当者を指定する。研究プロジェクトでは、データ前処理、統計分析、図表作成などのタスクをIssueとして管理できる。

\subsection{Pull Request によるコードレビュー}
\index{Pull Request@Pull Request}

Pull Request（PR）は、ブランチの変更をメインブランチに統合する前に、レビューを行う仕組みである。

\begin{enumerate}
\item 新しいブランチを作成する
\item ブランチ上でコードを書いてコミットする
\item GitHub上でPull Requestを作成する
\item 共同研究者がコードをレビューする
\item レビューが通ったらマージする
\end{enumerate}

\subsection{研究コードの共有とライセンス}
\index{らいせんす@ライセンス}

研究コードを公開する場合、適切なライセンスを選択することが重要である。

\begin{table}[H]
\centering
\begin{tabular}{lll}
\toprule
ライセンス & 特徴 & 推奨場面 \\
\midrule
MIT & 最も自由。商用利用・改変・再配布可能 & 研究コードの一般公開 \\
Apache 2.0 & MIT + 特許条項 & ツール・ライブラリの公開 \\
GPL v3 & 派生物も同じライセンスで公開必須 & 強いオープンソースを望む場合 \\
CC BY 4.0 & クリエイティブ・コモンズ & データセットの公開 \\
\bottomrule
\end{tabular}
\end{table}

研究コードには\textbf{MITライセンス}が最も広く使われている。特別な事情がなければMITが無難な選択肢である。

\section*{演習問題}

\begin{enumerate}
\item[\textbf{演習8-1}] \textbf{GitHubアカウントとリポジトリの作成}\\
GitHubアカウントを作成し、GitHub Educationに登録しなさい。\texttt{my-first-repo}という名前のリポジトリを作成し、GitHub Desktopでクローンした後、テキストファイルを1つ追加してコミット・プッシュしなさい。スクリーンショットを撮影して提出すること。

\item[\textbf{演習8-2}] \textbf{AIでWebページを作成する}\\
AIチャットツールを使って、自己紹介Webページを作成しなさい。自分の情報を含むプロンプトを作成し、AIにHTML/CSSコードを生成させ、VS Codeで編集してLive Serverでプレビューすること。最低2回はAIにカスタマイズを依頼して改善しなさい。

\item[\textbf{演習8-3}] \textbf{GitHub Pagesで公開する}\\
演習8-2で作成したWebページをGitHub Pagesで公開しなさい。公開URLにアクセスして表示を確認し、公開URLとスクリーンショットを提出すること。

\item[\textbf{演習8-4}] \textbf{README.md の作成}\\
自分の研究プロジェクト（または架空の研究プロジェクト）について、AIを活用してREADME.mdを作成しなさい。プロジェクト概要、インストール方法、使い方（コード例付き）、ファイル構成、ライセンスを含めること。

\item[\textbf{演習8-5}] \textbf{共同作業の実践}\\
2人以上のグループで共有リポジトリを作成し、Issueを3つ以上作成して担当者を割り当て、各メンバーがブランチを作成してPull Requestとコードレビューを行いなさい。作業の流れと学んだことをレポートにまとめること。
\end{enumerate}

\putbib
\end{bibunit}

\begin{bibunit}
\chapter{就活文書の作成}
\label{chap9}

\begin{chapterbox}
\textbf{【この章で学ぶこと】}
\begin{itemize}
\item AIを活用してガクチカ・ESを効果的に作成するワークフローを習得する
\item STAR法による経験の構造化技法を身につける
\item 研究内容を非専門家に分かりやすく伝える技法を理解する
\end{itemize}
\end{chapterbox}
\vspace{5mm}

\section{大学院生の就職活動とAI}
\index{しゅうしょくかつどう@就職活動}

\subsection*{M1の就活スケジュール}

大学院修士課程の学生にとって、就職活動は研究と並行して進めなければならない重要なタスクである。

\begin{table}[H]
\centering
\begin{tabular}{ll}
\toprule
時期 & 活動内容 \\
\midrule
M1 4月〜6月 & 自己分析、業界研究の開始 \\
M1 6月〜8月 & サマーインターンシップ応募・参加 \\
M1 9月〜11月 & 秋冬インターンシップ、OB/OG訪問 \\
M1 12月〜2月 & 本選考に向けたES準備、早期選考開始 \\
M1 3月〜 & 本選考エントリー開始（広報解禁） \\
M2 4月〜6月 & 面接・内定（選考解禁は6月1日が建前） \\
\bottomrule
\end{tabular}
\end{table}

\subsection*{企業が大学院生に求めるもの}

企業が大学院生を採用する際に注目するポイントは、大きく分けて2つある。

\begin{enumerate}
\item \textbf{専門性} --- 専門知識そのものだけでなく、「深く学ぶ力」の証明。研究を通じて培った分析力、論理的思考力。技術的な課題に粘り強く取り組む姿勢。
\item \textbf{汎用スキル} --- 課題発見力、仮説構築力、コミュニケーション力、プロジェクト管理力、主体性。
\end{enumerate}

重要なのは、企業は「研究の専門知識をそのまま使いたい」のではなく、「研究を通じて培われた思考プロセスと行動特性」を評価しているということである。

\subsection*{AIを就活で活用する意義と注意点}

\begin{itembox}[l]{AI活用の意義と注意点}
\textbf{意義}
\begin{itemize}
\item 自己分析の壁打ち相手として、客観的な視点を得られる
\item 文章の構造化・推敲を効率的に行える
\item 想定質問の生成により、面接準備の抜け漏れを減らせる
\item 専門用語の平易化など、表現の調整に役立つ
\end{itemize}

\textbf{注意点}
\begin{itemize}
\item AIが生成した文章をそのまま提出すると、面接で自分の言葉で語れなくなる
\item 複数の就活生が同じAIを使うため、没個性的な文章になりやすい
\item AIは「あなたの経験」を知らない。事実の入力は必ず自分で行う必要がある
\item 企業によってはAI使用を検出する仕組みを導入している場合がある
\end{itemize}
\end{itembox}

本章で提唱するアプローチは、\textbf{「AIに書かせる」のではなく「AIと一緒に磨く」}というスタンスである。

\section{ガクチカとは何か}
\index{がくちか@ガクチカ}

\subsection*{「学生時代に力を入れたこと」の本質}

「ガクチカ」とは「学生時代に力を入れたこと」の略称で、ほぼすべての企業のESや面接で問われる定番質問である。ガクチカが問うているのは、\textbf{「何をやったか」ではなく「どのように取り組んだか」}である。

\subsection*{企業が見ているポイント}

\begin{table}[H]
\centering
\begin{tabular}{lll}
\toprule
ポイント & 説明 & 具体的に見ていること \\
\midrule
課題発見力 & 状況の中から問題を見つける力 & 何を課題と認識したか \\
行動力 & 自ら考えて動く力 & どのような工夫をしたか \\
成果 & 行動の結果として何が得られたか & 定量的・定性的な成果 \\
学び & 経験から何を学んだか & 自己成長、気づき \\
再現性 & 入社後も同様に活躍できるか & 学んだことをどう活かすか \\
\bottomrule
\end{tabular}
\end{table}

\subsection*{研究をガクチカにする場合の特有の課題}

\begin{enumerate}
\item \textbf{専門用語の壁}: 研究内容をそのまま書くと、専門外の採用担当者には理解できない
\item \textbf{「やらされた」感の払拭}: 与えられたテーマの中で、自分がどのような課題意識を持ち、どう工夫したかを明確にする必要がある
\item \textbf{成果の表現}: 研究の社会的意義や、その成果が何に繋がるのかまで含めて説明する必要がある
\item \textbf{個人の貢献の明確化}: 研究室での研究はチームで行うことも多く、自分の独自の貢献を明確にしなければならない
\end{enumerate}

\section{AIを使ったガクチカ作成ワークフロー}
\index{STAR法@STAR法}

ガクチカ作成を4つのステップに分け、各ステップでAIをどのように活用するかを解説する。

\subsection{Step 1: 素材の棚卸し --- AIに壁打ちしてもらう}

最初のステップは、自分の経験を漏れなく洗い出すことである。AIを壁打ち相手にすることで、忘れていた経験や見落としていた側面を引き出せる。

\begin{promptbox}
\# 役割\\
あなたは大学院生の就職活動を支援するキャリアカウンセラーです。

\# 依頼\\
私は修士課程の学生で、ガクチカの素材を探しています。以下の情報をもとに、私の経験を深掘りする質問を10個生成してください。

\# 私の情報
\begin{itemize}
\item 専攻: [例: 機械工学]
\item 研究テーマ: [例: 自動運転における安全性評価手法の研究]
\item 研究期間: [例: 2024年4月〜現在]
\item その他の活動: [例: TAの経験、学会発表1回、研究室のゼミ運営]
\end{itemize}

\# 質問の方向性: 研究の中での困難や工夫、チームワーク、自発的な取り組み、失敗と学び、スキルアップの努力に関する質問
\end{promptbox}

このステップでは、AIに文章を書かせるのではなく、AIに「質問してもらう」ことが重要である。

\subsection{Step 2: STAR法での構造化}

素材が見つかったら、STAR法\index{STAR法@STAR法}で構造化する。STAR法は、経験を説得力のある形で整理するためのフレームワークである。

\begin{table}[H]
\centering
\begin{tabular}{lll}
\toprule
要素 & 内容 & ガクチカにおける役割 \\
\midrule
\textbf{S}ituation（状況） & どんな環境・背景だったか & 読み手に場面をイメージさせる \\
\textbf{T}ask（課題） & 何が課題・目標だったか & 取り組みの意義を示す \\
\textbf{A}ction（行動） & 具体的に何をしたか & 自分の主体性・工夫を示す \\
\textbf{R}esult（結果） & どんな成果が得られたか & 行動の効果を証明する \\
\bottomrule
\end{tabular}
\end{table}

\subsection{Step 3: AIで文章化（400字程度）}

STAR構造ができたら、これを読みやすい文章に変換する。

\begin{promptbox}
\# 役割\\
あなたはESライティングの専門家です。

\# 依頼\\
以下のSTAR構造を、400字以内のガクチカ文章にまとめてください。

\# STAR構造\\
{[Step 2で作成したSTAR構造をここに貼る]}

\# 文章の条件
\begin{itemize}
\item 400字以内（厳守）
\item 最初の1文で「何に力を入れたか」が分かること
\item 専門用語は最小限にし、非専門家でも理解できること
\item 数字を使って具体性を出すこと
\item 最後に「学び」や「今後への展望」を含めること
\end{itemize}
\end{promptbox}

\begin{itembox}[l]{Before/Afterの例}
\textbf{Before（構造化前の素材メモ）}:

自動運転の安全性の研究をしています。GSNというのを使って安全性を評価する手法を提案しました。論文を50本くらい読みました。国際学会に通りました。

\textbf{After（AIと磨いた完成稿）}:

自動運転システムの安全性を論理的に証明する研究に力を入れました。修士研究として取り組んだのは、「なぜこのシステムは安全と言えるのか」を体系的に説明する枠組みの構築です。従来の安全性評価は、確認項目が断片的で全体像が見えにくいという課題がありました。私は50本以上の先行研究を精読して課題を整理した上で、安全性の論拠を木構造で可視化する手法と、システムの状態を継続的に記録する仕組みを組み合わせた新しいアプローチを考案しました。週1回の研究ミーティングでは、自ら論点を整理した資料を作成し、議論を主導しました。その結果、この成果をまとめた論文が安全性工学分野の国際学会に採択され、研究室初の海外口頭発表を実現しました。この経験を通じて、先行研究の徹底的な調査に基づく課題設定と、仮説を検証しながら新しい解決策を構築する力を身につけました。（350字）
\end{itembox}

\subsection{Step 4: AIで推敲・添削}

文章化ができたら、採用担当者の視点でAIに添削してもらう。

\begin{promptbox}
\# 役割\\
あなたはメーカーの採用担当者（技術系）です。年間500本以上のESを読んでいます。

\# 依頼\\
以下のガクチカ文章を、採用担当者の視点で添削してください。

\# 添削の観点
\begin{enumerate}
\item 【第一印象】最初の1文で興味を持てるか
\item 【課題の明確さ】何が課題だったかが明確か
\item 【行動の具体性】何をしたかが具体的に書かれているか
\item 【主体性】「自分が」やったことが伝わるか
\item 【成果】成果が具体的か（定量情報があるか）
\item 【再現性】入社後も活かせそうな能力が見えるか
\item 【読みやすさ】文章は読みやすいか、冗長な部分はないか
\end{enumerate}

\# 出力形式\\
各観点について、○/△/×の評価と、具体的な改善提案
\end{promptbox}

添削のサイクルは2〜3回繰り返すと効果的である。重要なのは、\textbf{AIの修正案をそのまま採用するのではなく、自分の言葉で書き直すこと}である。面接では必ずガクチカの深掘りがあり、自分の言葉で語れなければ見抜かれる。

\section{研究を非専門家に伝える技法}

\subsection*{エレベーターピッチ（30秒で研究を説明）}
\index{えれべーたーぴっち@エレベーターピッチ}

エレベーターピッチとは、エレベーターに乗っている30秒程度の時間で、自分の研究を相手に伝える技法である。

\begin{itembox}[l]{30秒エレベーターピッチの構造}
\begin{enumerate}
\item \textbf{一言で言うと}（1文）：研究を一言で表現
\item \textbf{なぜ重要か}（1〜2文）：社会的な文脈での意義
\item \textbf{何をしたか}（1〜2文）：アプローチの概要
\item \textbf{何が分かったか}（1文）：主要な成果
\end{enumerate}
\end{itembox}

\subsection*{技術的な内容の平易化テクニック}

\begin{table}[H]
\centering
\begin{tabular}{ll}
\toprule
専門的な表現 & 平易な表現（比喩） \\
\midrule
Goal Structuring Notation & 安全性の根拠を木構造で整理する図 \\
形式検証 & 数学的に「絶対に間違いがない」と証明すること \\
機械学習モデルの過学習 & 丸暗記したけど応用問題が解けない状態 \\
API & ソフトウェア同士が会話するための共通言語 \\
\bottomrule
\end{tabular}
\end{table}

\begin{itemize}
\item \textbf{テクニック1: 日常の比喩を使う} --- 専門概念を身近なものに例える
\item \textbf{テクニック2: 「何が嬉しいか」を先に言う} --- How（何をしたか）よりWhy/What（なぜ重要か、何が嬉しいか）を先に伝える
\item \textbf{テクニック3: 数字で具体化する} --- 「多くの先行研究」ではなく「50本以上の論文を精読」
\end{itemize}

\section{ES全体でのAI活用}
\index{えんとりーしーと@エントリーシート}

\subsection*{志望動機の作成}

志望動機は、ガクチカ以上に「自分の言葉」が求められる設問である。ポイントは、AIに「文章」ではなく「論理構造」を出力させることである。構造ができたら、自分の言葉で文章化する。

\subsection*{自己PRの作成}

自己PRはガクチカと役割が異なる。

\begin{itemize}
\item \textbf{ガクチカ}: 具体的なエピソードを通じて能力を示す（帰納的）
\item \textbf{自己PR}: 自分の強みを宣言し、エピソードで裏付ける（演繹的）
\end{itemize}

\subsection*{注意: AIに丸投げしない}

\textbf{ESは面接の「台本」である。}ESに書いた内容は、面接で必ず深掘りされる。AIを使って作成する場合でも、以下を心がけること。

\begin{enumerate}
\item \textbf{素材（事実・経験）は必ず自分で入力する} --- AIは知らないことを捏造する
\item \textbf{AIの出力を鵜呑みにしない} --- 自分の経験と合っているか確認する
\item \textbf{最終的な文章は自分の言葉に置き換える} --- 音読して違和感がないか確認する
\item \textbf{面接で話す練習をする} --- ESを見ずに、自分の言葉で3分間話せるか試す
\end{enumerate}

\section{面接準備でのAI活用}

\subsection*{想定質問の生成}

面接準備で最も効果的なAI活用は、想定質問の生成である。

\begin{promptbox}
\# 役割\\
あなたはメーカー（技術系）の面接官です。10年以上の面接経験があります。

\# 依頼\\
以下のESの内容を読み、面接で聞きそうな質問を15個生成してください。

\# 質問のカテゴリ
\begin{itemize}
\item ガクチカの深掘り（5問）
\item 志望動機の深掘り（3問）
\item 研究内容に関する質問（3問）
\item 価値観・人柄に関する質問（2問）
\item 逆質問のヒント（2問）
\end{itemize}

各質問に対して、質問文、この質問の意図、回答のポイントを記載してください。
\end{promptbox}

\subsection*{模擬面接（AIに面接官役をさせる）}

ChatGPTやClaudeなどの対話型AIを使って、模擬面接を行うことができる。1回につき1つの質問をし、回答に対してフィードバックを述べてから次の質問に進む形式が効果的である。

\subsection*{面接後の振り返り}

面接が終わった後の振り返りにもAIは活用できる。質問と自分の回答を記録し、各回答の論理構造、具体性、質問の意図への対応をAIに分析してもらう。

\vspace{5mm}
\begin{itembox}[l]{本章のキーポイント}
\begin{enumerate}
\item \textbf{就活文書は「AIに書かせる」のではなく「AIと一緒に磨く」} --- 素材は自分、構造化と推敲にAIを活用する
\item \textbf{STAR法は汎用的な構造化ツール} --- Situation→Task→Action→Resultの順で整理すれば説得力が増す
\item \textbf{非専門家への説明は「Why/What」から始める} --- 「何をしたか」より「なぜ重要か」を先に伝える
\item \textbf{ESは面接の台本} --- AIの出力をそのまま使わず、自分の言葉で語れるようにする
\item \textbf{AIを使った模擬面接は効果的だが、人間との練習も併用する}
\end{enumerate}
\end{itembox}

\section*{演習問題}

\begin{enumerate}
\item[\textbf{演習9-1}] \textbf{素材の棚卸し}\\
本章のStep 1のプロンプトを使い、AIに自分の経験を引き出す質問を生成してもらいなさい。生成された質問に対してすべて回答し、ガクチカの素材となりそうなエピソードを3つ選びなさい。

\item[\textbf{演習9-2}] \textbf{STAR法による構造化}\\
演習9-1で選んだエピソードの1つを、STAR法で構造化しなさい。AIを使ってもよいが、最終的な構造は自分の言葉で記述すること。Situation, Task, Action, Resultそれぞれについて、3〜5文で記述しなさい。

\item[\textbf{演習9-3}] \textbf{ガクチカ文章の作成と推敲}\\
演習9-2のSTAR構造をもとに、400字以内のガクチカ文章を作成しなさい。AIに下書きを作成させた後、Step 4の添削プロンプトで少なくとも2回の推敲サイクルを行い、Before（初稿）とAfter（最終稿）を比較して提出しなさい。

\item[\textbf{演習9-4}] \textbf{エレベーターピッチの作成}\\
自分の研究について、30秒で説明できるエレベーターピッチを作成しなさい。以下の2つのバージョンを作成すること。
\begin{itemize}
\item 版A: 同じ分野の研究者向け（専門用語使用可）
\item 版B: 文系の採用担当者向け（専門用語不使用）
\end{itemize}
実際に声に出して読み、30秒以内に収まるか確認しなさい。

\item[\textbf{演習9-5}] \textbf{模擬面接の実施}\\
本章のプロンプトを使い、AIと模擬面接を実施しなさい。8問以上のやり取りを行い、AIからのフィードバックをもとに、自分の回答の改善点を3つ挙げなさい。
\end{enumerate}

\putbib
\end{bibunit}

\begin{bibunit}
\chapter{研究の新規性分析と総合演習}
\label{chap10}

\begin{chapterbox}
\textbf{【この章で学ぶこと】}
\begin{itemize}
\item AIを使って自分の研究の新規性を客観的に分析・言語化できる
\item これまでの章の成果物を統合した研究ポートフォリオを完成させる
\item AI活用の全体像を振り返り、今後の活用方針を立てられる
\end{itemize}
\end{chapterbox}
\vspace{5mm}

\section{なぜ新規性の言語化が重要か}
\index{しんきせい@新規性}

研究の価値は、突き詰めれば「何が新しいのか」に集約される。しかし、多くの大学院生にとって、自分の研究の新規性を明確に言語化することは難しい。新規性の言語化が求められる場面は、研究生活の中で繰り返し登場する。

\subsection*{論文投稿（査読対応）}

学術論文を投稿すると、査読者から必ず問われるのが「この研究の新規性は何か（What is the novelty of this work?）」である。Introductionで先行研究との差分が明確に記述されていなければ、どれほど実験結果が良くても、論文は採択されない。

\subsection*{学会発表（Introductionでの位置づけ）}

学会発表のスライドでは、冒頭のIntroductionで「既存研究の何が足りなくて、自分の研究が何を解決するのか」を明確に示す必要がある。

\subsection*{就職面接での研究説明}

第\ref{chap9}章で述べた通り、就職面接では「あなたの研究のオリジナリティは何ですか？」という質問が頻出である。

\subsection*{研究計画書の作成}

博士課程への進学を志望する場合や、科研費などの研究費申請を行う場合、研究計画書の中で「提案する研究の新規性」を説得力を持って記述しなければならない。

これらすべての場面において、\textbf{自分の研究の新規性を客観的に分析し、適切な言葉で表現する能力}が求められる。

\section{新規性分析の3つの観点}

研究の新規性は、一つの軸だけで語れるものではない。ここでは、新規性を分析するための3つの観点を紹介する。

\subsection{観点1: 方法論的新規性 --- 新しい手法・アプローチの提案}
\index{ほうほうろんてきしんきせい@方法論的新規性}

最も典型的な新規性の形態である。既存の手法では解決できない問題に対して、新しいアルゴリズム、フレームワーク、モデル、評価手法などを提案するものである。

方法論的新規性を主張する際の表現パターン:
\begin{itemize}
\item ``We propose a novel approach that combines X and Y to address Z''
\item ``Unlike existing methods that ..., our method ...''
\item ``To the best of our knowledge, this is the first work that ...''
\end{itemize}

\subsection{観点2: 実践的新規性 --- 実世界への適用・検証}
\index{じっせんてきしんきせい@実践的新規性}

必ずしも手法そのものが新しくなくても、既存手法を新しい領域・対象に適用し、その有効性を検証することは重要な新規性である。

実践的新規性を主張する際の表現パターン:
\begin{itemize}
\item ``We apply X to the domain of Y for the first time''
\item ``We validate the effectiveness of X through a real-world case study of Y''
\item ``Our empirical evaluation reveals that ...''
\end{itemize}

\subsection{観点3: 時間的観点 --- 最新の研究動向の中での位置づけ}

研究は常に時間軸の中に存在する。最新の技術動向、規格の改定、社会的ニーズの変化に対応することも、重要な新規性の一形態である。

時間的観点での新規性を主張する際の表現パターン:
\begin{itemize}
\item ``With the recent advancement of X, there is an urgent need for Y''
\item ``In light of the latest revision of standard Z, we propose ...''
\item ``The emerging challenge of X necessitates a new approach to Y''
\end{itemize}

\subsection*{3つの観点の関係}

多くの優れた研究は、これら3つの観点のうち複数を兼ね備えている。

\begin{table}[H]
\centering
\begin{tabular}{lll}
\toprule
観点 & 問いかけ & 強い場合の特徴 \\
\midrule
方法論的新規性 & 手法のどこが新しいか？ & 新アルゴリズム、新フレームワーク \\
実践的新規性 & どこに適用したのが新しいか？ & 新領域への適用、大規模実験 \\
時間的観点 & なぜ「今」この研究が必要か？ & 最新動向への対応、新規格対応 \\
\bottomrule
\end{tabular}
\end{table}

\section{AIを使った新規性分析ワークフロー}

自分の研究の新規性を客観的に分析するために、3つのステップからなるワークフローを提案する。

\subsection{Step 1: 研究の位置づけマッピング}

\begin{promptbox}
\# 役割\\
あなたは[研究分野名]の専門家であり、学術論文の査読経験が豊富な研究者です。

\# 依頼\\
以下の研究の学術的な位置づけを分析してください。

\# 研究の概要
\begin{itemize}
\item 分野: [例: 自動運転の安全性工学]
\item テーマ: [例: GSNと安全状態レポートを統合した安全性保証手法の提案]
\item 主要なキーワード: [例: Safety Case, GSN, Autonomous Driving, ODD]
\item 使用した手法: [例: GSNによる安全ケース構築、安全状態レポートとの統合]
\item 主な結果: [例: ケーススタディで有効性を検証]
\end{itemize}

\# 分析してほしい内容
\begin{enumerate}
\item この研究が属する学術分野・サブ分野の俯瞰図
\item 関連する研究コミュニティ（学会、ジャーナル）
\item この分野の主要な研究潮流（3〜5つ）
\item 各潮流における自分の研究の位置づけ
\item 隣接する研究分野との接点
\end{enumerate}
\end{promptbox}

\textbf{注意}: AIの回答はあくまで「仮説」として扱うこと。特に、具体的な論文名や著者名が挙げられた場合は、必ず実際にその論文が存在するか確認する必要がある（第3章で学んだハルシネーション対策）。

\subsection{Step 2: 既存研究との差分分析}

\begin{promptbox}
\# 役割\\
あなたは学術論文の査読者です。

\# 依頼\\
以下の「自分の研究」と「先行研究」を比較し、差分を分析してください。

\# 分析してほしい内容
\begin{enumerate}
\item 各先行研究と自分の研究の共通点
\item 各先行研究と自分の研究の相違点
\item 先行研究全体に共通する「限界」で、自分の研究が解決しているもの
\item 比較表（行: 比較項目、列: 各研究）
\end{enumerate}
\end{promptbox}

このような比較表は、論文のIntroductionやRelated Workセクションで非常に効果的である。AIに下書きを作成させた後、自分で正確性を検証し、必要に応じて修正する。

\subsection{Step 3: 新規性の言語化}

\begin{promptbox}
\# 役割\\
あなたは学術英語ライティングの専門家です。

\# 依頼\\
新規性分析の結果を、英語論文のIntroductionの最後に置く「本研究の貢献」パラグラフとしてまとめてください。

\# 条件
\begin{itemize}
\item 3文で構成（各文が1つの貢献に対応）
\item 文1: 方法論的新規性（何を提案したか）
\item 文2: 実践的新規性（何で検証したか）
\item 文3: 時間的観点（なぜ今重要か）
\item ``The main contributions of this paper are as follows:'' で始めること
\end{itemize}
\end{promptbox}

\section{総合演習: 研究ポートフォリオの完成}

本章の後半は、これまでの章で作成してきた成果物を統合し、研究ポートフォリオとして完成させる総合演習である。

\subsection*{これまでの成果物の振り返り}

\begin{table}[H]
\centering
\begin{tabular}{lll}
\toprule
章 & 成果物 & 内容 \\
\midrule
第3章 & 文献調査レポート & 先行研究のサーベイ \\
第4章 & 英語Abstract & 英文要旨（150〜250語） \\
第5章 & 発表スライド & 学会発表やゼミ発表用 \\
第6章 & 研究コード & 実験・分析に使用したコード \\
第7章 & データ分析・グラフ & データの可視化と分析結果 \\
第8章 & GitHub Webページ & 研究成果を公開するWebページ \\
第9章 & ガクチカ & 就職活動用のES文章 \\
第10章 & 新規性分析レポート & 研究の新規性を整理した文書 \\
\bottomrule
\end{tabular}
\end{table}

\subsection*{GitHubリポジトリでの整理と公開}

これらの成果物を、一つのGitHubリポジトリとして体系的に整理する。

\begin{verbatim}
my-research-portfolio/
├── README.md               ← リポジトリの概要
├── docs/
│   ├── literature-review/  ← 文献調査レポート
│   ├── abstract/           ← 英語Abstract
│   ├── novelty-analysis/   ← 新規性分析レポート
│   └── job-search/         ← ガクチカ（公開範囲に注意）
├── slides/                 ← 発表スライド
├── src/                    ← 研究コード
├── data/                   ← データ分析・グラフ
├── .gitignore
└── LICENSE
\end{verbatim}

\section{AI活用の振り返りと今後}

\subsection*{AIが特に効果的だった場面}

\begin{table}[H]
\centering
\begin{tabular}{lll}
\toprule
場面 & 効果 & 理由 \\
\midrule
文章の構造化 & 非常に高い & パターン認識と整理が得意 \\
英語表現の改善 & 高い & 大量の英語テキストで学習 \\
コードの雛形生成 & 高い & 定型的なパターンの再現が得意 \\
アイデアの壁打ち & 高い & 多角的な視点を提供できる \\
想定質問の生成 & 高い & 多様なパターンを網羅的に生成 \\
\bottomrule
\end{tabular}
\end{table}

\subsection*{AIの限界が明らかだった場面}

\begin{table}[H]
\centering
\begin{tabular}{lll}
\toprule
場面 & 限界 & 対処法 \\
\midrule
事実の正確性 & ハルシネーションのリスク & 必ず原典で検証 \\
最新情報の把握 & 学習データの時点で止まる & 最新論文は自分で確認 \\
研究の独創性 & 既存パターンの組み合わせが限界 & 本質的な着想は人間の役割 \\
実験の設計 & 分野固有の暗黙知が不足 & 指導教員や先輩に相談 \\
感情・動機の表現 & 表面的になりがち & 自分の言葉で書く \\
\bottomrule
\end{tabular}
\end{table}

\subsection*{今後のAI進化への対応}

\textbf{ツールは変わっても、原理は変わらない}。本書で学んだ以下の原理は、どのようなAIツールが登場しても適用できる。

\begin{itembox}[l]{AI活用の4つの普遍的原理}
\begin{enumerate}
\item \textbf{明確な指示の重要性} --- ROCF（Role, Output, Context, Format）の考え方は、あらゆるAIツールに応用できる。
\item \textbf{段階的な作業プロセス} --- 複雑なタスクを段階に分けて進めることで、人間が各段階で検証・修正できる。
\item \textbf{批判的検証の習慣} --- 99\%正確なAIの残り1\%の誤りを見つけられるのは、専門知識を持つ人間だけである。
\item \textbf{人間の判断が必要な領域の認識} --- 研究の方向性を決める、仮説を立てる、実験結果を解釈する、倫理的な判断をする——これらはAIが「代替」できない。
\end{enumerate}
\end{itembox}

\subsection*{「AIを使いこなす」研究者になるために}

AIの普及により、「AIを使える人」と「使えない人」の間に生産性の格差が生まれている。しかし、より本質的な格差は、\textbf{「AIに使われる人」と「AIを使いこなす人」}の間にある。

\textbf{「AIに使われる人」の特徴}:
\begin{itemize}
\item AIの出力をそのまま使う
\item AIなしでは作業できなくなる
\item AIの回答が正しいかどうか判断できない
\end{itemize}

\textbf{「AIを使いこなす人」の特徴}:
\begin{itemize}
\item AIの出力を批判的に検証し、取捨選択する
\item AIを使って自分の思考を加速・拡張する
\item AIの限界を理解し、適切な場面で活用する
\end{itemize}

AIは道具である。素晴らしい道具であるが、あくまで道具である。あなたの研究は、あなたにしかできない知的営みである。AIはその営みを支える強力な助手であり、本書がその助手との付き合い方を学ぶ一助となれば幸いである。

\vspace{5mm}
\begin{itembox}[l]{本章のキーポイント}
\begin{enumerate}
\item \textbf{新規性は3つの観点で分析する} --- 方法論的新規性、実践的新規性、時間的観点
\item \textbf{新規性分析は段階的に行う} --- 位置づけマッピング→差分分析→言語化の3ステップ
\item \textbf{成果物は統合して初めて価値が最大化する} --- GitHubリポジトリとして体系的に整理する
\item \textbf{ツールは変わっても原理は変わらない} --- 明確な指示、段階的プロセス、批判的検証、人間の判断
\item \textbf{「AIを使いこなす」研究者を目指す} --- AIは助手であり、研究の本質は人間の知的営みにある
\end{enumerate}
\end{itembox}

\section*{演習問題}

\begin{enumerate}
\item[\textbf{演習10-1}] \textbf{新規性の3観点分析}\\
自分の研究（または取り組んでいるテーマ）について、方法論的新規性、実践的新規性、時間的観点の3つの観点から新規性を分析しなさい。各観点について、3〜5文で記述すること。

\item[\textbf{演習10-2}] \textbf{先行研究との比較表の作成}\\
自分の研究テーマに関する先行研究を3本以上取り上げ、比較表の形式で自分の研究との差分を整理しなさい。比較項目は5つ以上設定すること。AIを使って下書きを作成してもよいが、各セルの内容の正確性は自分で検証すること。

\item[\textbf{演習10-3}] \textbf{貢献記述の作成}\\
演習10-1と10-2の結果をもとに、英語論文のIntroductionで使える「本研究の貢献」パラグラフ（3文）を作成しなさい。日本語訳も併記すること。

\item[\textbf{演習10-4}] \textbf{研究ポートフォリオの構築}\\
本書の各章で作成した成果物を含むGitHubリポジトリを作成しなさい。README.mdを英語で記述し、リポジトリの構成、各成果物の概要、実行方法（コードがある場合）を記載すること。

\item[\textbf{演習10-5}] \textbf{AI活用方針書の作成}\\
本書で学んだ内容を踏まえ、今後の研究活動におけるAI活用方針書（A4用紙1枚程度）を作成しなさい。以下の項目を含めること。
\begin{itemize}
\item 自分の研究活動でAIを活用する場面のリスト
\item 各場面でのAI活用の具体的な方法
\item AIに任せない（人間が判断する）領域の明記
\item AI使用の倫理的ガイドライン（自分なりの基準）
\item 半年後に見直す予定日
\end{itemize}
\end{enumerate}

\putbib
\end{bibunit}

% 付録
\appendix
\chapter{主要AIツールリファレンス}
\label{app:tools}

\begin{quote}
\textbf{注記}: AIツールの開発・更新は非常に速いペースで進んでいる。本付録の情報は2026年3月時点のものであり、最新のツール情報・料金体系については、本書のWebサポートページを参照されたい。
\end{quote}

\section{対話型AI（汎用）}
\index{ChatGPT@ChatGPT}\index{Claude@Claude}\index{Gemini@Gemini}

\begin{table}[H]
\centering
\footnotesize
\resizebox{\textwidth}{!}{%
\begin{tabular}{llll}
\toprule
ツール名 & 開発元 & 料金 & 得意分野 \\
\midrule
ChatGPT & OpenAI & 無料版あり/Plus \$20/月 & 汎用的な対話、文章・コード生成 \\
Claude & Anthropic & 無料版あり/Pro \$20/月 & 長文の分析・生成、論理的推論 \\
Gemini & Google & 無料版あり/Advanced \$20/月 & Google連携、マルチモーダル \\
Copilot & Microsoft & 無料版あり/Pro \$20/月 & Office連携、Web検索統合 \\
\bottomrule
\end{tabular}}
\end{table}

\textbf{選び方のポイント}: 研究用途では、長文の分析能力と正確性が重要である。複数のツールを試し、自分の研究分野で最も精度の高い回答を返すツールを見つけることを推奨する。

\section{検索特化型AI}
\index{Perplexity@Perplexity}\index{Consensus@Consensus}\index{Elicit@Elicit}

\begin{table}[H]
\centering
\footnotesize
\resizebox{\textwidth}{!}{%
\begin{tabular}{llll}
\toprule
ツール名 & 開発元 & 料金 & 得意分野 \\
\midrule
Perplexity & Perplexity AI & 無料版あり/Pro \$20/月 & 出典付き検索、最新情報の収集 \\
Consensus & Consensus & 無料版あり/Premium \$12/月 & 学術論文に特化した検索・要約 \\
Elicit & Elicit & 無料版あり/Plus \$12/月 & 学術論文の検索・分析・要約 \\
\bottomrule
\end{tabular}}
\end{table}

\textbf{選び方のポイント}: 文献調査の初期段階では、学術論文に特化したConsensusやElicitが有効である。一般的な情報収集にはPerplexityが便利だが、学術的な正確性を求める場合はGoogle Scholarとの併用を推奨する。

\section{論文分析・執筆支援AI}
\index{Semantic Scholar@Semantic Scholar}\index{Writefull@Writefull}\index{Grammarly@Grammarly}

\begin{table}[H]
\centering
\footnotesize
\resizebox{\textwidth}{!}{%
\begin{tabular}{llll}
\toprule
ツール名 & 開発元 & 料金 & 得意分野 \\
\midrule
Semantic Scholar & Allen AI & 無料 & 論文検索、引用分析 \\
Connected Papers & Connected Papers & 無料版あり/月額\$3〜 & 論文の関連性グラフ可視化 \\
Writefull & Writefull & 無料版あり/Premium 16ユーロ/月 & 学術英語の文法・表現チェック \\
Grammarly & Grammarly & 無料版あり/Premium \$12/月 & 英文の文法・スタイルチェック \\
DeepL Write & DeepL & 無料版あり/Pro 8ユーロ/月 & 英文の改善・書き換え提案 \\
\bottomrule
\end{tabular}}
\end{table}

\textbf{選び方のポイント}: 論文執筆では、Writefullのように学術英語に特化したツールが有効である。GrammarlyやDeepL Writeは一般的な英文チェックに優れるが、学術特有の表現については専門ツールが勝る。

\section{コード支援AI}
\index{GitHub Copilot@GitHub Copilot}\index{Claude Code@Claude Code}\index{Cursor@Cursor}

\begin{table}[H]
\centering
\footnotesize
\resizebox{\textwidth}{!}{%
\begin{tabular}{llll}
\toprule
ツール名 & 開発元 & 料金 & 得意分野 \\
\midrule
GitHub Copilot & GitHub/OpenAI & 無料枠あり/Pro \$10/月 & コード補完・生成 \\
Claude Code & Anthropic & Pro/Maxプランに含まれる & コード生成、デバッグ \\
Cursor & Cursor & 無料版あり/Pro \$20/月 & AI統合エディタ \\
Google Colab AI & Google & 無料版あり/Pro \$12/月 & Pythonコード支援 \\
\bottomrule
\end{tabular}}
\end{table}

\textbf{選び方のポイント}: 研究でプログラミングを行う場合、GitHub Copilotは最も広く使われている。学生は無料で利用できるため、まずこれを試すことを推奨する。より高度なコード生成やプロジェクト全体の理解が必要な場合は、Claude CodeやCursorが有効である。

\section{デザイン・プレゼンテーションAI}
\index{Canva@Canva}\index{Gamma@Gamma}\index{Mermaid@Mermaid}

\begin{table}[H]
\centering
\footnotesize
\begin{tabular}{lp{2cm}p{2.5cm}p{3cm}}
\toprule
ツール名 & 開発元 & 料金 & 得意分野 \\
\midrule
Canva（AI機能） & Canva & 無料版あり/Pro \$13/月 & スライド・ポスターデザイン \\
Gamma & Gamma & 無料版あり/Plus \$8/月 & AI自動スライド生成 \\
Beautiful.ai & Beautiful.ai & \$12/月〜 & デザイン自動調整スライド \\
Mermaid & オープンソース & 無料 & テキストから図表生成 \\
\bottomrule
\end{tabular}
\end{table}

\textbf{選び方のポイント}: 学会発表用のスライドは、\LaTeX{} BeamerやPowerPointで作成するのが一般的だが、Gammaを使って初期構成を自動生成し、その後手動で調整するワークフローも効果的である。図表の生成にはMermaidが便利で、Markdown文書に直接埋め込める。

\chapter{プロンプトテンプレート集}
\label{app:prompts}

本付録では、本書で紹介したプロンプトのエッセンスをテンプレートとして整理する。各テンプレートには、ROCF（Role, Output, Context, Format）の各要素がどこに対応するかを注釈として示している。

\begin{itembox}[l]{ROCF要素の凡例}
\begin{itemize}
\item \textbf{【R】}Role（役割）: AIに与える役割設定
\item \textbf{【O】}Output（出力）: 求める出力の内容
\item \textbf{【C】}Context（文脈）: 背景情報や条件
\item \textbf{【F】}Format（形式）: 出力の形式指定
\end{itemize}
\end{itembox}

\section{文献調査用プロンプト}

\begin{promptbox}
\textbf{テンプレート1: サーベイの起点を見つける}

【R】あなたは[分野名]の研究者で、文献調査の専門家です。

【O】以下のテーマに関するサーベイ論文・レビュー論文を\\
探すための検索戦略を提案してください。

【C】テーマ: [テーマ] / 知りたいこと: [具体的な問い] /\\
すでに知っていること: [既知の情報]

【F】推奨する検索キーワード（英語）5セット /\\
推奨する学術データベース / 注意すべき関連分野
\end{promptbox}

\begin{promptbox}
\textbf{テンプレート2: 論文の要約を依頼する}

【R】あなたは学術論文の読解に精通した研究アシスタントです。

【O】以下の論文情報を、構造化された要約にまとめて\\
ください。

【C】タイトル: [タイトル] / 著者: [著者] /\\
概要: [Abstractまたは要点]

【F】研究の目的（2文以内）/ 手法の概要（3文以内）/\\
主な結果（箇条書き3点）/ 限界・今後の課題（2点）/\\
自分の研究との関連（※ここは自分で記入すること）
\end{promptbox}

\begin{promptbox}
\textbf{テンプレート3: 研究動向の整理}

【R】あなたは[分野名]の動向に詳しい研究者です。

【O】以下のキーワードに関する研究動向を整理してください。

【C】キーワード: [キーワード] / 期間: 過去[5]年間 /\\
対象: [学会名/ジャーナル名]

【F】主要な研究潮流（3〜5つ）/ 最近のトレンド・転換点 /\\
未解決の課題・今後の方向性\\
※ 挙げられた論文は必ず実在を確認すること
\end{promptbox}

\begin{promptbox}
\textbf{テンプレート4: 論文間の関係性の整理}

【R】あなたは文献レビューの専門家です。

【O】以下の複数の論文の関係性を整理してください。

【C】論文1: [タイトル, 著者, 年] --- [要点]\\
論文2: [タイトル, 著者, 年] --- [要点]\\
論文3: [タイトル, 著者, 年] --- [要点]

【F】各論文の共通点と相違点を表形式で /\\
研究の発展の流れ（時系列）/ 未解決の問題
\end{promptbox}

\begin{promptbox}
\textbf{テンプレート5: 批判的読解の支援}

【R】あなたは厳格な査読者です。

【O】以下の論文の主張について、批判的な視点からの質問を\\
生成してください。

【C】論文の主張: [主張] / 使用された手法: [手法] /\\
実験の規模: [データ数等]

【F】以下の観点から各3問ずつ: 方法論的な妥当性 /\\
結果の解釈 / 一般化可能性
\end{promptbox}

\section{英語論文用プロンプト}

\begin{promptbox}
\textbf{テンプレート1: Abstractの英訳・改善}

【R】あなたは学術英語ライティングの専門家です。\\{}[分野名]の論文に精通しています。

【O】以下の日本語要旨を、学術英語のAbstractとして\\
書き直してください。

【C】[日本語の要旨] / 対象ジャーナル/学会: [名称] /\\
字数制限: [語数]語以内

【F】構成: Background → Purpose → Methods →\\
Results → Conclusion / 学術英語の慣習に従った表現 /\\
出力後、使用した主要な学術表現のリストを付記
\end{promptbox}

\begin{promptbox}
\textbf{テンプレート2: Introduction の段落構成}

【R】あなたは学術論文のライティングコーチです。

【O】以下の情報をもとに、Introductionの段落構成案を\\
作成してください。

【C】研究分野: [分野] / 研究の背景: [2-3文] /\\
既存研究の課題: [箇条書き] / 本研究の提案: [1-2文]

【F】各段落の役割 / 使うべき英語表現の候補（2-3個）/\\
段落間の接続の仕方
\end{promptbox}

\begin{promptbox}
\textbf{テンプレート3: 英文の校正・改善}

【R】あなたはネイティブの学術英語校正者です。

【O】以下の英文を校正し、より学術的に洗練された表現に\\
改善してください。

【C】[校正対象の英文] / 論文のセクション:\\
{[}Introduction / Methods / Results / Discussion] /\\
分野: [分野名]

【F】修正箇所を明示（元の表現→修正後の表現）/\\
各修正の理由を簡潔に説明 / 全体的な改善提案
\end{promptbox}

\begin{promptbox}
\textbf{テンプレート4: 査読コメントへの回答文作成}

【R】あなたは学術論文の査読対応に精通した研究者です。

【O】以下の査読コメントに対する回答文の下書きを\\
作成してください。

【C】査読コメント: [コメント] /\\
対応方針: [どのように対応するか]

【F】英語で記述 / ``We thank the reviewer for...'' で\\
始める丁寧な書き出し / 論文の修正箇所を明示
\end{promptbox}

\begin{promptbox}
\textbf{テンプレート5: 関連研究セクションの構成}

【R】あなたは学術論文の Related Work セクションの\\
構成アドバイザーです。

【O】以下の先行研究リストを、Related Workセクション\\
として効果的に構成する方法を提案してください。

【C】[先行研究のリスト] / 本研究のテーマ: [テーマ]

【F】サブセクションの分け方 /\\
各サブセクションで述べるべき内容 /\\
先行研究から本研究への橋渡しの書き方
\end{promptbox}

\section{コード生成用プロンプト}

\begin{promptbox}
\textbf{テンプレート1: 分析コードの雛形生成}

【R】あなたはPythonのデータ分析に精通したプログラマーです。

【O】以下の分析タスクを実行するPythonコードを\\
生成してください。

【C】目的: [例: CSVファイルからデータを読み込み、\\
基本統計量を算出] / データの形式: [例: CSV、列名はA, B, C] /\\
使用ライブラリ: [例: pandas, matplotlib]

【F】コメント付きの読みやすいコード / 各ステップの説明 /\\
エラーハンドリングを含む
\end{promptbox}

\begin{promptbox}
\textbf{テンプレート2: 既存コードのリファクタリング}

【R】あなたはコードレビューの専門家です。

【O】以下のコードをリファクタリングしてください。

【C】[リファクタリング対象のコード] /\\
問題点: [例: 可読性が低い、重複が多い]

【F】リファクタリング後のコード /\\
変更点のリスト（何を、なぜ変更したか）
\end{promptbox}

\begin{promptbox}
\textbf{テンプレート3: テストコードの生成}

【R】あなたはソフトウェアテストの専門家です。

【O】以下の関数に対するユニットテストを生成してください。

【C】[テスト対象の関数] / テストフレームワーク: [例: pytest]

【F】正常系テスト（3ケース以上）/\\
異常系テスト（2ケース以上）/\\
エッジケーステスト（2ケース以上）
\end{promptbox}

\begin{promptbox}
\textbf{テンプレート4: エラーのデバッグ支援}

【R】あなたはPythonのデバッグに精通したエンジニアです。

【O】以下のエラーの原因と修正方法を教えてください。

【C】コード: [エラーが発生するコード] /\\
エラーメッセージ: [エラーメッセージ] /\\
実行環境: [例: Python 3.11, macOS]

【F】エラーの原因の説明（初学者にも分かるように）/\\
修正後のコード / 同様のエラーを防ぐためのヒント
\end{promptbox}

\begin{promptbox}
\textbf{テンプレート5: README.mdの生成}

【R】あなたはオープンソースプロジェクトのドキュメント専門家です。

【O】以下のプロジェクト情報から、GitHub用の\\
README.mdを生成してください。

【C】プロジェクト名: [名前] / 概要: [1-2文] /\\
使用言語: [例: Python] / 主な機能: [箇条書き]

【F】英語で記述 / 構成: Title → Badge → Overview →\\
Installation → Usage → Project Structure → License
\end{promptbox}

\section{データ分析用プロンプト}

\begin{promptbox}
\textbf{テンプレート1: 分析計画の立案}

【R】あなたはデータサイエンティストです。

【O】以下のデータと研究目的に対する分析計画を\\
提案してください。

【C】研究目的: [目的] /\\
データの概要: [データの種類、サイズ、変数] /\\
仮説: [検証したい仮説]

【F】推奨する分析手法（3つ以上）と理由 /\\
分析の手順（ステップバイステップ）/ 必要な前処理 /\\
結果の解釈の仕方
\end{promptbox}

\begin{promptbox}
\textbf{テンプレート2: グラフの改善提案}

【R】あなたはデータ可視化の専門家です。

【O】以下のグラフの改善提案をしてください。

【C】グラフの種類: [例: 棒グラフ] /\\
何を示したいか: [意図] / 現在の問題: [例: 見づらい] /\\
使用ツール: [例: matplotlib]

【F】改善ポイント（5つ以上）/ 改善後のコード例 /\\
なぜその改善が効果的かの説明
\end{promptbox}

\begin{promptbox}
\textbf{テンプレート3: 統計的検定の選択}

【R】あなたは統計学の専門家です。

【O】以下の研究状況に適した統計的検定手法を\\
推薦してください。

【C】比較したいもの: [例: 2群の平均値の差] /\\
データの性質: [例: 正規分布かどうか、対応の有無] /\\
有意水準: [例: 0.05]

【F】推薦する検定手法とその理由 / 前提条件の確認方法 /\\
Pythonでの実装コード / 結果の解釈の仕方 /\\
論文での記載例（英語）
\end{promptbox}

\section{スライド作成用プロンプト}

\begin{promptbox}
\textbf{テンプレート1: スライド構成の立案}

【R】あなたは学会発表のプレゼンテーションコーチです。

【O】以下の研究内容を、[発表時間]分の学会発表\\
スライドの構成案にしてください。

【C】研究タイトル: [タイトル] / 対象学会: [学会名] /\\
聴衆: [同分野の研究者/広い分野/実務者] /\\
研究の概要: [3-5文]

【F】各スライドについて: スライドタイトル /\\
含めるべき内容 / 推定所要時間 / 注意点
\end{promptbox}

\begin{promptbox}
\textbf{テンプレート2: スライドの文字量削減}

【R】あなたはプレゼンテーションデザインの専門家です。

【O】以下のスライドの内容を、視覚的に見やすくなるよう\\
文字量を削減してください。

【C】現在のスライド内容: [テキスト] / 最も伝えたいこと: [1文]

【F】削減後のテキスト（元の50\%以下の文字数）/\\
キーワードの強調提案 / 図表化できる部分の提案
\end{promptbox}

\begin{promptbox}
\textbf{テンプレート3: Q\&A想定質問の準備}

【R】あなたは[分野名]の学会に参加する研究者です。

【O】以下の発表内容に対して、Q\&Aセッションで出そうな\\
質問を生成してください。

【C】発表タイトル: [タイトル] / 発表の概要: [3-5文] /\\
使用した手法: [手法] / 主な結果: [結果]

【F】以下のカテゴリに分けて各3問:\\
手法に関する質問 / 結果の解釈に関する質問 /\\
今後の展望に関する質問 / 想定される厳しい質問\\
各質問に回答のポイントを付記
\end{promptbox}

\section{ガクチカ・ES用プロンプト}

\begin{promptbox}
\textbf{テンプレート1: 経験の棚卸し}

【R】あなたは大学院生専門のキャリアカウンセラーです。

【O】以下の情報をもとに、ガクチカの素材を引き出す\\
ための深掘り質問を10個生成してください。

【C】専攻: [専攻] / 研究テーマ: [テーマ] /\\
研究以外の活動: [箇条書き] /\\
大学院生活で苦労したこと: [箇条書き]

【F】各質問に対して: 質問文 /\\
この質問が引き出そうとしているポイント / 回答のヒント
\end{promptbox}

\begin{promptbox}
\textbf{テンプレート2: STAR法での構造化}

【R】あなたはESライティングの専門コーチです。

【O】以下の経験メモを、STAR法で構造化してください。

【C】[経験の箇条書きメモ] / 志望業界: [業界]

【F】S（状況）: [3文以内] / T（課題）: [2文以内] /\\
A（行動）: [3-4文] /\\
R（結果）: [2文以内、定量的な成果を含む]\\
＋各要素の改善ポイント
\end{promptbox}

\begin{promptbox}
\textbf{テンプレート3: 採用担当者視点での添削}

【R】あなたは[業界名]の採用担当者で、\\
年間500本以上のESを読んでいます。

【O】以下のガクチカを、7つの観点で評価・添削して\\
ください。

【C】[ガクチカの文章] / 志望職種: [職種]

【F】以下の各観点について ○/△/× と改善提案:\\
(1) 第一印象 (2) 課題の明確さ (3) 行動の具体性\\
(4) 主体性の伝わり方 (5) 成果の定量性\\
(6) 入社後の再現性 (7) 全体の読みやすさ\\
＋修正案（修正箇所を明示）
\end{promptbox}

\chapter{学術英語表現集}
\label{app:english}

論文の各セクションで使える英語表現を整理する。これらの表現は、AIに英文生成を依頼する際の参考にもなるし、AIの出力をチェックする際の基準にもなる。

\section{Introduction で使える表現}

\subsection*{研究背景の記述}

\begin{table}[H]
\centering
\small
\begin{tabular}{p{4.5cm}p{7cm}}
\toprule
日本語の意図 & 英語表現 \\
\midrule
近年、〜が注目されている & In recent years, ... has attracted considerable attention. \\
〜は重要な課題である & ... remains a critical challenge in the field of ... \\
〜の需要が高まっている & There is a growing demand for ... \\
〜は広く研究されてきた & ... has been extensively studied in the literature. \\
〜の進展に伴い & With the rapid advancement of ... \\
〜は〜において重要な役割を果たす & ... plays a crucial role in ... \\
〜はますます重要になっている & ... has become increasingly important. \\
〜の分野では & In the domain of ... \\
〜が実用化されつつある & ... is being deployed in real-world applications. \\
〜に関する研究が活発化している & A growing body of research has focused on ... \\
\bottomrule
\end{tabular}
\end{table}

\subsection*{研究目的の記述}

\begin{table}[H]
\centering
\small
\begin{tabular}{p{4.5cm}p{7cm}}
\toprule
日本語の意図 & 英語表現 \\
\midrule
本研究では〜を提案する & In this paper, we propose ... \\
本研究の目的は〜である & The aim of this study is to ... \\
本研究は〜に取り組む & This work addresses the problem of ... \\
〜を明らかにすることを目指す & We aim to elucidate ... \\
本研究の主な貢献は以下の通りである & The main contributions of this paper are as follows: \\
我々の知る限り、初めて & To the best of our knowledge, this is the first study to ... \\
\bottomrule
\end{tabular}
\end{table}

\subsection*{論文の構成の記述}

\begin{table}[H]
\centering
\small
\begin{tabular}{p{4.5cm}p{7cm}}
\toprule
日本語の意図 & 英語表現 \\
\midrule
本論文は以下のように構成される & The remainder of this paper is organized as follows. \\
第2節では〜について述べる & Section 2 describes ... \\
第3節では〜を紹介する & Section 3 presents ... \\
最後に、第6節で結論を述べる & Finally, Section 6 concludes the paper. \\
\bottomrule
\end{tabular}
\end{table}

\section{Methods で使える表現}

\subsection*{手法の説明}

\begin{table}[H]
\centering
\small
\begin{tabular}{p{4.5cm}p{7cm}}
\toprule
日本語の意図 & 英語表現 \\
\midrule
提案手法の概要を述べる & We provide an overview of the proposed method. \\
〜のアルゴリズムを以下に示す & The algorithm for ... is presented below. \\
本研究では〜を採用した & In this study, we adopted ... \\
〜を〜に適用した & We applied ... to ... \\
〜は以下の3ステップから成る & ... consists of the following three steps. \\
〜に基づいて & Based on ... / Building upon ... \\
〜を拡張して & By extending ... / We extend ... to ... \\
図Xは〜の概要を示す & Figure X illustrates the overview of ... \\
〜と定義する & We define ... as ... \\
〜を仮定する & We assume that ... \\
\bottomrule
\end{tabular}
\end{table}

\subsection*{実験条件の記述}

\begin{table}[H]
\centering
\small
\begin{tabular}{p{4.5cm}p{7cm}}
\toprule
日本語の意図 & 英語表現 \\
\midrule
実験環境は以下の通りである & The experimental setup is as follows. \\
〜を用いて評価を行った & We conducted the evaluation using ... \\
パラメータは〜に設定した & The parameter was set to ... \\
〜を比較手法として用いた & We used ... as a baseline for comparison. \\
実験は〜回繰り返した & The experiment was repeated ... times. \\
〜の条件下で & Under the condition of ... \\
\bottomrule
\end{tabular}
\end{table}

\section{Results で使える表現}

\subsection*{結果の記述}

\begin{table}[H]
\centering
\small
\begin{tabular}{p{4.5cm}p{7cm}}
\toprule
日本語の意図 & 英語表現 \\
\midrule
表Xは〜の結果を示す & Table X shows the results of ... \\
図Xに〜を示す & Figure X presents ... \\
結果から〜が分かる & The results indicate that ... \\
〜が観察された & ... was observed. \\
〜は〜を上回った & ... outperformed ... \\
〜において有意な差が見られた & A significant difference was found in ... \\
〜は〜と一致する & ... is consistent with ... \\
予想通り、〜であった & As expected, ... \\
興味深いことに、〜であった & Interestingly, ... \\
〜は〜\%の改善を示した & ... showed an improvement of ...\% over ... \\
\bottomrule
\end{tabular}
\end{table}

\subsection*{比較の記述}

\begin{table}[H]
\centering
\small
\begin{tabular}{p{4.5cm}p{7cm}}
\toprule
日本語の意図 & 英語表現 \\
\midrule
〜と比較して & Compared with ... / In comparison with ... \\
〜に対して & In contrast to ... / As opposed to ... \\
同様に & Similarly, ... / Likewise, ... \\
一方で & On the other hand, ... / Meanwhile, ... \\
特に & In particular, ... / Notably, ... \\
\bottomrule
\end{tabular}
\end{table}

\section{Discussion/Conclusion で使える表現}

\subsection*{考察}

\begin{table}[H]
\centering
\small
\begin{tabular}{p{4.5cm}p{7cm}}
\toprule
日本語の意図 & 英語表現 \\
\midrule
この結果は〜を示唆している & These results suggest that ... \\
〜の理由として考えられるのは & A possible explanation for ... is that ... \\
この知見は〜と整合する & This finding is in line with ... \\
〜の観点から考えると & From the perspective of ... \\
〜は〜によって説明できる & ... can be attributed to ... \\
さらなる調査が必要である & Further investigation is warranted. \\
\bottomrule
\end{tabular}
\end{table}

\subsection*{限界の記述}

\begin{table}[H]
\centering
\small
\begin{tabular}{p{4.5cm}p{7cm}}
\toprule
日本語の意図 & 英語表現 \\
\midrule
本研究にはいくつかの限界がある & This study has several limitations. \\
本研究の限界として & A limitation of this study is that ... \\
〜は本研究の範囲外である & ... is beyond the scope of this study. \\
データセットの規模が限られている & The dataset size is limited. \\
一般化には注意が必要である & Caution is needed when generalizing ... \\
\bottomrule
\end{tabular}
\end{table}

\subsection*{展望・結論}

\begin{table}[H]
\centering
\small
\begin{tabular}{p{4.5cm}p{7cm}}
\toprule
日本語の意図 & 英語表現 \\
\midrule
今後の研究では〜を検討する & In future work, we plan to investigate ... \\
〜への拡張が期待される & Extension to ... is a promising direction. \\
本研究は〜に貢献する & This study contributes to ... \\
結論として & In conclusion, ... / To summarize, ... \\
本研究で提案した手法は〜に有用である & The proposed method is useful for ... \\
\bottomrule
\end{tabular}
\end{table}

\section{接続表現・論理展開表現}

\subsection*{追加・列挙}

\begin{table}[H]
\centering
\small
\begin{tabular}{ll}
\toprule
機能 & 表現 \\
\midrule
さらに & Furthermore, / Moreover, / In addition, / Additionally, \\
第一に…第二に & First, ... Second, ... / Firstly, ... Secondly, ... \\
同様に & Similarly, / Likewise, / In the same vein, \\
\bottomrule
\end{tabular}
\end{table}

\subsection*{対比・逆接}

\begin{table}[H]
\centering
\small
\begin{tabular}{ll}
\toprule
機能 & 表現 \\
\midrule
しかし & However, / Nevertheless, / Nonetheless, \\
一方で & On the other hand, / Conversely, / In contrast, \\
〜にもかかわらず & Despite ... / In spite of ... / Notwithstanding ...\footnote{やや古風な表現。現代の論文では Despite や In spite of が好まれる。} \\
〜であるが & Although ... / While ... / Whereas ... \\
\bottomrule
\end{tabular}
\end{table}

\subsection*{因果・理由}

\begin{table}[H]
\centering
\small
\begin{tabular}{ll}
\toprule
機能 & 表現 \\
\midrule
したがって & Therefore, / Thus, / Hence, / Consequently, \\
〜のために & Due to ... / Owing to ... / Because of ... \\
その結果 & As a result, / As a consequence, \\
〜を考慮すると & Considering ... / Given that ... \\
\bottomrule
\end{tabular}
\end{table}

\subsection*{例示・具体化}

\begin{table}[H]
\centering
\small
\begin{tabular}{ll}
\toprule
機能 & 表現 \\
\midrule
例えば & For example, / For instance, / e.g., \\
具体的には & Specifically, / More specifically, / In particular, \\
すなわち & That is, / Namely, / i.e., \\
\bottomrule
\end{tabular}
\end{table}

\subsection*{要約・結論}

\begin{table}[H]
\centering
\small
\begin{tabular}{ll}
\toprule
機能 & 表現 \\
\midrule
要するに & In summary, / To summarize, / In brief, \\
結論として & In conclusion, / To conclude, \\
全体として & Overall, / On the whole, / Taken together, \\
\bottomrule
\end{tabular}
\end{table}

\chapter{15回授業計画}
\label{app:syllabus}

本書を教科書として半期15回の授業で使用する場合の授業計画を示す。各回90分の授業を想定し、講義と演習のバランスを考慮している。

\section{授業計画表}

\footnotesize
\begin{longtable}{ccp{1.8cm}p{3.3cm}p{1.8cm}}
\toprule
回 & 対応章 & テーマ & 授業の進め方 & 課題 \\
\midrule
\endfirsthead
\toprule
回 & 対応章 & テーマ & 授業の進め方 & 課題 \\
\midrule
\endhead
\midrule
\multicolumn{5}{r}{（次ページに続く）}\\
\endfoot
\bottomrule
\endlastfoot

1 & 第1章 & 研究活動と生成AI入門 & 【講義60分】AI活用の意義と倫理、批判的AI活用の考え方を解説。【演習30分】各自が普段使っているAIツールとその使い方を共有するワークショップ。 & 自分の研究活動におけるAI活用の現状と課題を400字でまとめる \\
\hline
2 & 第2章前半 & プロンプトエンジニアリング基礎 & 【講義40分】ROCF（Role, Output, Context, Format）の解説と実例紹介。【演習50分】同じタスクに対して異なるプロンプトを試し、出力の違いを比較する実践演習。 & 自分の研究テーマに関するプロンプトをROCFで設計し、3パターン作成 \\
\hline
3 & 第2章後半 & プロンプトエンジニアリング実践 & 【講義30分】Few-shot, Chain-of-Thought, 段階的プロンプトの解説。【演習60分】研究テーマに関する実践的なプロンプト作成と改善のサイクルを体験。 & 最も効果的だったプロンプトとその改善過程をレポートにまとめる \\
\hline
4 & 第3章 & AIを活用した文献調査 & 【講義30分】AI時代の文献調査戦略、ハルシネーション対策。【演習60分】自分の研究テーマについて、AIを使って文献サーベイの初期段階を実施。 & 文献調査レポート（先行研究5本以上の要約と比較表） \\
\hline
5 & 第4章 & 英語論文の読解と執筆支援 & 【講義40分】学術英語の構造、AIを使った英文作成・校正の方法。【演習50分】自分の研究のAbstractを日本語→英語で作成する実践。 & 自分の研究の英語Abstract（150--250語）を作成 \\
\hline
6 & 第5章 & プレゼンテーション準備 & 【講義30分】効果的な学会発表の構成、AIを使ったスライド作成。【演習60分】5分間の研究紹介スライド（5--7枚）の構成をAIと共に設計。 & 研究紹介スライドの完成版を提出 \\
\hline
7 & 第6章前半 & プログラミングとAI（基礎） & 【講義30分】AI支援コーディングの原則、GitHub Copilot等の使い方。【演習60分】自分の研究に関連するデータ処理コードをAIと共に作成。 & 作成したコードとAIとのやり取りの記録を提出 \\
\hline
8 & 第6章後半 & プログラミングとAI（応用） & 【講義20分】コードレビュー、テスト、デバッグへのAI活用。【演習40分】前回のコードのリファクタリングとテスト作成。【発表30分】中間発表（3分×12名程度）。 & コードを改善し、README付きで提出 \\
\hline
9 & 第7章 & データ分析・可視化とAI & 【講義30分】AIを使ったデータ分析の計画立案、グラフ作成の原則。【演習60分】自分の研究データ（またはサンプルデータ）を使った分析と可視化の実践。 & データ分析結果のグラフ3枚と各グラフの解釈を記述したレポート \\
\hline
10 & 第8章 & GitHub Webページの作成 & 【講義20分】GitHub Pagesの仕組み、研究成果のWeb公開の意義。【演習70分】自分の研究紹介Webページの作成（GitHub Pages）。 & GitHub PagesによるWebページを公開し、URLを提出 \\
\hline
11 & 第9章前半 & ガクチカ・ES作成（基礎） & 【講義30分】就活スケジュール、企業の評価ポイント、STAR法の解説。【演習60分】AIを使った経験の棚卸しとSTAR法での構造化。 & STAR法で構造化した経験メモ（2テーマ分） \\
\hline
12 & 第9章後半 & ガクチカ・ES作成（実践） & 【講義20分】文章化のテクニック、面接との連動。【演習70分】ガクチカの文章化→AI添削→修正のサイクルを2回以上実施。ペアで相互レビュー。 & ガクチカ完成版（400字以内）、Before/After比較付き \\
\hline
13 & 第10章前半 & 研究の新規性分析 & 【講義30分】新規性の3つの観点、言語化の方法。【演習60分】自分の研究について新規性分析ワークフローを実施。 & 新規性分析レポート（テンプレートに沿って作成） \\
\hline
14 & 第10章後半 & 研究ポートフォリオの統合 & 【講義20分】ポートフォリオの構成、GitHubリポジトリの整理方法。【演習70分】これまでの成果物をGitHubリポジトリとして統合・整理。 & 研究ポートフォリオ（GitHubリポジトリ）の完成版を提出 \\
\hline
15 & 総合 & 最終発表と振り返り & 【発表60分】研究ポートフォリオのプレゼンテーション（5分×12名程度）。【振り返り30分】AI活用の全体振り返り、今後の活用方針の策定。 & AI活用方針書（A4用紙1枚）を最終課題として提出 \\
\end{longtable}

\section{授業運営のヒント}

\subsection*{事前準備}

\begin{itemize}
\item 受講者全員がChatGPTまたはClaudeのアカウントを持っていることを確認する（第1回の前に案内）
\item GitHubアカウントの作成も第1回の前に案内する
\item GitHub Copilotの学生無料枠の申請は手続きに時間がかかるため、早めに案内する
\end{itemize}

\subsection*{演習の進め方}

\begin{itemize}
\item 演習時間は、受講者が実際にAIと対話しながら作業する時間である。教員は巡回してサポートする
\item 可能であれば、受講者同士のペアワーク・グループワークを積極的に取り入れる。他の人のプロンプトや使い方を見ることが大きな学びとなる
\item 「うまくいかなかった例」も学びの素材として共有する
\end{itemize}

\subsection*{評価方法の案}

\begin{table}[H]
\centering
\small
\begin{tabular}{llp{6cm}}
\toprule
評価項目 & 配分 & 評価基準 \\
\midrule
各回の課題 & 40\% & 期限内の提出、指示への対応度 \\
中間発表 & 10\% & 研究内容の伝達力、スライドの質 \\
研究ポートフォリオ & 30\% & 成果物の質と統合度、GitHubの整理状況 \\
最終発表 & 10\% & プレゼンテーション力、質疑応答 \\
AI活用方針書 & 10\% & 振り返りの深さ、今後への具体的な方針 \\
\bottomrule
\end{tabular}
\end{table}

\subsection*{ポートフォリオの評価ルーブリック}

\begin{table}[H]
\centering
\footnotesize
\begin{tabular}{lp{1.8cm}p{1.8cm}p{1.8cm}p{1.8cm}}
\toprule
観点 & A（優秀） & B（良好） & C（合格） & D（要改善） \\
\midrule
成果物の網羅性 & 全章の成果物が揃っている & 大半の成果物が揃っている & 必須の成果物は揃っている & 成果物が不足している \\
\hline
成果物の質 & 高品質で研究内容が正確に反映 & 一定の質を持つ & 最低限の要件を満たしている & 質に問題がある \\
\hline
AI活用の適切さ & 効果的に活用しつつ批判的検証ができている & 概ね適切に検証している & AIを使っているが検証が不十分 & AIの出力を無批判に使用 \\
\hline
GitHubの整理 & README充実、構造明確、再現性確保 & 構造が概ね整理されている & 最低限の構造がある & 整理されていない \\
\hline
オリジナリティ & 独自の工夫や視点が随所に見られる & 一部に独自の工夫がある & テンプレート通り & 独自性がない \\
\bottomrule
\end{tabular}
\end{table}

\section{授業の変形パターン}

本書は15回の授業を想定しているが、以下のような変形も可能である。

\subsection*{8回集中講座の場合}

\begin{table}[H]
\centering
\begin{tabular}{cl}
\toprule
回 & 内容 \\
\midrule
1 & 第1章+第2章: AI基礎とプロンプトエンジニアリング \\
2 & 第3章+第4章: 文献調査と英語論文 \\
3 & 第5章: プレゼンテーション \\
4 & 第6章: プログラミング \\
5 & 第7章: データ分析 \\
6 & 第8章+第9章: Web公開と就活文書 \\
7 & 第10章: 新規性分析とポートフォリオ \\
8 & 最終発表 \\
\bottomrule
\end{tabular}
\end{table}

\subsection*{自主ゼミ（週1回、2時間）の場合}

各回の演習部分を充実させ、参加者同士のディスカッションを中心に進める。講義部分は事前に本書の該当章を読んできてもらい、演習とディスカッションに時間を割く反転授業形式が効果的である。


% 索引
\printindex

\end{document}
